% Le courant alternatif — PCT 3ème — Cours signé OnBuch+
% Programme officiel MINESEC. Compiler avec : tectonic N12-courant-alternatif.tex
\newcommand{\DOCMATIERE}{PCT}
\newcommand{\DOCNIVEAU}{3ème}
\newcommand{\DOCMODULE}{Module 2 : Électricité}
\newcommand{\DOCLECON}{N12}
\newcommand{\DOCTITRE}{Le courant alternatif}
% OnBuch+ — préambule « cours signés » ENRICHI (compilé avec tectonic / XeLaTeX)
% Système visuel « Pop » d'OnBuch+ : fond crème (jamais blanc), bordures encre
% épaisses, ombres « dures » décalées, titres Archivo Black en capitales.
% Version MATHÉMATIQUES : graphes (pgfplots/tikz), boîtes pédagogiques
% (définition, propriété, méthode, attention, le savais-tu, exercice résolu,
% exercice, TP, bilan), page de garde et signature OnBuch+ ancrée en bas.
%
% Macros attendues dans main.tex AVANT \input{preamble} :
%   \DOCMATIERE  \DOCNIVEAU  \DOCMODULE  \DOCLECON (numéro)  \DOCTITRE
\documentclass[11pt]{article}

\usepackage{fontspec}
\usepackage[french]{babel}
\usepackage[a4paper,margin=2cm,top=3.4cm,bottom=2.8cm,headheight=1.4cm,footskip=1.3cm]{geometry}
\usepackage{amsmath,amssymb,amsfonts}
\usepackage{mathtools} % \xmapsto, \coloneqq... souvent utilisés par les modèles
\usepackage{cancel} % simplifications barrées (bilans de forces)
% \abs{z} (module, valeur absolue) et \norm{u} (norme), utilisés par les modèles
\providecommand{\abs}{}\let\abs\relax\DeclarePairedDelimiter{\abs}{\lvert}{\rvert}
\providecommand{\norm}{}\let\norm\relax\DeclarePairedDelimiter{\norm}{\lVert}{\rVert}
\usepackage[table]{xcolor} % [table] : fournit aussi \rowcolors (alternance de lignes)
\usepackage{pagecolor}
\usepackage{tikz}
\usetikzlibrary{arrows.meta,positioning,calc,shapes.geometric,shapes.misc,shapes.arrows,decorations.pathmorphing,decorations.pathreplacing,decorations.markings,patterns,shadows,mindmap,trees,fit,backgrounds,matrix,intersections,angles,quotes}
\usepackage{pgfplots}
\usepackage[version=4]{mhchem} % équations nucléaires \ce{^{235}_{92}U}
\usepackage[european,siunitx]{circuitikz} % schémas électriques
\pgfplotsset{compat=1.18}
\usepgfplotslibrary{fillbetween,groupplots} % aires entre courbes (calcul intégral)
\usepackage{enumitem}
\usepackage{fancyhdr}
% [most] charge toutes les bibliothèques tcolorbox (listings, minted, raster,
% documentation...) et gonfle inutilement la mémoire TeX sur un document long
% (~300 boîtes) : on ne charge que ce qui est réellement utilisé.
\usepackage[skins,breakable]{tcolorbox}
\usepackage{graphicx}
\usepackage{colortbl}
\usepackage{booktabs}
\usepackage{tabularx}
\usepackage{multirow}
\usepackage{diagbox} % case d'en-tête diagonale des tableaux à double entrée
% Colonnes extensibles centrée / gauche / droite (C, L, R), souvent utilisées par les modèles
\newcolumntype{C}{>{\centering\arraybackslash}X}
\newcolumntype{L}{>{\raggedright\arraybackslash}X}
\newcolumntype{R}{>{\raggedleft\arraybackslash}X}
\usepackage{array}
\usepackage{multicol}
\usepackage{siunitx}
% parse-numbers=false : accepte aussi \SI{2,5 \times 10^{-3}}{...}
\sisetup{locale=FR,output-decimal-marker={,},group-minimum-digits=4,parse-numbers=false}
\DeclareSIUnit\ohm{\ensuremath{\symup{\Omega}}} % U+2126 absent de la police
\DeclareSIUnit\division{div} % réglages d'oscilloscope (ms/div, V/div)
\DeclareSIUnit\tr{tr} % tours (vitesse de rotation, ex. \SI{1500}{\tr\per\minute})
\DeclareSIUnit\molar{M} % concentration molaire (ex. \SI{3}{\molar}, \SI{1}{\micro\molar})
\DeclareSIUnit\an{an} % année (le modèle confond parfois avec \year, un compteur TeX) — ex. \SI{5}{\kilo\meter\per\an}
\DeclareSIUnit\FCFA{FCFA} % franc CFA (ex. \SI{500}{\FCFA\per\kilo\gram})
\DeclareSIUnit\degreeBrix{\textdegree Brix} % teneur en sucre d'un sirop/jus (réfractométrie)
\DeclareSIUnit\month{mois} % le modèle utilise parfois \month, un compteur TeX, comme unité de durée
\DeclareSIUnit\microliter{\micro\liter} % le modèle écrit parfois \microliter en un seul token
\DeclareSIUnit\million{millions} % dénombrement (ex. \SI{15}{\million\per\mL} spermatozoïdes)
\usepackage{qrcode}
% \degree : commande fréquemment utilisée par les modèles pour un angle ou une
% température en dehors de \SI{}{} (non fournie par les paquets chargés).
\providecommand{\degree}{^{\circ}}
% \qed (fin de démonstration), utilisé par les modèles sans amsthm
\providecommand{\qed}{\hfill$\square$}
% \barre : double barre des valeurs interdites dans un tableau de variations (tkz-tab)
\providecommand{\barre}{\ensuremath{\|}}
% environnement proof (amsthm non chargé) utilisé par les modèles
\ifdefined\proof\else\newenvironment{proof}[1][Démonstration]{\par\noindent\textit{#1.}\ }{\qed\par}\fi
\usepackage{lastpage}
\usepackage{hyperref}
\hypersetup{hidelinks}

% ============================================================
% Palette « Pop » OnBuch+ — mêmes hex que la classe P (plus_ui.dart)
% ============================================================
\definecolor{poporange}{HTML}{F59321}
\definecolor{popdark}{HTML}{E07A0C}
\definecolor{popcream}{HTML}{FFF6E8}
\definecolor{popink}{HTML}{1C1714}
\definecolor{poppurple}{HTML}{7A5AE0}
\definecolor{popgreen}{HTML}{2E9E6A}
\definecolor{poppink}{HTML}{E0457B}
\definecolor{popblue}{HTML}{2F7DE1}
\definecolor{popgold}{HTML}{C98A12}
\definecolor{popmuted}{HTML}{8A7F76}
\definecolor{popline}{HTML}{EADFD2}
\definecolor{popgray}{HTML}{8A7F76}
\colorlet{popgrey}{popgray}
% Alias tolérants (le modèle nomme parfois les couleurs différemment)
\colorlet{popwhite}{white}
\colorlet{popblack}{popink}
\colorlet{popred}{poppink}
\colorlet{popyellow}{popgold}
% Teintes claires (fonds de tableaux/encadrés) — le modèle nomme parfois
% les couleurs avec un suffixe L (ex. popblueL) sans les avoir définies.
\colorlet{poporangeL}{poporange!20}
\colorlet{popdarkL}{popdark!20}
\colorlet{poppurpleL}{poppurple!20}
\colorlet{popgreenL}{popgreen!20}
\colorlet{poppinkL}{poppink!20}
\colorlet{popblueL}{popblue!20}
\colorlet{popgoldL}{popgold!20}
\colorlet{popredL}{popred!20}
\colorlet{popgrayL}{popgray!20}
% Teintes claires pour fonds de tableaux / zones
\colorlet{popblueL}{popblue!10!white}
\colorlet{poporangeL}{poporange!14!white}
\colorlet{popgreenL}{popgreen!12!white}
\colorlet{poppurpleL}{poppurple!12!white}
\colorlet{poppinkL}{poppink!12!white}
\colorlet{popgoldL}{popgold!14!white}

\pagecolor{popcream}

% ============================================================
% Typographie
% ============================================================
\defaultfontfeatures{Ligatures=TeX}
\setmainfont{PlusJakartaSans-Medium.ttf}[Path=../../../fonts/,BoldFont=PlusJakartaSans-Bold.ttf]
% Maths en sans-serif (Fira Math) avec chiffres et lettres droites de Plus
% Jakarta, pour que les formules se fondent dans le texte.
\usepackage[math-style=ISO,bold-style=ISO]{unicode-math}
\setmathfont{FiraMath-Regular.otf}[Path=../../../fonts/]
\setmathfont{PlusJakartaSans-Medium.ttf}[Path=../../../fonts/,range={up/num,up/{Latin,latin},\mathup}]
\usepackage{icomma} % virgule décimale sans espace en mode math
% Caractères mathématiques Unicode absents des polices de texte (Plus Jakarta,
% Archivo Black) : rendus par leur équivalent mathématique, en texte comme en
% formule — sinon ils s'affichent en carré vide (ex. « Arithmétique dans ℤ »).
\usepackage{newunicodechar}
\newunicodechar{ℝ}{\ensuremath{\mathbb{R}}}
\newunicodechar{ℤ}{\ensuremath{\mathbb{Z}}}
\newunicodechar{ℂ}{\ensuremath{\mathbb{C}}}
\newunicodechar{ℕ}{\ensuremath{\mathbb{N}}}
\newunicodechar{ℚ}{\ensuremath{\mathbb{Q}}}
\newunicodechar{𝔻}{\ensuremath{\mathbb{D}}}
\newunicodechar{α}{\ensuremath{\alpha}}
\newunicodechar{β}{\ensuremath{\beta}}
\newunicodechar{γ}{\ensuremath{\gamma}}
\newunicodechar{δ}{\ensuremath{\delta}}
\newunicodechar{ε}{\ensuremath{\varepsilon}}
\newunicodechar{θ}{\ensuremath{\theta}}
\newunicodechar{λ}{\ensuremath{\lambda}}
\newunicodechar{μ}{\ensuremath{\mu}}
\newunicodechar{σ}{\ensuremath{\sigma}}
\newunicodechar{τ}{\ensuremath{\tau}}
\newunicodechar{φ}{\ensuremath{\varphi}}
\newunicodechar{ψ}{\ensuremath{\psi}}
\newunicodechar{ω}{\ensuremath{\omega}}
\newunicodechar{Σ}{\ensuremath{\Sigma}}
% majuscules produites par \MakeUppercase dans les titres (α-aminés -> Α)
\newunicodechar{Α}{\ensuremath{\alpha}}
\newunicodechar{Β}{\ensuremath{\beta}}
\newunicodechar{Γ}{\ensuremath{\Gamma}}
\newunicodechar{Δ}{\ensuremath{\Delta}}
\newunicodechar{Ε}{\ensuremath{\varepsilon}}
\newunicodechar{Θ}{\ensuremath{\Theta}}
\newunicodechar{Λ}{\ensuremath{\Lambda}}
\newunicodechar{Μ}{\ensuremath{\mu}}
\newunicodechar{Τ}{\ensuremath{\tau}}
\newunicodechar{Φ}{\ensuremath{\Phi}}
\newunicodechar{Ψ}{\ensuremath{\Psi}}
\newunicodechar{Ω}{\ensuremath{\Omega}}
\newunicodechar{↦}{\ensuremath{\mapsto}}
\newunicodechar{∈}{\ensuremath{\in}}
\newunicodechar{∉}{\ensuremath{\notin}}
\newunicodechar{∧}{\ensuremath{\wedge}}
\newunicodechar{⇔}{\ensuremath{\Leftrightarrow}}
\newunicodechar{⟺}{\ensuremath{\Longleftrightarrow}}
\newunicodechar{⇒}{\ensuremath{\Rightarrow}}
\newunicodechar{⟹}{\ensuremath{\Longrightarrow}}
\newunicodechar{∘}{\ensuremath{\circ}}
\newunicodechar{∪}{\ensuremath{\cup}}
\newunicodechar{∩}{\ensuremath{\cap}}
\newunicodechar{≡}{\ensuremath{\equiv}}
\newunicodechar{⊂}{\ensuremath{\subset}}
\newunicodechar{⊥}{\ensuremath{\perp}}
\newunicodechar{∃}{\ensuremath{\exists}}
\newunicodechar{∀}{\ensuremath{\forall}}
\newunicodechar{∛}{\ensuremath{\sqrt[3]{\,}}}
\newunicodechar{∜}{\ensuremath{\sqrt[4]{\,}}}
\newunicodechar{⃗}{\ensuremath{{}^{\to}}}
\newunicodechar{ⁿ}{\ifmmode^{n}\else\textsuperscript{\ensuremath{n}}\fi}
\newunicodechar{ˣ}{\ifmmode^{x}\else\textsuperscript{\ensuremath{x}}\fi}
\newunicodechar{ᵇ}{\ifmmode^{b}\else\textsuperscript{\ensuremath{b}}\fi}
\newunicodechar{ʳ}{\ifmmode^{r}\else\textsuperscript{\ensuremath{r}}\fi}
\newunicodechar{ᵘ}{\ifmmode^{u}\else\textsuperscript{\ensuremath{u}}\fi}
\newunicodechar{ʸ}{\ifmmode^{y}\else\textsuperscript{\ensuremath{y}}\fi}
\newunicodechar{ᵃ}{\ifmmode^{a}\else\textsuperscript{\ensuremath{a}}\fi}
\newunicodechar{ᵉ}{\ifmmode^{e}\else\textsuperscript{\ensuremath{e}}\fi}
\newunicodechar{ᵏ}{\ifmmode^{k}\else\textsuperscript{\ensuremath{k}}\fi}
\newunicodechar{ᵖ}{\ifmmode^{p}\else\textsuperscript{\ensuremath{p}}\fi}
\newunicodechar{ᴺ}{\ifmmode^{N}\else\textsuperscript{\ensuremath{N}}\fi}
\newunicodechar{ᶜ}{\ifmmode^{c}\else\textsuperscript{\ensuremath{c}}\fi}
\newunicodechar{ʰ}{\ifmmode^{h}\else\textsuperscript{\ensuremath{h}}\fi}
\newunicodechar{ᵠ}{\ifmmode^{q}\else\textsuperscript{\ensuremath{q}}\fi}
\newunicodechar{ₙ}{\ifmmode_{n}\else\textsubscript{\ensuremath{n}}\fi}
\newunicodechar{ₐ}{\ifmmode_{a}\else\textsubscript{\ensuremath{a}}\fi}
\newunicodechar{ᵢ}{\ifmmode_{i}\else\textsubscript{\ensuremath{i}}\fi}
\newunicodechar{ᵣ}{\ifmmode_{r}\else\textsubscript{\ensuremath{r}}\fi}
\newunicodechar{ₖ}{\ifmmode_{k}\else\textsubscript{\ensuremath{k}}\fi}
\newunicodechar{ᵦ}{\ifmmode_{\beta}\else\textsubscript{\ensuremath{\beta}}\fi}
\newunicodechar{ₓ}{\ifmmode_{x}\else\textsubscript{\ensuremath{x}}\fi}
\newfontfamily\popdisplay{ArchivoBlack-Regular.ttf}[Path=../../../fonts/]
\newfontfamily\poplogo{SpaceGrotesk-Bold.ttf}[Path=../../../fonts/]
\newfontfamily\popsemi{PlusJakartaSans-SemiBold.ttf}[Path=../../../fonts/]
\newfontfamily\popxbold{PlusJakartaSans-ExtraBold.ttf}[Path=../../../fonts/]
\newfontfamily\popxboldls{PlusJakartaSans-ExtraBold.ttf}[Path=../../../fonts/,LetterSpace=3.0]

\setlist[enumerate]{leftmargin=1.7em,itemsep=5pt,topsep=4pt}
\setlist[itemize]{leftmargin=1.7em,itemsep=4pt,topsep=4pt,label={\color{poporange}\textbullet}}
\setlength{\parindent}{0pt}
\setlength{\parskip}{5pt}
\renewcommand{\arraystretch}{1.25}

% pgfplots : style Pop par défaut
\pgfplotsset{
  every axis/.append style={axis line style={popink,line width=1pt},tick style={popink},
    grid style={popline,line width=0.5pt},label style={font=\small},tick label style={font=\footnotesize}},
  popcurve/.style={poporange,line width=1.8pt,no markers,smooth},
}
% Même style disponible dans un \draw TikZ simple (hors pgfplots)
\tikzset{popcurve/.style={poporange,line width=1.8pt}}
\tikzset{popgrid/.style={popline,very thin}}
\colorlet{popcurve}{poporange} % le nom du style est parfois utilisé comme couleur

% ============================================================
% Pill / squiggle
% ============================================================
\newcommand{\poppill}[3]{%
  \tikz[baseline=-0.55ex]\node[draw=popink,line width=1.2pt,rounded corners=2.6pt,%
    fill=#2,inner xsep=6.5pt,inner ysep=3pt]{%
    \popxboldls\fontsize{7}{7}\selectfont\color{#3}\MakeUppercase{#1}};%
}
\newcommand{\popsquiggle}[1]{%
  \tikz[baseline=-0.4ex]{
    \draw[#1,line width=2.6pt,line cap=round]
      plot [smooth] coordinates {(0,0.09) (0.34,0.01) (0.68,0.21) (1.02,0.01) (1.36,0.10)};
  }%
}
% Mot-clé surligné (marqueur orange sous le texte)
\newcommand{\cle}[1]{{\popxbold\color{popdark}#1}}

% ============================================================
% En-tête / pied
% ============================================================
\pagestyle{fancy}
\fancyhf{}
\renewcommand{\headrulewidth}{0pt}
\renewcommand{\footrulewidth}{0pt}
\lhead{\raisebox{-0.15ex}{\poplogo\fontsize{13}{13}\selectfont\color{popink}OnBuch\color{poporange}+}}
\rhead{\poppill{\DOCMATIERE\ \textbullet\ \DOCNIVEAU\ \textbullet\ Leçon \DOCLECON}{white}{popink}}
\lfoot{\popxbold\fontsize{7.6}{7.6}\selectfont\color{popmuted}\MakeUppercase{Cours signé OnBuch+}\ \textbullet\ {\popsemi\DOCTITRE}}
\rfoot{\tikz[baseline=-0.6ex]\node[rounded corners=8.5pt,fill=popink,minimum height=17pt,minimum width=17pt,inner xsep=5pt,inner sep=0]{\popxbold\fontsize{8}{8}\selectfont\color{popcream}\thepage\,/\,\pageref*{LastPage}};}

% ============================================================
% Titres
% ============================================================
\newcommand{\chaptitre}[2]{%
  \begin{tcolorbox}[enhanced,colback=poporange,colframe=popink,boxrule=2.6pt,arc=11pt,
      drop shadow=popink,left=15pt,right=15pt,top=12pt,bottom=11pt,before skip=3pt,after skip=18pt]
    \poppill{#2}{popink}{popcream}\par\vspace{9pt}
    {\raggedright\hyphenpenalty=10000\exhyphenpenalty=10000\popdisplay\fontsize{22}{24}\selectfont\color{popcream}\MakeUppercase{#1}\par}
  \end{tcolorbox}
}
% Section numérotée : pastille orange avec le numéro + titre Archivo
\newcounter{popsec}
\newcommand{\coursec}[1]{%
  \stepcounter{popsec}\setcounter{popsub}{0}%
  \par\vspace{16pt}\noindent
  \begin{minipage}{\linewidth}
  \tikz[baseline=-0.7ex]\node[draw=popink,line width=1.6pt,fill=poporange,rounded corners=4pt,minimum size=22pt,inner sep=0]{\popdisplay\fontsize{12}{12}\selectfont\color{popcream}\thepopsec};%
  \hspace{8pt}{\popdisplay\fontsize{15}{17}\selectfont\color{popink}\MakeUppercase{#1}}\par
  \vspace{4pt}\noindent{\color{poporange}\rule{36pt}{3.6pt}}
  \end{minipage}\par\nopagebreak\vspace{8pt}
}
\newcounter{popsub}
\newcommand{\courssub}[1]{%
  \stepcounter{popsub}%
  \par\vspace{9pt}\noindent{\popxbold\fontsize{11.5}{13}\selectfont\color{popink}{\color{poporange}\thepopsec.\thepopsub}\hspace{6pt}#1}\par\nopagebreak\vspace{4pt}
}

% ============================================================
% Boîtes pédagogiques — même recette : fond blanc, bordure épaisse colorée,
% ombre dure encre, badge plein. Titre optionnel [#1].
% ============================================================
\tcbset{popbase/.style={breakable,enhanced,colback=white,boxrule=2.3pt,arc=9pt,
  shadow={3pt}{-3pt}{0pt}{popink},left=14pt,right=14pt,top=11pt,bottom=11pt,before skip=11pt,after skip=15pt,
  fontupper=\popsemi\fontsize{10.6}{14.5}\selectfont\color{popink},
  fontlower=\popsemi\fontsize{10.6}{14.5}\selectfont\color{popink}}}
% Coche dessinée (le glyphe ✓ n'existe pas dans les polices du document)
\newcommand{\popcheck}[1]{\tikz[baseline=-0.5ex]\draw[#1,line width=1.6pt,line cap=round,line join=round](-0.13,0.0)--(-0.04,-0.09)--(0.13,0.1);}
\providecommand{\coche}{\popcheck{popgreen}}
\providecommand{\resultat}[1]{\par\smallskip\noindent\fcolorbox{popgreen}{popgreenL}{\parbox{\dimexpr\linewidth-2\fboxsep-2\fboxrule\relax}{\textbf{Résultat.} #1}}\par\smallskip}
\newcommand{\popboxhead}[3]{\poppill{#1}{#2}{white}\ifx\relax#3\relax\else\hspace{6pt}{\popxbold\fontsize{10.6}{12}\selectfont\color{popink}#3}\fi\par\vspace{7pt}}

\newtcolorbox{aretenir}[1][]{popbase,colframe=popgreen,before upper={\popboxhead{À retenir}{popgreen}{#1}}}
\newtcolorbox{exemplebox}[1][]{popbase,colframe=poporange,before upper={\popboxhead{Exemple}{poporange}{#1}}}
\newtcolorbox{definition}[1][]{popbase,colframe=popblue,before upper={\popboxhead{Définition}{popblue}{#1}}}
\newtcolorbox{propriete}[1][]{popbase,colframe=poppurple,before upper={\popboxhead{Propriété}{poppurple}{#1}}}
\newtcolorbox{methode}[1][]{popbase,colframe=popgold,before upper={\popboxhead{Méthode}{popgold}{#1}}}
\newtcolorbox{attention}[1][]{popbase,colframe=poppink,before upper={\popboxhead{Attention}{poppink}{#1}}}
\newtcolorbox{experience}[1][]{popbase,colframe=popblue,colback=popblueL,before upper={\popboxhead{Expérience}{popblue}{#1}}}
% « Le savais-tu ? » : carte violette pleine (contexte, culture, Cameroun)
\newtcolorbox{savaistu}[1][]{popbase,colback=poppurple,colframe=popink,
  fontupper=\popsemi\fontsize{10.4}{14.2}\selectfont\color{white},
  before upper={\poppill{Le savais-tu\,?}{white}{poppurple}\ifx\relax#1\relax\else\hspace{6pt}{\popxbold\color{white}#1}\fi\par\vspace{7pt}}}
% Exercice résolu : énoncé en haut, \tcblower puis la solution sur fond crème
\newcounter{popexo}\newcounter{popexr}
\newtcolorbox{exoresolu}[1][]{popbase,colframe=poporange,colbacklower=popcream,
  segmentation style={popink,line width=1.2pt,dashed},
  before upper={\stepcounter{popexr}\popboxhead{Exercice résolu \thepopexr}{poporange}{#1}},
  before lower={\poppill{Solution}{popink}{popcream}\par\vspace{6pt}}}
% Exercice (non résolu en place) — la correction est dans la partie Corrigés
\newtcolorbox{exercice}[1][]{popbase,colframe=popink,
  before upper={\stepcounter{popexo}\popboxhead{Exercice \thepopexo}{popink}{#1}}}
% Corrigé
\newtcolorbox{corrige}[1][]{popbase,colframe=popgreen,colback=popgreenL,
  before upper={\popboxhead{Corrigé}{popgreen}{#1}}}
% Objectifs / prérequis / bilan
\newtcolorbox{objectifs}{popbase,colframe=popink,colback=poporangeL,
  before upper={\popboxhead{Objectifs de la leçon}{popink}{}}}
\newtcolorbox{prerequis}{popbase,colframe=popink,colback=popblueL,
  before upper={\popboxhead{Prérequis}{popblue}{}}}
\newtcolorbox{bilan}{popbase,colframe=popink,colback=white,boxrule=2.8pt,
  before upper={\popboxhead{Fiche bilan}{poporange}{L'essentiel à connaître pour l'examen}}}
% Figure encadrée avec légende
\newtcolorbox{popfigure}[1][]{enhanced,breakable=false,colback=white,colframe=popline,boxrule=1.4pt,arc=8pt,
  left=10pt,right=10pt,top=10pt,bottom=8pt,before skip=10pt,after skip=12pt,halign=center,
  lower separated=false,halign lower=center,
  fontlower=\popsemi\fontsize{9}{11.5}\selectfont\color{popmuted},
  #1}
\newcommand{\legende}[1]{\par\vspace{6pt}{\popsemi\fontsize{9}{11.5}\selectfont\color{popmuted}#1}}

% ============================================================
% Signature OnBuch+
% ============================================================
\newcommand{\poplogoinline}[1]{{\poplogo\fontsize{#1}{#1}\selectfont\color{popink}OnBuch\color{poporange}+}}
\newcommand{\signatureonbuch}{%
  \begin{tcolorbox}[enhanced,colback=popink,colframe=popink,boxrule=2.2pt,arc=10pt,
      drop shadow=poporange,left=14pt,right=14pt,top=10pt,bottom=10pt,before skip=0pt,after skip=0pt]
    \tikz[baseline=-0.7ex]\node[circle,fill=popgreen,draw=popcream,line width=1.3pt,minimum size=22pt,inner sep=0]{\popcheck{white}};%
    \hspace{9pt}\begin{minipage}[c]{0.62\linewidth}
      {\popxbold\fontsize{12}{13}\selectfont\color{popcream}Signé\ \textbullet\ {\poplogo OnBuch\color{poporange}+}}\par\vspace{2pt}
      {\raggedright\popsemi\fontsize{8}{10.5}\selectfont\color{popline}\MakeUppercase{Équipe pédagogique OnBuch+}\\\MakeUppercase{Conforme au programme officiel MINESEC}\par}
    \end{minipage}\hfill
    {\poplogo\fontsize{22}{22}\selectfont\color{popcream}OnBuch\color{poporange}+}
  \end{tcolorbox}%
}

% ============================================================
% Page de garde
% #1 titre · #2 matière · #3 niveau · #4 module · #5 n° leçon · #6 durée ·
% #7 description / objectifs (texte ou itemize)
% ============================================================
\newcommand{\pagegarde}[7]{%
  \thispagestyle{empty}
  \newgeometry{a4paper,margin=1.7cm,top=1.6cm,bottom=1.4cm}%
  \begin{tikzpicture}[remember picture,overlay]
    \fill[poporange!18] ([xshift=-3.2cm,yshift=-3.6cm]current page.north east) circle (4.6cm);
    \fill[poppurple!14] ([xshift=2.4cm,yshift=7.6cm]current page.south west) circle (3.8cm);
  \end{tikzpicture}%
  {\poplogo\fontsize{30}{30}\selectfont\color{popink}OnBuch\color{poporange}+}\hfill
  \raisebox{0.6ex}{\poppill{Cours signé \textbullet\ Premium}{popink}{popcream}}\par
  \vspace{1pt}\popsquiggle{poppurple}\par
  \vspace{8pt}
  \begin{tcolorbox}[enhanced,colback=poporange,colframe=popink,boxrule=2.8pt,arc=13pt,
      drop shadow=popink,left=18pt,right=18pt,top=12pt,bottom=13pt,after skip=10pt]
    \poppill{#2 \textbullet\ #3}{popink}{popcream}\hspace{5pt}\poppill{#4}{white}{popink}\par\vspace{8pt}
    {\popdisplay\fontsize{14}{14}\selectfont\color{popink}LEÇON #5}\par\vspace{4pt}
    {\raggedright\hyphenpenalty=10000\exhyphenpenalty=10000\popdisplay\fontsize{25}{27}\selectfont\color{popcream}\MakeUppercase{#1}\par}
    \vspace{8pt}
    {\popxbold\fontsize{9.5}{9.5}\selectfont\color{popink}\MakeUppercase{Durée conseillée : #6}}
  \end{tcolorbox}
  \begin{tcolorbox}[enhanced,colback=white,colframe=popink,boxrule=2.2pt,arc=10pt,
      drop shadow=popink,left=16pt,right=16pt,top=10pt,bottom=10pt,after skip=8pt]
    \poppill{Dans cette leçon}{poppurple}{white}\par\vspace{5pt}
    \popsemi\fontsize{9.3}{12.6}\selectfont\color{popink}#7
  \end{tcolorbox}
  \noindent\begin{minipage}[t]{0.487\linewidth}
    \begin{tcolorbox}[enhanced,colback=popgreen,colframe=popink,boxrule=2.2pt,arc=10pt,
        drop shadow=popink,left=11pt,right=11pt,top=10pt,bottom=10pt,height=3.1cm,valign=center]
      \tcbox[colback=white,colframe=popink,boxrule=1.3pt,arc=5pt,boxsep=0pt,nobeforeafter,
        left=3pt,right=3pt,top=3pt,bottom=3pt,valign=center]{\qrcode[height=1.9cm]{https://play.google.com/store/apps/details?id=com.luvvix.onbuch&hl=fr}}%
      \hspace{8pt}\begin{minipage}[c]{0.5\linewidth}
        {\popdisplay\fontsize{11}{12.5}\selectfont\color{white}\MakeUppercase{Télécharge OnBuch}}\par\vspace{4pt}
        {\raggedright\popsemi\fontsize{8.4}{11}\selectfont\color{white}Gratuit \textbullet\ Google Play\\Tuteur IA, annales, concours, résultats\par}
      \end{minipage}
    \end{tcolorbox}
  \end{minipage}\hfill
  \begin{minipage}[t]{0.487\linewidth}
    \begin{tcolorbox}[enhanced,colback=poppurple,colframe=popink,boxrule=2.2pt,arc=10pt,
        drop shadow=popink,left=11pt,right=11pt,top=10pt,bottom=10pt,height=3.1cm,valign=center]
      \tikz[baseline=-0.7ex]\node[circle,fill=white,draw=popink,line width=1.3pt,minimum size=1.6cm,inner sep=0]{\popdisplay\fontsize{24}{24}\selectfont\color{poppurple}+};%
      \hspace{8pt}\begin{minipage}[c]{0.6\linewidth}
        {\popdisplay\fontsize{11}{12.5}\selectfont\color{white}\MakeUppercase{Passe à OnBuch+}}\par\vspace{4pt}
        {\raggedright\popsemi\fontsize{8.4}{11}\selectfont\color{white}Tous les cours signés, Tuteur IA illimité, fiches PDF — onglet OnBuch+ de l'app\par}
      \end{minipage}
    \end{tcolorbox}
  \end{minipage}
  \vspace{8pt}
  \popcontenu
  \vfill
  \signatureonbuch
  \restoregeometry
  \newpage
}

% Rangée « Ce cours contient » (page de garde)
\newcommand{\poptile}[3]{% #1 couleur #2 gros texte #3 légende
  \begin{tcolorbox}[enhanced,colback=#1,colframe=popink,boxrule=2pt,arc=9pt,shadow={2.5pt}{-2.5pt}{0pt}{popink},
      left=6pt,right=6pt,top=9pt,bottom=9pt,halign=center,valign=center,height=1.75cm,width=\linewidth,nobeforeafter]
    {\popdisplay\fontsize{12.5}{13}\selectfont\color{white}\MakeUppercase{#2}}\par\vspace{3pt}
    {\centering\popsemi\fontsize{7.8}{9.5}\selectfont\color{white}#3\par}
  \end{tcolorbox}}
\newcommand{\popcontenu}{%
  \par\noindent{\popxboldls\fontsize{7.5}{7.5}\selectfont\color{popmuted}CE COURS SIGNÉ CONTIENT}\par\vspace{7pt}
  \noindent\begin{minipage}[t]{0.235\linewidth}\poptile{popblue}{Cours}{complet, illustré, pas à pas}\end{minipage}\hfill
  \begin{minipage}[t]{0.235\linewidth}\poptile{poporange}{Méthodes}{et exercices résolus}\end{minipage}\hfill
  \begin{minipage}[t]{0.235\linewidth}\poptile{poppink}{Exercices}{type examen + corrigés}\end{minipage}\hfill
  \begin{minipage}[t]{0.235\linewidth}\poptile{popgreen}{Fiche bilan}{l'essentiel à retenir}\end{minipage}\par}

% Sommaire
\newtcolorbox{sommaire}{enhanced,colback=white,colframe=popink,boxrule=2.2pt,arc=10pt,
  drop shadow=popink,left=18pt,right=18pt,top=13pt,bottom=13pt,before skip=6pt,after skip=20pt,
  before upper={\poppill{Sommaire}{poppurple}{white}\par\vspace{9pt}\popsemi\fontsize{10.6}{14.5}\selectfont\color{popink}}}

% Signature de fin de cours — poussée tout en bas de la dernière page
\newcommand{\findecours}{%
  \par\vfill\nopagebreak
  \begin{center}\popsquiggle{poporange}\end{center}\vspace{4pt}
  \signatureonbuch
}
\AtBeginDocument{\def\llbracket{\mathopen{[\mkern-3.5mu[}}\def\rrbracket{\mathclose{]\mkern-3.5mu]}}} % intervalles d'entiers (glyphe ⟦ absent de la police)
% Glyphes absents de Fira Math : redéfinis à partir de glyphes disponibles
\AtBeginDocument{%
  \def\setminus{\mathbin{\backslash}}\let\smallsetminus\setminus
  \def\vdots{\mathord{\vbox{\baselineskip3.2pt\lineskiplimit0pt\kern4pt\hbox{.}\hbox{.}\hbox{.}}}}%
  \def\ddots{\mathinner{\mkern1mu\raise6.5pt\hbox{.}\mkern2mu\raise3.6pt\hbox{.}\mkern2mu\raise0.7pt\hbox{.}\mkern1mu}}%
  \def\triangle{\mathord{\scalebox{0.9}{$\symup{\Delta}$}}}%
  \def\nabla{\mathord{\raisebox{\depth}{\scalebox{1}[-1]{$\symup{\Delta}$}}}}%
  \def\checkmark{\popcheck{popdark}}%
  \def\star{\mathbin{\ast}}\def\bigstar{\mathord{\ast}}%
  \def\diamond{\mathbin{\scalebox{0.6}{$\Diamond$}}}%
  \def\bigcup{\mathop{\mathchoice{\raisebox{-0.3em}{\scalebox{1.7}{$\cup$}}}{\scalebox{1.25}{$\cup$}}{\cup}{\cup}}\displaylimits}%
  \def\bigcap{\mathop{\mathchoice{\raisebox{-0.3em}{\scalebox{1.7}{$\cap$}}}{\scalebox{1.25}{$\cap$}}{\cap}{\cap}}\displaylimits}%
  \def\lesseqqgtr{\mathrel{\lessgtr}}\def\bigoplus{\mathop{\oplus}\displaylimits}%
}
\providecommand{\ssi}{si et seulement si} % abréviation fréquente
\AtBeginDocument{\providecommand{\wideparen}{\overparen}} % arc de cercle

%% FIGFIX-BEGIN v4 — correctifs globaux des figures (docs/onbuch-generation/tools/figfix)
\usepackage{adjustbox}\usepackage{etoolbox}
% toute figure TikZ/pgfplots plus large que la ligne est réduite à la largeur disponible
\BeforeBeginEnvironment{tikzpicture}{\begin{adjustbox}{max width=\linewidth}}
\AfterEndEnvironment{tikzpicture}{\end{adjustbox}}
% légendes pgfplots lisibles par-dessus les courbes
\pgfplotsset{every axis/.append style={legend style={fill=white,fill opacity=0.92,text opacity=1,draw=black!15,inner sep=3pt}}}
% étiquettes d'arêtes (syntaxe edge["..."]) avec fond blanc
\tikzset{every edge quotes/.append style={fill=white,inner sep=1.2pt}}
%% FIGFIX-END
\begin{document}

\pagegarde{Le courant alternatif}{PCT}{3ème}{Module 2 : Électricité}{N12}{1 séance (1 h chacune)}{Cette leçon fait découvrir le courant alternatif, ses caractéristiques et sa différence avec le courant continu. Elle explique ensuite comment l'énergie électrique est produite dans les centrales, transportée par les lignes haute tension, transformée puis distribuée aux usagers, en insistant sur les règles de sécurité.
\vspace{4pt}
\begin{multicols}{2}\small
\begin{itemize}
\item À la fin de cette leçon, je sais définir le courant alternatif et le distinguer du courant continu.
\item À la fin de cette leçon, je sais citer les grandeurs caractéristiques d'une tension alternative (période, fréquence, valeur maximale, valeur efficace).
\item À la fin de cette leçon, je sais interpréter l'oscillogramme d'une tension alternative et en déduire la période et la fréquence.
\item À la fin de cette leçon, je sais décrire le principe de production de l'énergie électrique dans une centrale (hydraulique, thermique).
\item À la fin de cette leçon, je sais expliquer le rôle du transformateur et justifier le transport de l'électricité en haute tension.
\item À la fin de cette leçon, je sais décrire le trajet de l'énergie électrique de la centrale jusqu'à l'installation domestique (schéma de la chaîne ENEO).
\end{itemize}\end{multicols}}

\begin{sommaire}
\begin{multicols}{2}
\begin{enumerate}
\item Courant continu et courant alternatif
\item Caractéristiques d'une tension alternative sinusoïdale
\item Production de l'énergie électrique : les centrales
\item Transport et distribution : lignes haute tension et transformateurs
\item De la centrale à l'habitation : la chaîne ENEO et le compteur
\item Sécurité avec le courant électrique domestique
\item Activité d'intégration
\item Méthodes et exercices résolus
\item Exercices
\item Corrigés des exercices
\item Fiche bilan
\end{enumerate}
\end{multicols}
\end{sommaire}

\begin{objectifs}
\begin{itemize}
\item À la fin de cette leçon, je sais définir le courant alternatif et le distinguer du courant continu.
\item À la fin de cette leçon, je sais citer les grandeurs caractéristiques d'une tension alternative (période, fréquence, valeur maximale, valeur efficace).
\item À la fin de cette leçon, je sais interpréter l'oscillogramme d'une tension alternative et en déduire la période et la fréquence.
\item À la fin de cette leçon, je sais décrire le principe de production de l'énergie électrique dans une centrale (hydraulique, thermique).
\item À la fin de cette leçon, je sais expliquer le rôle du transformateur et justifier le transport de l'électricité en haute tension.
\item À la fin de cette leçon, je sais décrire le trajet de l'énergie électrique de la centrale jusqu'à l'installation domestique (schéma de la chaîne ENEO).
\item À la fin de cette leçon, je sais appliquer les consignes de sécurité liées à l'utilisation du courant domestique.
\end{itemize}
\end{objectifs}

\begin{prerequis}
\begin{itemize}
\item Le courant continu : générateur (pile), intensité et tension continues (leçons de 4ème et début de 3ème)
\item Notions de tension électrique et d'intensité, unités volt (V) et ampère (A)
\item Le circuit électrique simple et les effets du courant (effet thermique, magnétique)
\item Notion d'énergie et de puissance électrique (P = U × I)
\item L'aimantation et l'action d'un aimant sur une bobine (induction, vue en technologie/physique)
\end{itemize}
\end{prerequis}

\newpage

\coursec{Courant continu et courant alternatif}

À Yaoundé comme à Ebolowa, dès que le soleil se couche, les maisons s'illuminent grâce au réseau ENEO. Pourtant, l'électricité qui arrive dans nos prises n'est pas celle qui sort de la pile de ta lampe torche : elle change de sens des milliers de fois par seconde ! Comment cette électricité « dansante » est-elle produite à Songloulou ou à Memve'ele ? Pourquoi voyage-t-elle sous des tensions de plusieurs centaines de milliers de volts sur les lignes haute tension, alors que nos prises n'affichent que 220 V ? Et pourquoi faut-il tant de prudence avec elle ?

Toutes ces questions tournent autour d'une idée simple mais fondamentale : il existe \cle{deux grandes familles de courant électrique}, le courant \cle{continu} et le courant \cle{alternatif}. Dans cette première partie, nous allons apprendre à les distinguer, à les observer et à les reconnaître dans la vie quotidienne.

\courssub{Rappel : le courant continu}

Depuis la classe de 4ème, tu as l'habitude d'alimenter tes circuits avec une \cle{pile} ou un \cle{accumulateur}. Reprenons les idées essentielles.

Quand tu relies une lampe aux deux bornes d'une pile plate de \SI{4,5}{V}, le courant sort toujours par la borne $+$ de la pile, traverse la lampe, et revient par la borne $-$. Il circule \textbf{toujours dans le même sens}, sans jamais s'inverser. De plus, tant que la pile n'est pas usée, la tension entre ses bornes garde \textbf{la même valeur au cours du temps} : elle est constante.

\begin{definition}[Courant continu]
Un \cle{courant continu} est un courant qui circule \textbf{toujours dans le même sens} dans un circuit. La \cle{tension continue} d'un générateur (pile, accumulateur) garde \textbf{une valeur constante} au cours du temps.
\end{definition}

\begin{aretenir}[Notation du courant continu]
Le courant continu est noté \cle{DC} (de l'anglais \emph{Direct Current}) ou symbolisé par le signe \cle{=} sur les appareils électriques. Exemples : une pile marquée « 1,5 V DC », un chargeur de téléphone qui affiche « Sortie : 5 V $\equiv$ ».
\end{aretenir}

Si l'on représentait la tension d'une pile en fonction du temps, on obtiendrait une simple droite horizontale : la tension ne bouge pas.

\begin{popfigure}
\begin{center}
\begin{tikzpicture}
\begin{axis}[
    width=0.85\linewidth, height=5.5cm,
    xlabel={Temps $t$}, ylabel={Tension $U$},
    xmin=0, xmax=10, ymin=0, ymax=6,
    xtick=\empty, ytick={4.5},
    yticklabels={\SI{4,5}{V}},
    axis lines=left,
    grid=major, grid style={popline!40},
]
\addplot[popcurve, domain=0:10, samples=50]{4.5};
\draw[popmuted, dashed] (axis cs:0,4.5) -- (axis cs:0,0);
\node[popblue, anchor=south west] at (axis cs:5.2,4.6) {\small $U$ reste constante};
\end{axis}
\end{tikzpicture}
\end{center}
\legende{Tension continue d'une pile de \SI{4,5}{V} en fonction du temps : la courbe est une droite horizontale. La tension garde la même valeur, et le courant circule toujours dans le même sens.}
\end{popfigure}

\courssub{Observation expérimentale : une lampe, deux sources}

\begin{experience}[La lampe brille-t-elle de la même façon ?]
\textbf{Matériel} : une lampe de faible puissance adaptée (ex. \SI{6}{V} \SI{0,5}{A}), une pile plate de \SI{4,5}{V}, un montage sécurisé branché sur le secteur \SI{220}{V} (réalisé par le professeur, jamais par l'élève seul), des fils de connexion.

\textbf{Protocole} :
\begin{enumerate}
\item Brancher la lampe sur la pile. Observer son éclat.
\item Brancher une lampe adaptée sur le secteur, via le montage sécurisé du professeur. Observer son éclat.
\end{enumerate}

\textbf{Observations} : dans les deux cas, la lampe brille de façon apparemment identique et stable. Rien, à l'œil nu, ne permet de deviner une différence.

\textbf{Interprétation} : pourtant, la \emph{nature} du courant n'est pas la même. Avec la pile, le courant circule toujours dans le même sens : c'est un courant \textbf{continu}. Avec le secteur ENEO, le courant change de sens de très nombreuses fois par seconde : c'est un courant \textbf{alternatif}. Nos yeux sont simplement incapables de percevoir ces changements, car ils sont beaucoup trop rapides.
\end{experience}

Cette expérience illustre une démarche d'investigation classique : l'observation directe ne suffit pas toujours. Pour « voir » la différence, il faudrait un appareil capable de suivre les variations rapides de la tension, comme l'\cle{oscilloscope} (que tu découvriras au lycée). En attendant, retenons les définitions.

\courssub{Le courant alternatif : définitions}

\begin{definition}[Courant alternatif]
Un \cle{courant alternatif} est un courant qui \textbf{change périodiquement de sens} au cours du temps : il circule un instant dans un sens, puis dans le sens opposé, et recommence ainsi de manière régulière.
\end{definition}

\begin{definition}[Tension alternative]
Une \cle{tension alternative} est une tension qui \textbf{change périodiquement de signe} au cours du temps : elle est tantôt positive, tantôt négative.
\end{definition}

Le courant distribué par ENEO n'est pas alternatif n'importe comment : sa tension varie en suivant une courbe mathématique appelée \cle{sinusoïde}. On parle de \cle{courant alternatif sinusoïdal}. La courbe monte doucement jusqu'à un maximum, redescend en passant par zéro, descend jusqu'à un minimum, puis remonte : c'est un cycle complet, qui se répète indéfiniment.

\begin{popfigure}
\begin{center}
\begin{tikzpicture}
\begin{axis}[
    width=0.9\linewidth, height=6.5cm,
    xlabel={Temps $t$}, ylabel={Tension $u(t)$},
    xmin=0, xmax=2.2, ymin=-1.6, ymax=1.8,
    xtick={1}, xticklabels={$T$},
    ytick={1,-1}, yticklabels={$U_{\text{max}}$, $-U_{\text{max}}$},
    axis lines=middle,
    grid=major, grid style={popline!40},
]
\addplot[popcurve, domain=0:2.2, samples=200]{sin(360*x)};
\draw[poporange, thick, <->] (axis cs:0.25,1.25) -- (axis cs:1.25,1.25) node[midway, above, poporange]{\small période $T$};
\draw[poporange, dashed] (axis cs:0.25,1) -- (axis cs:0.25,1.25);
\draw[poporange, dashed] (axis cs:1.25,1) -- (axis cs:1.25,1.25);
\draw[popgreen, thick, <->] (axis cs:1.55,0) -- (axis cs:1.55,1) node[midway, right, popgreen]{\small $U_{\text{max}}$};
\end{axis}
\end{tikzpicture}
\end{center}
\legende{Tension alternative sinusoïdale $u(t)$ : la tension change périodiquement de signe. La \textbf{période} $T$ est la durée d'un cycle complet ; $U_{\text{max}}$ est la valeur maximale atteinte par la tension.}
\end{popfigure}

\begin{aretenir}[Le courant du réseau ENEO]
Le courant distribué par ENEO au Cameroun est un courant \cle{alternatif sinusoïdal} :
\begin{itemize}
\item de \cle{tension efficace} $U = \SI{220}{V}$ (la valeur indiquée sur les appareils et les factures) ;
\item de \cle{fréquence} $f = \SI{50}{Hz}$, c'est-à-dire que le courant effectue \textbf{50 cycles complets par seconde}.
\end{itemize}
La période vaut donc $T = \dfrac{1}{f} = \dfrac{1}{50} = \SI{0,02}{s}$ : chaque cycle ne dure que deux centièmes de seconde !
\end{aretenir}

\begin{aretenir}[Notation du courant alternatif]
Le courant alternatif est noté \cle{AC} (de l'anglais \emph{Alternating Current}) ou symbolisé par le signe \cle{$\sim$} sur les appareils. Exemple : « Entrée : 220 V $\sim$ 50 Hz » sur l'étiquette d'un téléviseur.
\end{aretenir}

\begin{attention}[La lampe alimentée en alternatif ne clignote pas !]
\textbf{Erreur fréquente} : croire qu'une lampe branchée sur le secteur « clignote » puisque le courant change de sens. C'est faux ! Le courant change de sens 100 fois par seconde (deux fois par cycle), soit une variation toutes les \SI{0,01}{s}. C'est beaucoup trop rapide pour être perçu par l'œil humain, qui ne distingue pas plus d'une vingtaine d'images par seconde. De plus, le filament de la lampe n'a pas le temps de refroidir entre deux passages du courant. La lampe nous paraît donc briller de façon \textbf{continue et stable}.
\end{attention}

\courssub{Les sources de courant dans la vie courante}

Sauras-tu maintenant classer les sources d'électricité qui t'entourent ? Voici les principales.

\begin{center}
\begin{tabularx}{\linewidth}{|l|X|X|}
\hline
\rowcolor{popblueL} \textbf{Source} & \textbf{Nature du courant} & \textbf{Symbole} \\
\hline
Pile (lampe torche, télécommande) & Continu & DC ou $=$ \\
\hline
Accumulateur (batterie de téléphone, de voiture) & Continu & DC ou $=$ \\
\hline
Prise murale ENEO & Alternatif sinusoïdal 220 V, 50 Hz & AC ou $\sim$ \\
\hline
Alternateur de centrale (Songloulou, Memve'ele) & Alternatif & AC ou $\sim$ \\
\hline
Dynamo de vélo (éclairage) & Alternatif & AC ou $\sim$ \\
\hline
Groupe électrogène & Alternatif & AC ou $\sim$ \\
\hline
\end{tabularx}
\end{center}

Remarque bien la logique : les sources \textbf{chimiques} (piles, accumulateurs), où le courant naît d'une réaction chimique, fournissent du courant \textbf{continu}. Les sources \textbf{mécaniques}, où une bobine tourne devant un aimant (alternateur, dynamo, groupe électrogène), fournissent du courant \textbf{alternatif} : à chaque demi-tour de la bobine, le sens du courant s'inverse naturellement. C'est pourquoi presque toute l'électricité produite dans le monde est alternative.

\begin{exemplebox}[Deux tensions à ne pas confondre]
La prise murale de ta maison, au Cameroun, délivre une tension \textbf{alternative} de valeur efficace \SI{220}{V} et de fréquence \SI{50}{Hz}. Une pile plate de \SI{4,5}{V} délivre une tension \textbf{continue} de \SI{4,5}{V}. Même symbole $U$, mais deux comportements totalement différents : l'une oscille 50 fois par seconde entre environ $+311$ V et $-311$ V, l'autre reste fixe à $+4,5$ V.
\end{exemplebox}

\begin{exoresolu}[Continu ou alternatif ? À toi de juger]
On relève les indications suivantes sur quatre appareils : (a) « 1,5 V DC » sur une pile ; (b) « 220 V $\sim$ 50 Hz » sur un fer à repasser ; (c) « Sortie : 12 V $\equiv$ » sur un adaptateur ; (d) « 6 V AC » sur la dynamo d'un vélo. Pour chacun, préciser la nature du courant utilisé et justifier.
\tcblower
\textbf{Solution.}
\begin{enumerate}
\item[(a)] Le symbole DC indique un courant \textbf{continu} : la pile fournit une tension constante de \SI{1,5}{V}.
\item[(b)] Le symbole $\sim$ indique un courant \textbf{alternatif} : le fer à repasser est conçu pour le secteur ENEO (\SI{220}{V}, \SI{50}{Hz}).
\item[(c)] Le symbole $\equiv$ (équivalent à DC) indique un courant \textbf{continu} : l'adaptateur transforme le courant alternatif du secteur en courant continu de \SI{12}{V} pour l'appareil.
\item[(d)] Le symbole AC indique un courant \textbf{alternatif} : la rotation de la roue fait tourner un aimant devant une bobine, ce qui produit naturellement un courant qui change de sens.
\end{enumerate}
\textbf{Méthode} : repérer d'abord le symbole (DC, $=$, $\equiv$ pour le continu ; AC, $\sim$ pour l'alternatif), puis conclure.
\end{exoresolu}

\begin{savaistu}[La « guerre des courants » : Tesla contre Edison]
À la fin du XIX\textsuperscript{e} siècle, deux inventeurs s'affrontent aux États-Unis. Thomas \cle{Edison} défend le courant continu, qu'il utilise dans ses premières centrales. Nikola \cle{Tesla}, lui, défend le courant alternatif. Le verdict de la technique est sans appel : le courant alternatif l'emporte, car sa tension peut être facilement élevée ou abaissée grâce au \cle{transformateur}. Or, transporter l'électricité sous très haute tension permet de la faire voyager sur des centaines de kilomètres avec très peu de pertes — impossible avec le courant continu de l'époque. C'est grâce à cette victoire de l'alternatif que l'électricité produite à Songloulou peut alimenter aujourd'hui des villes éloignées du barrage. Nous étudierons le transformateur plus loin dans cette leçon.
\end{savaistu}

\begin{attention}[Pièges de rédaction au BEPC]
\begin{itemize}
\item Ne dis pas « le courant alternatif s'arrête et repart » : il \textbf{change de sens}, c'est différent. Il circule en permanence, tantôt dans un sens, tantôt dans l'autre.
\item Ne confonds pas \textbf{tension maximale} $U_{\text{max}}$ et \textbf{tension efficace} $U$ : la valeur \SI{220}{V} affichée sur les appareils est la tension \emph{efficace}, pas la valeur maximale de la sinusoïde.
\item « 50 Hz » signifie \textbf{50 cycles par seconde}, et non « 50 changements de sens par seconde » : à chaque cycle, le courant change de sens \emph{deux} fois.
\item Une pile usée voit sa tension \emph{diminuer lentement} : ce n'est pas pour autant un courant alternatif, car le sens du courant ne s'inverse jamais.
\end{itemize}
\end{attention}


\coursec{Caractéristiques d'une tension alternative sinusoïdale}

Dans la section précédente, tu as découvert que le courant délivré par une pile est \cle{continu} : il circule toujours dans le même sens. En revanche, le courant fourni par les prises de courant ENEO est \cle{alternatif} : il change de sens de façon régulière, des dizaines de fois par seconde. Mais comment le vérifier ? Le courant est invisible, et un simple voltmètre ne nous dit pas tout. Il nous faut un instrument capable de \emph{dessiner} la tension au fil du temps : l'\cle{oscilloscope}. Cette section t'apprend à lire le dessin qu'il produit et à en extraire quatre grandeurs essentielles : la période, la fréquence, la tension maximale et la tension efficace.

\courssub{Observation à l'oscilloscope : la sinusoïde}

\begin{experience}[Observer une tension alternative à l'oscilloscope]
\textbf{Matériel :} un oscilloscope, un générateur basse fréquence (GBF) réglé sur « tension alternative sinusoïdale », une pile plate, des fils de connexion.

\textbf{Protocole :}
\begin{enumerate}
\item Brancher la sortie du GBF sur la voie d'entrée de l'oscilloscope.
\item Régler la sensibilité verticale et le balayage pour obtenir une courbe lisible.
\item Observer l'écran, puis recommencer en branchant une pile à la place du GBF.
\end{enumerate}

\textbf{Observations :}
\begin{itemize}
\item Avec la pile, on observe une \emph{droite horizontale} : la tension est constante au cours du temps.
\item Avec le GBF, on observe une courbe régulière \emph{en forme de vague}, qui monte, redescend, passe sous l'axe horizontal, puis recommence à l'identique.
\end{itemize}

\textbf{Interprétation :} la tension délivrée par le GBF varie au cours du temps : elle prend tour à tour des valeurs positives et négatives, et le même motif se répète indéfiniment. Cette courbe mathématique particulière s'appelle une \cle{sinusoïde}. La tension du secteur ENEO présente la même allure.
\end{experience}

\begin{definition}[Tension alternative sinusoïdale]
Une tension est dite \cle{alternative sinusoïdale} lorsque sa représentation graphique en fonction du temps est une courbe en forme de vague appelée \cle{sinusoïde}, qui change de signe de façon régulière et dont un même motif se répète indéfiniment.
\end{definition}

\begin{popfigure}
\begin{center}
\begin{tikzpicture}[scale=0.9]
% Corps de l'oscilloscope
\draw[fill=popdark, rounded corners=4pt] (0,0) rectangle (8,5);
% Ecran
\draw[fill=popcream, rounded corners=2pt] (0.5,1.2) rectangle (5.3,4.5);
% Quadrillage ecran
\draw[popline, step=0.6] (0.5,1.2) grid (5.3,4.5);
\draw[popmuted] (0.5,2.85) -- (5.3,2.85);
% Sinusoide sur l'ecran
\draw[popblue, very thick, domain=0.5:5.3, samples=100] plot (\x, {2.85 + 1.1*sin((\x-0.5)*260)});
% Boutons
\draw[fill=popmuted] (5.9,3.9) circle (0.45);
\draw[fill=popmuted] (7.1,3.9) circle (0.45);
\draw[fill=popmuted] (5.9,2.6) circle (0.45);
\draw[fill=popmuted] (7.1,2.6) circle (0.45);
\draw[fill=popmuted] (5.9,1.3) circle (0.45);
\draw[fill=popmuted] (7.1,1.3) circle (0.45);
% Traits de repere sur boutons
\draw[white, thick] (5.9,3.9) -- (5.9,4.25);
\draw[white, thick] (7.1,3.9) -- (7.35,3.9);
\draw[white, thick] (5.9,2.6) -- (6.2,2.9);
\draw[white, thick] (7.1,2.6) -- (7.1,2.95);
\draw[white, thick] (5.9,1.3) -- (5.6,1.3);
\draw[white, thick] (7.1,1.3) -- (7.35,1.55);
% Etiquettes
\node[white, font=\small] at (5.9,4.75) {V/div};
\node[white, font=\small] at (7.1,4.75) {ms/div};
\node[white, font=\footnotesize] at (6.5,0.55) {entrée};
% Fiche d'entree
\draw[fill=poporange] (6.2,0.1) rectangle (6.8,0.35);
\end{tikzpicture}
\end{center}
\legende{Schéma d'un oscilloscope. L'écran quadrillé affiche la sinusoïde. Les deux réglages essentiels : la \textbf{sensibilité verticale} (en V/div), qui fixe la valeur en volts d'une division verticale, et le \textbf{balayage} (en ms/div), qui fixe la durée d'une division horizontale.}
\end{popfigure}

\begin{attention}[L'oscilloscope ne se branche pas n'importe comment]
Au laboratoire, on observe la tension d'un GBF ou d'un montage basse tension. \textbf{On ne branche jamais un oscilloscope directement sur une prise 220 V} sans dispositif adapté : danger de mort. La tension du secteur est dangereuse ; seul le professeur ou un technicien qualifié réalise cette manipulation.
\end{attention}

\courssub{La période T : le temps d'un motif}

Regarde la sinusoïde : elle est faite d'un même motif (une « bosse » vers le haut suivie d'une « bosse » vers le bas) qui se répète sans fin. Le temps que met ce motif à se reproduire est une grandeur fondamentale.

\begin{definition}[Période]
La \cle{période}, notée $T$, d'une tension alternative est la durée du plus petit motif qui se répète. Son unité est la \cle{seconde} (symbole : s). On utilise souvent la milliseconde : $\SI{1}{ms} = \SI{0,001}{s}$.
\end{definition}

Sur l'oscillogramme, la période se mesure \emph{horizontalement} : on compte le nombre de divisions occupées par un motif complet, puis on multiplie par le balayage (la durée d'une division).

\courssub{La fréquence f : le nombre de motifs par seconde}

Si un motif dure $T$ secondes, combien de motifs s'écoulent en une seconde ? La réponse est une division : $1$ divisé par $T$. C'est la fréquence.

\begin{definition}[Fréquence]
La \cle{fréquence}, notée $f$, d'une tension alternative est le nombre de périodes par seconde. Son unité est le \cle{hertz} (symbole : Hz). Elle est liée à la période par la relation :
\[ f = \frac{1}{T} \qquad \text{et donc} \qquad T = \frac{1}{f} \]
avec $T$ en secondes (s) et $f$ en hertz (Hz). $1$ Hz correspond à $1$ période par seconde.
\end{definition}

\begin{exemplebox}[Le réseau ENEO : 50 Hz]
Au Cameroun, la tension du réseau ENEO a une fréquence $f = \SI{50}{Hz}$ : le motif complet se répète $50$ fois par seconde, ce qui correspond à $100$ alternances (le courant change de sens $100$ fois par seconde). Sa période vaut :
\[ T = \frac{1}{f} = \frac{1}{50} = \SI{0,02}{s} = \SI{20}{ms}. \]
C'est beaucoup trop rapide pour que nos yeux le perçoivent : c'est pourquoi une lampe branchée sur le secteur semble briller en continu, alors que le courant qui la traverse s'inverse $100$ fois par seconde.
\end{exemplebox}

\courssub{Tension maximale et tension efficace}

\textbf{La tension maximale $U_{max}$.}\par
Sur la sinusoïde, la tension atteint un sommet (la crête) puis redescend. La hauteur de cette crête, mesurée \emph{verticalement} depuis l'axe horizontal, est la tension maximale.

\begin{definition}[Tension maximale]
La \cle{tension maximale}, notée $U_{max}$, est la valeur maximale atteinte par la tension alternative au cours d'une période. C'est l'\emph{amplitude} de la sinusoïde. Elle s'exprime en volts (V).
\end{definition}

\textbf{La tension efficace $U$.}\par
Quand on branche un voltmètre (en position AC, c'est-à-dire alternatif) aux bornes d'une prise, il n'affiche pas $U_{max}$ : il affiche une valeur plus petite, appelée \cle{tension efficace}. C'est elle qui caractérise l'effet « utile » de la tension : une tension alternative de valeur efficace $U$ produit le même échauffement dans un appareil (fer à repasser, cuisinière) qu'une tension continue de même valeur.

\begin{propriete}[Relation entre tension efficace et tension maximale — admise]
Pour une tension alternative \emph{sinusoïdale}, la tension efficace $U$ et la tension maximale $U_{max}$ sont liées par :
\[ U = \frac{U_{max}}{\sqrt{2}} \qquad \text{ou encore} \qquad U_{max} = U \times \sqrt{2} \]
avec $\sqrt{2} \approx 1,414$. Le voltmètre en position \textbf{AC} mesure toujours la valeur \emph{efficace}. Cette relation est admise au niveau 3ème : aucune démonstration n'est exigible au BEPC.
\end{propriete}

\begin{exemplebox}[Pourquoi « 220 V » cache en réalité 311 V]
Pour le réseau ENEO, la valeur affichée est $U = \SI{220}{V}$ (valeur efficace). La tension maximale vaut alors :
\[ U_{max} = U \times \sqrt{2} = 220 \times 1,414 \approx \SI{311}{V}. \]
La tension de la prise atteint donc en réalité $311$ volts à chaque crête, $50$ fois par seconde !
\end{exemplebox}

\begin{attention}[Ne pas confondre valeur maximale et valeur efficace]
C'est l'erreur la plus fréquente au BEPC :
\begin{itemize}
\item La valeur « $\SI{220}{V}$ » inscrite sur les appareils (fer à repasser, téléviseur, chargeur) est la valeur \cle{efficace}.
\item La valeur lue sur l'oscillogramme (hauteur de la crête) est la valeur \cle{maximale}.
\item Les deux sont liées par $U = \dfrac{U_{max}}{\sqrt{2}}$. Ne mélange jamais les deux dans un calcul !
\end{itemize}
\end{attention}

\begin{popfigure}
\begin{center}
\begin{tikzpicture}
\begin{axis}[
  width=0.95\linewidth, height=6.5cm,
  axis lines=middle,
  xlabel={$t$ (ms)}, ylabel={$u$ (V)},
  xlabel style={at={(ticklabel* cs:1)}, anchor=west},
  ylabel style={at={(ticklabel* cs:1)}, anchor=south},
  xmin=-2, xmax=26, ymin=-380, ymax=420,
  xtick={0,5,10,15,20,25},
  ytick={-311,0,311},
  yticklabels={$-311$, $0$, $311$},
  grid=both, grid style={popline, very thin},
  major grid style={popmuted!50},
  tick label style={font=\small},
]
% Sinusoide : Umax = 311 V, T = 20 ms
\addplot[popcurve, domain=0:25, samples=300] {311*sin(360*x/20)};
% Fleche Umax (verticale)
\draw[<->, poporange, very thick] (axis cs:5,0) -- (axis cs:5,311)
  node[midway, right, font=\small, text=poporange]{$U_{max}$};
% Fleche T (horizontale)
\draw[<->, popgreen, very thick] (axis cs:0,-360) -- (axis cs:20,-360)
  node[midway, below, font=\small, text=popgreen]{$T = \SI{20}{ms}$};
% Marquage d'un motif
\draw[dashed, poppurple] (axis cs:20,-311) -- (axis cs:20,340);
\node[font=\footnotesize, text=poppurple, anchor=south west] at (axis cs:20.3,320) {un motif complet};
\end{axis}
\end{tikzpicture}
\end{center}
\legende{Oscillogramme de la tension du réseau ENEO : sinusoïde de période $T = \SI{20}{ms}$ et de tension maximale $U_{max} \approx \SI{311}{V}$ (soit $U = \SI{220}{V}$ efficaces). Réglages : balayage $\SI{5}{ms/div}$, sensibilité verticale $\SI{100}{V/div}$ : un motif occupe $4$ divisions horizontales et la crête culmine à environ $3,1$ divisions.}
\end{popfigure}

\courssub{Exploiter un oscillogramme : la méthode}

Au BEPC, on te donne souvent un oscillogramme avec deux réglages : le \cle{balayage} (en ms/div) et la \cle{sensibilité verticale} (en V/div). Tu dois en tirer $T$, $f$, $U_{max}$ et $U$.

\begin{methode}[Exploiter un oscillogramme en 4 étapes]
\begin{enumerate}
\item \textbf{Mesurer la période} : compter le nombre $x$ de divisions horizontales occupées par \emph{un motif complet}, puis calculer :
\[ T = x \times \text{balayage} \quad (\text{en ms, à convertir en s}). \]
\item \textbf{Calculer la fréquence} : $f = \dfrac{1}{T}$ avec $T$ en secondes ; $f$ est en hertz.
\item \textbf{Mesurer la tension maximale} : compter le nombre $y$ de divisions verticales entre l'axe horizontal et la crête, puis calculer :
\[ U_{max} = y \times \text{sensibilité verticale} \quad (\text{en V}). \]
\item \textbf{Calculer la tension efficace} : $U = \dfrac{U_{max}}{\sqrt{2}}$.
\end{enumerate}
\end{methode}

\begin{exoresolu}[Reconnaître le réseau ENEO sur un oscillogramme]
\textbf{Énoncé.} Sur l'écran d'un oscilloscope, on observe qu'un motif complet de la sinusoïde occupe $4$ divisions horizontales et que la crête atteint $3,1$ divisions verticales. Les réglages sont : balayage $= \SI{5}{ms/div}$ ; sensibilité verticale $= \SI{100}{V/div}$.
\begin{enumerate}
\item Déterminer la période $T$ puis la fréquence $f$ de cette tension.
\item Déterminer la tension maximale $U_{max}$ puis la tension efficace $U$.
\item Cette tension peut-elle être celle du réseau ENEO ?
\end{enumerate}
\tcblower
\textbf{Solution détaillée.}
\begin{enumerate}
\item \textbf{Période :} un motif occupe $4$ divisions et chaque division vaut $\SI{5}{ms}$ :
\[ T = 4 \times 5 = \SI{20}{ms} = \SI{0,02}{s}. \]
\textbf{Fréquence :}
\[ f = \frac{1}{T} = \frac{1}{0,02} = \SI{50}{Hz}. \]
\item \textbf{Tension maximale :} la crête est à $3,1$ divisions et chaque division vaut $\SI{100}{V}$ :
\[ U_{max} = 3,1 \times 100 = \SI{310}{V}. \]
\textbf{Tension efficace :}
\[ U = \frac{U_{max}}{\sqrt{2}} = \frac{310}{1,414} \approx \SI{219}{V} \approx \SI{220}{V}. \]
\item On trouve $f = \SI{50}{Hz}$ et $U \approx \SI{220}{V}$ : ce sont exactement les caractéristiques du réseau ENEO. C'est donc bien la tension du secteur.
\end{enumerate}
\end{exoresolu}

\begin{attention}[Les trois pièges classiques de l'exercice « oscillogramme »]
\begin{itemize}
\item \textbf{Oublier la conversion} : le balayage est en ms/div, mais la formule $f = \dfrac{1}{T}$ exige $T$ en \emph{secondes}. $\SI{20}{ms} = \SI{0,02}{s}$, pas $20$ s !
\item \textbf{Mesurer une demi-période} : un motif complet, c'est une bosse vers le haut \emph{et} une bosse vers le bas. Si tu ne comptes qu'une bosse, tu trouves une période deux fois trop petite.
\item \textbf{Confondre crête et crête-à-crête} : $U_{max}$ se mesure de l'axe horizontal jusqu'à la crête, \emph{pas} du sommet le plus haut au sommet le plus bas (cela donnerait $2\,U_{max}$).
\end{itemize}
\end{attention}

\begin{aretenir}[L'essentiel de la section]
\begin{itemize}
\item Une tension alternative sinusoïdale est représentée par une \cle{sinusoïde} (courbe en vague) observée à l'\cle{oscilloscope}.
\item La \cle{période} $T$ est la durée d'un motif complet ; elle s'exprime en secondes (s).
\item La \cle{fréquence} $f$ est le nombre de périodes par seconde : $f = \dfrac{1}{T}$, en hertz (Hz).
\item La \cle{tension maximale} $U_{max}$ est la hauteur de la crête ; la \cle{tension efficace} $U$, mesurée par le voltmètre en position AC, vérifie $U = \dfrac{U_{max}}{\sqrt{2}}$.
\item Réseau ENEO : $f = \SI{50}{Hz}$, $T = \SI{20}{ms}$, $U = \SI{220}{V}$, $U_{max} \approx \SI{311}{V}$.
\item Sur un oscillogramme : $T = (\text{divisions horizontales}) \times (\text{balayage})$ et $U_{max} = (\text{divisions verticales}) \times (\text{sensibilité})$.
\end{itemize}
\end{aretenir}

\begin{exercice}[À toi de jouer]
Sur un oscillogramme, un motif complet occupe $2,5$ divisions horizontales avec un balayage de $\SI{2}{ms/div}$, et la crête atteint $2$ divisions verticales avec une sensibilité de $\SI{5}{V/div}$.
\begin{enumerate}
\item Calculer la période $T$ et la fréquence $f$ de cette tension.
\item Calculer la tension maximale $U_{max}$ et la tension efficace $U$.
\end{enumerate}
\end{exercice}

\begin{corrige}[Exercice : exploitation d'un oscillogramme]
\begin{enumerate}
\item Période : $T = 2,5 \times 2 = \SI{5}{ms} = \SI{0,005}{s}$. Fréquence : $f = \dfrac{1}{T} = \dfrac{1}{0,005} = \SI{200}{Hz}$.
\item Tension maximale : $U_{max} = 2 \times 5 = \SI{10}{V}$. Tension efficace : $U = \dfrac{U_{max}}{\sqrt{2}} = \dfrac{10}{1,414} \approx \SI{7,1}{V}$.
\end{enumerate}
\end{corrige}

\begin{savaistu}[Pourquoi 50 Hz ?]
Le choix de $50$ Hz en Afrique et en Europe (contre $60$ Hz en Amérique du Nord) est historique : il date des premières centrales électriques de la fin du XIX\textsuperscript{e} siècle. Cette fréquence est un compromis : assez élevée pour que les lampes ne clignotent pas visiblement, assez basse pour limiter les pertes dans les lignes et les moteurs. Tous les appareils vendus au Cameroun sont conçus pour fonctionner sous $\SI{220}{V}$ -- $\SI{50}{Hz}$ : vérifie toujours cette mention sur la plaque signalétique avant de brancher un appareil importé !
\end{savaistu}

\coursec{Production de l'énergie électrique : les centrales}

Dans la section précédente, nous avons décrit la tension alternative sinusoïdale : sa période, sa fréquence de \SI{50}{Hz}, sa valeur efficace de \SI{220}{V}. Mais une question reste en suspens : \emph{d'où vient cette tension alternative} ? Qui la « fabrique » et comment ? Quand tu branches ta lampe sur la prise, l'électricité a parcouru un long chemin depuis une \cle{centrale électrique}. C'est ce chemin que nous allons découvrir maintenant, en commençant par une expérience étonnante : produire du courant avec... un simple aimant.

\courssub{L'expérience fondatrice : l'induction électromagnétique}

\textbf{Intuition.} Tu connais la dynamo de vélo : quand la petite molette tourne contre le pneu, la lampe s'allume. Personne n'a branché de pile ! Le courant apparaît grâce au \emph{mouvement}. C'est exactement le principe de toutes les centrales électriques du monde.

\begin{experience}[Un aimant, une bobine, et le courant apparaît]
\textbf{Matériel :} une bobine de fil de cuivre (plusieurs centaines de spires), un aimant droit, un galvanomètre (ou une DEL), des fils de connexion.

\textbf{Protocole :}
\begin{enumerate}
\item Relier les deux bornes de la bobine au galvanomètre. Aucun générateur n'est branché.
\item Approcher rapidement le pôle nord de l'aimant de la bobine, puis l'éloigner. Observer l'aiguille.
\item Recommencer en tenant l'aimant \emph{immobile} à l'intérieur de la bobine. Observer.
\item Recommencer en tenant l'aimant fixe et en déplaçant la bobine. Observer.
\end{enumerate}

\textbf{Observations :}
\begin{itemize}
\item Quand l'aimant \emph{entre} dans la bobine, l'aiguille dévie dans un sens ; quand il \emph{sort}, elle dévie dans l'\emph{autre sens}.
\item Aimant immobile, même placé dans la bobine : l'aiguille reste à zéro. \textbf{Pas de mouvement, pas de courant.}
\item C'est le \emph{mouvement relatif} qui compte : déplacer la bobine devant l'aimant fixe produit le même effet.
\item Plus le mouvement est rapide, plus la déviation est grande.
\end{itemize}

\textbf{Interprétation :} le mouvement relatif de l'aimant et de la bobine fait apparaître une \emph{tension} aux bornes de la bobine. Comme cette tension change de sens quand le sens du mouvement change, elle est \emph{alternative}. C'est le phénomène d'\cle{induction électromagnétique}.
\end{experience}

\begin{propriete}[Induction électromagnétique — établie expérimentalement]
Le \cle{mouvement relatif} d'un aimant et d'une bobine produit une \cle{tension alternative} aux bornes de la bobine : c'est le phénomène d'\cle{induction électromagnétique}.
\begin{itemize}
\item Si l'aimant et la bobine sont immobiles l'un par rapport à l'autre, la tension est \textbf{nulle} : \emph{sans mouvement, pas de tension}.
\item Le sens de la tension dépend du sens du mouvement : c'est pourquoi une rotation continue produit une tension \emph{alternative} (qui change de sens périodiquement).
\end{itemize}
\end{propriete}

\begin{popfigure}
\begin{center}
\begin{tikzpicture}[scale=1,>=Stealth]
% Bobine
\draw[popblue,thick] (3.2,-0.5) rectangle (4.6,0.5);
\foreach \x in {3.35,3.55,3.75,3.95,4.15,4.35}{
  \draw[popblue] (\x,-0.5) arc[start angle=-90,end angle=90,x radius=0.06,y radius=0.5];
}
\node[popblue,below] at (3.9,-0.6) {\small Bobine};
% Aimant
\draw[fill=popink] (0.6,-0.5) rectangle (1.35,0.5);
\draw[fill=popmuted!40] (1.35,-0.5) rectangle (2.1,0.5);
\node[white] at (0.97,0) {\small \textbf{S}};
\node at (1.72,0) {\small \textbf{N}};
\node[below] at (1.35,-0.6) {\small Aimant en mouvement};
% Flèche mouvement
\draw[->,very thick,poporange] (2.2,0.9) -- (3.1,0.9) node[midway,above]{\small mouvement};
% Fils vers galvanomètre
\draw[popdark,thick] (3.2,-0.5) -- (3.2,-1.6) -- (5.6,-1.6);
\draw[popdark,thick] (4.6,-0.5) -- (4.6,-1.1) -- (6.6,-1.1) -- (6.6,-1.6);
% Galvanomètre
\draw[fill=popcream,thick,popdark] (6.1,-1.6) circle (0.5);
\node at (6.1,-1.6) {\small \textbf{G}};
\draw[popdark] (6.1,-1.6) -- (6.35,-1.35);
\node[below] at (6.1,-2.2) {\small Galvanomètre};
\end{tikzpicture}
\end{center}
\legende{Expérience d'induction : quand l'aimant se déplace vers la bobine (ou s'en éloigne), le galvanomètre G détecte un courant. Aimant immobile : l'aiguille reste à zéro.}
\end{popfigure}

\begin{attention}[Le piège classique]
Beaucoup d'élèves écrivent : « l'aimant dans la bobine produit du courant ». C'est \textbf{faux} ! Un aimant \emph{immobile}, même placé au cœur de la bobine, ne produit \textbf{aucune} tension. C'est le \cle{mouvement} (relatif) qui est indispensable. Retiens la formule : \emph{pas de mouvement, pas de tension}.
\end{attention}

\courssub{L'alternateur : le cœur de toute centrale}

Pour produire du courant en continu, il faut entretenir un mouvement \emph{permanent}. La solution technique : faire \emph{tourner} un aimant devant une bobine (ou une bobine devant un aimant). La rotation fait passer l'aimant alternativement dans un sens puis dans l'autre : la tension produite est alors \emph{alternative sinusoïdale}, exactement celle que nous avons étudiée à la section précédente. Cette machine s'appelle l'\cle{alternateur}.

\begin{definition}[Alternateur]
Un \cle{alternateur} est un appareil qui convertit l'\cle{énergie mécanique de rotation} en \cle{énergie électrique alternative}, grâce au phénomène d'induction électromagnétique.
\end{definition}

\begin{exemplebox}[De la dynamo de vélo au groupe électrogène]
\begin{itemize}
\item La \emph{dynamo} de vélo est un mini-alternateur : la rotation de la molette fait tourner un aimant devant une bobine.
\item Le \emph{groupe électrogène} que l'on entend dans les quartiers pendant les délestages contient un moteur à essence qui fait tourner un alternateur : il fournit du $220$ V alternatif, comme ENEO.
\item Dans une grande centrale, le principe est identique, mais l'alternateur est gigantesque et c'est une \emph{turbine} qui le fait tourner.
\end{itemize}
\end{exemplebox}

La question devient alors : \emph{qu'est-ce qui fait tourner la turbine ?} La réponse dépend du type de centrale.

\courssub{La centrale hydraulique : la force de l'eau}

\textbf{Observation de la vie courante.} Qui n'a jamais vu la puissance d'un cours d'eau en crue ? L'eau qui tombe ou qui s'écoule possède une énergie considérable. Le Cameroun, avec ses grands fleuves (Sanaga, Ntem, Bénoué), a compris très tôt qu'il pouvait transformer cette énergie en électricité.

\begin{definition}[Centrale hydraulique]
Une \cle{centrale hydraulique} est une centrale qui utilise l'énergie de l'eau (chute ou retenue d'un barrage) pour faire tourner une \cle{turbine} couplée à un \cle{alternateur}.
\end{definition}

\textbf{Fonctionnement, étape par étape :}
\begin{enumerate}
\item Un \emph{barrage} retient l'eau d'un fleuve en hauteur : l'eau accumulée possède une énergie (dite \emph{de position}).
\item L'eau est lâchée dans une \emph{conduite forcée} : elle dévale la pente et gagne une grande vitesse.
\item Le jet d'eau frappe les pales de la \emph{turbine}, qui se met à tourner.
\item La turbine entraîne l'\emph{alternateur}, qui produit une tension alternative.
\item Un \emph{transformateur} élève la tension pour le transport (nous y reviendrons à la section suivante).
\end{enumerate}

\begin{popfigure}
\begin{center}
\begin{tikzpicture}[scale=1,>=Stealth,font=\small]
% Barrage
\fill[popblue!25] (0,0) -- (0,2.2) -- (1.6,2.2) -- (1.6,0) -- cycle;
\draw[thick,popdark] (1.6,0) -- (1.6,2.6) -- (1.9,2.6) -- (1.9,0);
\node[popblue] at (0.8,1.1) {Retenue};
\node[popblue] at (0.8,0.7) {d'eau};
\node[below,align=center] at (1.75,-0.15) {Barrage};
% Conduite forcée
\draw[thick,popblue,->] (1.9,1.9) -- (3.6,0.55);
\draw[thick,popblue] (1.9,1.6) -- (3.6,0.25);
\node[popblue,above,rotate=-38] at (2.75,1.35) {Conduite forcée};
% Turbine
\draw[fill=poporangeL,thick,poporange] (4.1,0.4) circle (0.45);
\foreach \a in {0,90,180,270}{
  \draw[thick,poporange] (4.1,0.4) -- ++(\a:0.42);
}
\node[below] at (4.1,-0.15) {Turbine};
% Arbre
\draw[thick,popdark] (4.55,0.4) -- (5.15,0.4);
% Alternateur
\draw[fill=popgreenL,thick,popgreen] (5.15,0) rectangle (6.35,0.8);
\node at (5.75,0.4) {Alternateur};
% Transformateur
\draw[fill=poppurpleL,thick,poppurple] (7.15,0) rectangle (8.05,0.8);
\node[poppurple] at (7.6,0.4) {\scriptsize Transfo};
\draw[thick,popdark,->] (6.35,0.4) -- (7.15,0.4);
\draw[thick,popdark,->] (8.05,0.4) -- (8.9,0.4) node[right]{Vers le réseau};
\end{tikzpicture}
\end{center}
\legende{Chaîne d'une centrale hydraulique : l'eau retenue par le barrage dévale la conduite forcée, fait tourner la turbine, qui entraîne l'alternateur. Le transformateur élève ensuite la tension.}
\end{popfigure}

\begin{exemplebox}[Le Cameroun, pays de l'hydroélectricité]
Le Cameroun tire l'\emph{essentiel} de son électricité de l'hydraulique. Les principales centrales sont :
\begin{itemize}
\item \textbf{Songloulou}, sur la Sanaga : la plus grande centrale du pays ;
\item \textbf{Edéa}, sur la Sanaga également, historiquement liée à l'usine ALUCAM ;
\item \textbf{Lagdo}, sur la Bénoué, dans le Nord ;
\item \textbf{Memve'ele}, sur le Ntem, dans le Sud.
\end{itemize}
\end{exemplebox}

\begin{savaistu}[Memve'ele : un géant récent]
Le barrage hydroélectrique de \textbf{Memve'ele}, sur le fleuve Ntem (région du Sud), inauguré récemment, ajoute plus de $200$ MW de puissance au réseau camerounais. À titre de comparaison, cela représente la consommation de centaines de milliers de foyers ! D'autres projets, comme le barrage de Nachtigal sur la Sanaga, renforcent encore cette filière.
\end{savaistu}

\courssub{La centrale thermique : la force de la vapeur}

\textbf{Intuition.} Quand l'eau bout dans une marmite, la vapeur soulève le couvercle. La chaleur peut donc produire un mouvement. Une centrale thermique exploite exactement cette idée, à très grande échelle.

\begin{definition}[Centrale thermique]
Une \cle{centrale thermique} est une centrale où la chaleur dégagée par la combustion d'un \cle{combustible} (gaz, fioul, charbon) transforme de l'eau en \cle{vapeur sous pression} ; cette vapeur fait tourner une turbine couplée à un alternateur.
\end{definition}

\textbf{Fonctionnement :} combustion du gaz ou du fioul $\rightarrow$ production de vapeur d'eau sous pression $\rightarrow$ rotation de la turbine $\rightarrow$ rotation de l'alternateur $\rightarrow$ tension alternative.

\begin{exemplebox}[La centrale de Kribi]
La centrale thermique de \textbf{Kribi}, dans la région du Sud, fonctionne au \emph{gaz naturel} (avec possibilité d'utiliser du fioul). Elle complète les centrales hydrauliques, surtout en saison sèche quand le niveau des retenues d'eau baisse. Le gaz utilisé provient des gisements off-shore du Cameroun.
\end{exemplebox}

\begin{attention}[Hydraulique ou thermique : ne pas confondre ce qui fait tourner la turbine]
Dans les deux cas, une \emph{turbine} entraîne un \emph{alternateur} : cette partie est identique. Ce qui change, c'est la \emph{source d'énergie primaire} :
\begin{itemize}
\item centrale \textbf{hydraulique} : c'est l'\emph{eau en mouvement} qui fait tourner la turbine ;
\item centrale \textbf{thermique} : c'est la \emph{vapeur d'eau sous pression} (produite par une combustion) qui fait tourner la turbine.
\end{itemize}
Au BEPC, on te demandera souvent de compléter une chaîne énergétique : vérifie toujours la première case !
\end{attention}

\courssub{Les autres sources d'énergie électrique}

L'hydraulique et le thermique dominent, mais d'autres filières se développent, notamment au Cameroun :

\begin{itemize}
\item \textbf{Le solaire photovoltaïque} : les panneaux solaires convertissent directement la lumière du Soleil en électricité. Attention : ils fournissent du courant \emph{continu} ! Un appareil appelé \cle{onduleur} le transforme ensuite en courant alternatif $220$ V pour alimenter la maison. Des centrales solaires (Guider, Maroua) alimentent le Grand Nord.
\item \textbf{L'éolien} : le vent fait tourner les pales d'une éolienne, qui entraîne un alternateur. C'est une « centrale hydraulique à air », en quelque sorte.
\item \textbf{La biomasse} : on brûle des déchets végétaux (bois, résidus agricoles) pour produire de la vapeur, comme dans une centrale thermique.
\end{itemize}

\courssub{Bilan énergétique : une centrale convertit, elle ne crée rien}

Résumons le parcours de l'énergie dans une centrale. Quel que soit le type de centrale, le schéma général est le même :

\begin{popfigure}
\begin{center}
\begin{tikzpicture}[scale=1,>=Stealth,font=\small,
  boite/.style={draw,thick,rounded corners=3pt,minimum height=0.9cm,align=center,minimum width=3cm}]
\node[boite,fill=popblueL,draw=popblue, align=center] (a) at (0,0){Énergie primaire\\ (hydraulique ou thermique)};
\node[boite,fill=poporangeL,draw=poporange, align=center] (b) at (4.6,0){Énergie mécanique\\ (turbine en rotation)};
\node[boite,fill=popgreenL,draw=popgreen, align=center] (c) at (9.2,0){Énergie électrique\\ (alternateur)};
\draw[->,very thick,popdark] (a) -- (b) node[midway,above]{turbine};
\draw[->,very thick,popdark] (b) -- (c) node[midway,above]{alternateur};
\end{tikzpicture}
\end{center}
\legende{Chaîne énergétique d'une centrale : l'énergie primaire (eau, chaleur de combustion) devient énergie mécanique grâce à la turbine, puis énergie électrique grâce à l'alternateur.}
\end{popfigure}

\begin{aretenir}[Bilan énergétique d'une centrale]
Dans toute centrale électrique, la chaîne de conversion est :
\[ \text{énergie primaire} \longrightarrow \text{énergie mécanique (turbine)} \longrightarrow \text{énergie électrique (alternateur)} \]
L'énergie primaire est \emph{hydraulique} (chute d'eau) ou \emph{thermique} (combustion produisant de la vapeur).
\end{aretenir}

\begin{attention}[Une centrale ne « fabrique » pas d'énergie à partir de rien]
On entend parfois : « la centrale produit de l'énergie ». Au sens strict, c'est un abus de langage ! D'après le principe de \cle{conservation de l'énergie}, l'énergie ne se crée pas : elle se \emph{convertit}. La centrale \emph{transforme} une énergie déjà existante (énergie de l'eau, énergie chimique du combustible) en énergie électrique. D'ailleurs, une partie de l'énergie est toujours perdue sous forme de chaleur : l'énergie électrique obtenue est \emph{inférieure} à l'énergie primaire consommée.
\end{attention}

\begin{exoresolu}[Compléter une chaîne énergétique]
\textbf{Énoncé.} La centrale de Songloulou utilise l'eau retenue par un barrage sur la Sanaga.
\begin{enumerate}
\item Nommer les deux organes principaux qui assurent les conversions d'énergie dans cette centrale.
\item Recopier et compléter la chaîne : énergie hydraulique $\rightarrow$ \ldots $\rightarrow$ \ldots
\item Pourquoi dit-on que l'alternateur est le « cœur » de la centrale ?
\end{enumerate}
\tcblower
\textbf{Solution détaillée.}
\begin{enumerate}
\item Les deux organes sont la \textbf{turbine} (entraînée par l'eau de la conduite forcée) et l'\textbf{alternateur} (entraîné par la turbine).
\item Chaîne complétée :
\[ \text{énergie hydraulique} \rightarrow \text{énergie mécanique (turbine)} \rightarrow \text{énergie électrique (alternateur)} \]
\item L'alternateur est le cœur de la centrale car c'est lui qui réalise la conversion finale et indispensable : sans lui, l'énergie mécanique de rotation ne deviendrait jamais énergie électrique. Par induction électromagnétique, il transforme la rotation en tension alternative.
\end{enumerate}
\end{exoresolu}

\begin{exercice}[À toi de jouer]
\begin{enumerate}
\item Un élève affirme : « Un aimant posé immobile dans une bobine produit du courant. » Corrige cette affirmation en t'appuyant sur l'expérience d'induction.
\item Complète le tableau suivant :
\begin{center}
\begin{tabularx}{\linewidth}{|l|X|X|}
\hline
\rowcolor{popblueL} \textbf{Type de centrale} & \textbf{Énergie primaire} & \textbf{Exemple camerounais} \\
\hline
Hydraulique & & \\
\hline
Thermique & & \\
\hline
\end{tabularx}
\end{center}
\item La centrale solaire de Guider produit du courant continu. Quel appareil faut-il ajouter pour alimenter les habitations en $220$ V alternatif ?
\end{enumerate}
\end{exercice}

\begin{corrige}[Exercice : À toi de jouer]
\begin{enumerate}
\item L'affirmation est fausse : l'expérience montre que le galvanomètre ne dévie \emph{que pendant le mouvement} de l'aimant (ou de la bobine). Aimant immobile : aucune tension, aucun courant. C'est le mouvement relatif qui produit la tension alternative.
\item Hydraulique : énergie de l'eau (chute ou retenue) ; exemple : Songloulou (ou Edéa, Lagdo, Memve'ele). Thermique : chaleur de combustion d'un combustible (gaz, fioul) ; exemple : centrale de Kribi.
\item Il faut un \textbf{onduleur}, qui convertit le courant continu des panneaux solaires en courant alternatif $220$ V.
\end{enumerate}
\end{corrige}

Nous savons désormais \emph{produire} une tension alternative grâce à l'alternateur d'une centrale. Mais Songloulou est à des centaines de kilomètres de ta maison : comment cette énergie voyage-t-elle sans trop de pertes ? C'est le rôle des lignes haute tension et des transformateurs, que nous étudierons dans la section suivante.

\coursec{Transport et distribution : lignes haute tension et transformateurs}

Dans la section précédente, nous avons vu comment les centrales électriques — comme la centrale hydroélectrique de Songloulou sur le fleuve Sanaga ou le barrage de Lagdo sur la Bénoué — produisent l'énergie électrique. Mais voici un problème : Songloulou se trouve à des centaines de kilomètres de Yaoundé, de Douala ou de Ngaoundéré. Comment l'électricité produite là-bas arrive-t-elle jusqu'à la prise de courant de ta maison ? Et pourquoi voit-on, le long des routes, d'immenses pylônes métalliques portant des câbles, puis de petits transformateurs gris accrochés aux poteaux des quartiers ? C'est toute la question du \cle{transport} et de la \cle{distribution} de l'énergie électrique, que nous allons résoudre pas à pas.

\courssub{Le problème du transport : les pertes par effet Joule}

\textbf{Une situation concrète.}\par
Imagine que tu transportes de l'eau d'un puits jusqu'à ton quartier par un long tuyau : plus le tuyau est long et étroit, plus tu perds d'eau par frottement. Pour l'électricité, c'est pareil : les câbles qui relient la centrale aux villes sont longs de plusieurs centaines de kilomètres, et même un bon conducteur comme le cuivre ou l'aluminium possède une résistance $R$ qui n'est pas nulle.

Or tu as appris qu'un conducteur traversé par un courant d'intensité $I$ s'échauffe : c'est \cle{l'effet Joule}. La puissance perdue en chaleur dans un câble de résistance $R$ est :
\[ P_{\text{perdue}} = R \cdot I^{2} \]

Cette énergie perdue ne profite à personne : elle chauffe inutilement les câbles, et c'est autant d'électricité payée par la société ENEO qui n'arrive jamais chez l'abonné. Sur des centaines de kilomètres, ces pertes peuvent devenir énormes si l'on n'y prend pas garde.

\begin{experience}[Observation : l'échauffement d'un fil]
\textbf{Matériel :} un générateur de tension continue réglable, un fil conducteur long et fin (fil de fer), un ampèremètre, un thermomètre (ou simplement le toucher prudent du fil).
\begin{enumerate}
\item Faire passer un courant faible (\SI{0,5}{A}) dans le fil. Observer la température : le fil reste presque froid.
\item Augmenter l'intensité à \SI{2}{A}, puis \SI{4}{A}. Observer à nouveau.
\end{enumerate}
\textbf{Observations :} à \SI{2}{A}, le fil devient tiède ; à \SI{4}{A}, il devient nettement chaud.
\textbf{Interprétation :} quand l'intensité double, la puissance dissipée $R \cdot I^{2}$ est multipliée par \textbf{quatre} (car $2^{2} = 4$). Quand elle est multipliée par 4, les pertes sont multipliées par 16.
\textbf{Conclusion :} les pertes par effet Joule augmentent très vite avec l'intensité, car elles dépendent du \cle{carré} de l'intensité. Pour transporter l'électricité loin, il faut donc absolument \textbf{diminuer l'intensité} dans les câbles.
\end{experience}

\courssub{La solution : transporter sous haute tension}

Comment diminuer l'intensité tout en transportant la même puissance ? Rappelons la relation entre puissance, tension et intensité en régime alternatif (pour les appareils usuels) :
\[ P = U \cdot I \]

La centrale fournit une puissance $P$ déterminée par les besoins des consommateurs : cette puissance est \textbf{imposée}. On ne peut donc jouer que sur $U$ et $I$, qui sont liés : si l'on augmente la tension $U$, l'intensité $I$ diminue automatiquement dans la même proportion.

\begin{methode}[Justifier le transport en haute tension]
Pour expliquer pourquoi on transporte l'électricité sous haute tension, raisonne en quatre étapes :
\begin{enumerate}
\item La puissance à transporter est fixée : $P = U \cdot I$ = constante.
\item Si l'on augmente fortement la tension $U$, alors l'intensité $I = \dfrac{P}{U}$ devient faible.
\item Les pertes par effet Joule valent $P_{\text{perdue}} = R \cdot I^{2}$ : si $I$ devient faible, les pertes deviennent très faibles (car elles dépendent de $I^{2}$).
\item Conclusion : \textbf{plus la tension de transport est élevée, moins on perd d'énergie en chaleur dans les câbles.} C'est pourquoi on utilise des lignes à \cle{haute tension}.
\end{enumerate}
\end{methode}

\begin{exemplebox}[Un calcul qui parle de lui-même]
Une centrale doit transporter une puissance $P = \SI{220}{kW}$.
\begin{itemize}
\item \textbf{Cas 1 : transport sous \SI{220}{V}.} L'intensité nécessaire serait :
\[ I = \frac{P}{U} = \frac{220\,000}{220} = \SI{1000}{A} \]
Mille ampères ! Les câbles chaufferaient énormément et une grande partie de l'énergie serait perdue.
\item \textbf{Cas 2 : transport sous \SI{220}{kV}.} L'intensité devient :
\[ I = \frac{P}{U} = \frac{220\,000}{220\,000} = \SI{1}{A} \]
Un seul ampère ! Les pertes par effet Joule sont divisées par $1000^{2} = 1\,000\,000$ : elles deviennent presque nulles.
\end{itemize}
En multipliant la tension par 1000, on a divisé les pertes par un million. Voilà pourquoi ENEO transporte l'électricité sous \SI{90}{kV} ou \SI{225}{kV} sur ses lignes haute tension.
\end{exemplebox}

\begin{aretenir}[L'idée centrale de la section]
À puissance transportée égale, \textbf{augmenter la tension diminue l'intensité}, ce qui réduit fortement les pertes par effet Joule ($P_{\text{perdue}} = R \cdot I^{2}$) et l'échauffement des câbles. Le transport de l'énergie électrique sur de longues distances se fait donc sous \cle{haute tension}.
\end{aretenir}

\courssub{Le transformateur : l'appareil qui change la tension}

Reste une question : comment passer de la tension produite par la centrale (quelques milliers de volts) à \SI{225}{kV} pour le transport, puis redescendre à \SI{220}{V} pour les habitations ? Il faut un appareil capable de modifier la valeur d'une tension : c'est le \cle{transformateur}.

\begin{definition}[Transformateur]
Un \cle{transformateur} est un appareil qui \textbf{modifie la valeur d'une tension alternative} :
\begin{itemize}
\item s'il \textbf{augmente} la tension, on dit qu'il est \cle{élévateur} de tension ;
\item s'il \textbf{diminue} la tension, on dit qu'il est \cle{abaisseur} de tension.
\end{itemize}
\end{definition}

\textbf{Structure d'un transformateur.}\par
Un transformateur est constitué de :
\begin{itemize}
\item deux \cle{bobines} (enroulements de fil de cuivre) : la bobine d'entrée, appelée \cle{primaire}, reliée à la tension $U_{1}$, et la bobine de sortie, appelée \cle{secondaire}, qui fournit la tension $U_{2}$ ;
\item un \cle{noyau de fer} doux, fermé, sur lequel les deux bobines sont enroulées.
\end{itemize}

\begin{popfigure}
\begin{center}
\begin{tikzpicture}[scale=1, every node/.style={font=\small}]
% Noyau de fer (cadre rectangulaire épais)
\draw[line width=6pt, popmuted] (0,0) rectangle (4,3);
\draw[line width=6pt, popmuted!60] (0.35,0.35) rectangle (3.65,2.65);
% Bobine primaire (à gauche)
\foreach \y in {0.5,0.8,1.1,1.4,1.7,2.0,2.3,2.5}{
  \draw[popblue, thick] (0.05,\y) arc[start angle=-90, end angle=90, x radius=0.12, y radius=0.15];
}
\node[popblue, font=\small\bfseries] at (-0.6,1.5) {Primaire};
\node[popblue] at (-0.6,1.1) {$U_{1}$};
% Bobine secondaire (à droite)
\foreach \y in {0.7,1.0,1.3,1.6,1.9,2.2}{
  \draw[poporange, thick] (3.95,\y) arc[start angle=270, end angle=90, x radius=0.12, y radius=0.15];
}
\node[poporange, font=\small\bfseries] at (4.7,1.5) {Secondaire};
\node[poporange] at (4.7,1.1) {$U_{2}$};
% Noyau label
\node[popdark] at (2,-0.35) {Noyau de fer};
% Flèche champ magnétique
\draw[-{Stealth}, popgreen, thick] (1.2,1.5) -- (2.8,1.5);
\node[popgreen] at (2,1.8) {induction};
\end{tikzpicture}
\end{center}
\legende{Schéma simplifié d'un transformateur : deux bobines (primaire en bleu, secondaire en orange) enroulées sur un même noyau de fer. La tension $U_{1}$ appliquée au primaire crée, par induction, une tension $U_{2}$ au secondaire.}
\end{popfigure}

\textbf{Principe de fonctionnement (savoir admis).}\par
Quand la bobine primaire est parcourue par un courant \textbf{alternatif}, elle crée dans le noyau de fer un champ magnétique variable. Ce champ magnétique variable induit à son tour une tension alternative dans la bobine secondaire : c'est le phénomène d'\cle{induction}. Selon le nombre de spires de chaque bobine, la tension de sortie $U_{2}$ est plus grande ou plus petite que la tension d'entrée $U_{1}$.

\begin{propriete}[Le transformateur exige le courant alternatif — admis]
Un transformateur \textbf{ne fonctionne qu'avec une tension alternative}. Alimenté en courant continu, il ne transforme rien (et peut même être endommagé, car le primaire se comporte alors comme un simple fil).
\medskip

C'est l'une des \textbf{raisons majeures du choix du courant alternatif} pour le réseau électrique : seul l'alternatif permet d'élever et d'abaisser facilement la tension grâce au transformateur, donc de transporter l'énergie sous haute tension avec peu de pertes.
\end{propriete}

\begin{attention}[Erreur fréquente]
Beaucoup d'élèves croient que « le courant qui arrive à la maison vient directement de la centrale, sans transformation ». C'est \textbf{faux} ! Entre la centrale et ta prise de courant, l'électricité subit \textbf{plusieurs transformations de tension} : élévation à la sortie de la centrale, puis abaissements successifs dans les postes de transformation et sur le transformateur de ton quartier. Au BEPC, cite toujours au moins ces deux étapes : transformateur élévateur puis transformateur abaisseur.
\end{attention}

\courssub{Les étapes du transport : de la centrale au quartier}

Suivons le voyage de l'électricité, de Songloulou jusqu'à une maison de Yaoundé.

\begin{enumerate}
\item \textbf{À la sortie de la centrale}, un \cle{transformateur élévateur} porte la tension à des valeurs très élevées : \SI{90}{kV} ou \SI{225}{kV} au Cameroun. C'est la tension des \cle{lignes haute tension}.
\item \textbf{Le transport} se fait par les lignes haute tension, ces gros câbles suspendus aux grands pylônes métalliques que l'on aperçoit le long des routes. Sous haute tension, l'intensité est faible : les pertes par effet Joule sont réduites.
\item \textbf{À l'approche des villes}, des \cle{postes de transformation} abaissent la tension (par exemple de \SI{90}{kV} vers \SI{15}{kV} ou \SI{33}{kV}).
\item \textbf{Dans les quartiers}, des \cle{transformateurs abaisseurs} — ces gros boîtiers gris fixés sur les poteaux — ramènent la tension à environ \SI{220}{V}, la tension utilisée par les habitations.
\item L'électricité arrive enfin au \cle{compteur} de l'abonné, puis au tableau de distribution de la maison.
\end{enumerate}

\begin{popfigure}
\begin{center}
\begin{tikzpicture}[scale=0.92, every node/.style={font=\scriptsize}]
% Centrale
\draw[fill=popblueL, draw=popblue, thick] (0,0) rectangle (1.6,1.2);
\draw[fill=popblueL, draw=popblue, thick] (0.2,1.2) -- (0.2,1.7) -- (0.5,1.7) -- (0.5,1.2);
\node[popdark] at (0.8,0.6) {Centrale};
\node[popdark] at (0.8,-0.3) {(Songloulou)};
% Transfo élévateur
\draw[fill=popgreenL, draw=popgreen, thick] (2.4,0.2) rectangle (3.4,1.0);
\node[popdark] at (2.9,0.6) {T$_{\text{élév.}}$};
\node[popdark] at (2.9,-0.3) {\SI{225}{kV}};
% Pylône 1
\begin{scope}[shift={(5.2,0)}]
\draw[popdark, thick] (0,0) -- (0.35,1.8) -- (0.7,0);
\draw[popdark, thick] (0.1,0.9) -- (0.6,0.9);
\draw[popdark, thick] (-0.15,1.5) -- (0.85,1.5);
\draw[popdark, thick] (-0.15,1.2) -- (0.85,1.2);
\node[popdark] at (0.35,-0.3) {Pylône HT};
\end{scope}
% Pylône 2
\begin{scope}[shift={(7.2,0)}]
\draw[popdark, thick] (0,0) -- (0.35,1.8) -- (0.7,0);
\draw[popdark, thick] (0.1,0.9) -- (0.6,0.9);
\draw[popdark, thick] (-0.15,1.5) -- (0.85,1.5);
\draw[popdark, thick] (-0.15,1.2) -- (0.85,1.2);
\end{scope}
% Lignes entre pylônes
\draw[poporange, thick] (3.4,1.5) -- (5.05,1.5);
\draw[poporange, thick] (6.05,1.5) -- (7.05,1.5);
\draw[poporange, thick] (7.9,1.5) -- (9.1,1.5);
\draw[poporange, thick] (3.4,1.2) -- (5.05,1.2);
\draw[poporange, thick] (6.05,1.2) -- (7.05,1.2);
\draw[poporange, thick] (7.9,1.2) -- (9.1,1.2);
% Transfo abaisseur de quartier
\draw[fill=poporangeL, draw=poporange, thick] (9.1,0.2) rectangle (10.1,1.0);
\node[popdark] at (9.6,0.6) {T$_{\text{abais.}}$};
\node[popdark] at (9.6,-0.3) {\SI{220}{V}};
% Compteur
\draw[fill=poppurpleL, draw=poppurple, thick] (10.9,0.25) rectangle (11.6,0.95);
\node[popdark] at (11.25,0.6) {Cpt};
\node[popdark] at (11.25,-0.3) {Compteur};
% Maison
\draw[fill=popcream, draw=popgold, thick] (12.4,0) rectangle (14,1.0);
\draw[fill=popgold!40, draw=popgold, thick] (12.3,1.0) -- (13.2,1.6) -- (14.1,1.0);
\node[popdark] at (13.2,0.45) {Maison};
% Liaisons
\draw[-{Stealth}, thick, popink] (1.6,0.6) -- (2.4,0.6);
\draw[-{Stealth}, thick, popink] (10.1,0.6) -- (10.9,0.6);
\draw[-{Stealth}, thick, popink] (11.6,0.6) -- (12.4,0.6);
% Accolade ligne HT
\node[poporange, font=\scriptsize\itshape] at (6.5,2.15) {Ligne haute tension : pertes faibles};
\end{tikzpicture}
\end{center}
\legende{La chaîne complète de transport et de distribution : la centrale alimente un transformateur élévateur (T$_{\text{élév.}}$) qui porte la tension à \SI{225}{kV} ; les lignes haute tension transportent l'énergie sur de longues distances ; un transformateur abaisseur (T$_{\text{abais.}}$) de quartier ramène la tension à \SI{220}{V} avant le compteur et l'habitation.}
\end{popfigure}

\begin{definition}[Réseau de transport et de distribution]
Le \cle{réseau ENEO} est l'ensemble des lignes (haute, moyenne et basse tension), des postes de transformation et des transformateurs qui acheminent l'électricité de la centrale jusqu'au compteur de l'abonné.
\end{definition}

\begin{savaistu}[Des câbles nus dans le ciel]
Les câbles des lignes haute tension ne sont \textbf{pas recouverts d'une gaine isolante} comme les fils de ta maison : c'est \textbf{l'air} qui joue le rôle d'isolant entre eux et par rapport au sol. Une gaine plastique serait d'ailleurs inutile et trop lourde à ces tensions. Conséquence : il ne faut \textbf{jamais} s'approcher d'une ligne haute tension, ni faire voler un cerf-volant à proximité, ni toucher un câble tombé au sol après un accident ou une tempête. À \SI{225}{kV}, l'arc électrique peut jaillir sans même qu'on touche le câble !
\end{savaistu}

\begin{exoresolu}[Calcul d'intensité dans une ligne haute tension]
\textbf{Énoncé.} La centrale de Songloulou transporte une puissance $P = \SI{45}{MW}$ sur une ligne haute tension de $U = \SI{225}{kV}$.
\begin{enumerate}
\item Calculer l'intensité $I$ du courant dans la ligne.
\item Calculer cette même intensité si le transport se faisait sous \SI{220}{V}. Conclure.
\end{enumerate}
\tcblower
\textbf{Solution détaillée.}
\begin{enumerate}
\item On utilise $P = U \cdot I$, donc :
\begin{align*}
I &= \frac{P}{U} = \frac{45\,000\,000}{225\,000} = \SI{200}{A}
\end{align*}
L'intensité dans la ligne haute tension est de \SI{200}{A}.
\item Sous \SI{220}{V} :
\begin{align*}
I' &= \frac{P}{U'} = \frac{45\,000\,000}{220} \approx \SI{204545}{A}
\end{align*}
Il faudrait plus de \SI{200000}{A}, ce qui est matériellement impossible : les câbles fondraient.
\textbf{Conclusion :} grâce à la haute tension, l'intensité est environ mille fois plus faible, et les pertes par effet Joule (en $R \cdot I^{2}$) sont environ un million de fois plus faibles. Le transformateur élévateur est donc indispensable au transport de l'énergie électrique.
\end{enumerate}
\end{exoresolu}

\begin{exercice}[À toi de jouer]
Une ligne transporte une puissance de \SI{11}{MW} sous une tension de \SI{90}{kV}.
\begin{enumerate}
\item Calculer l'intensité du courant dans la ligne (arrondir à l'unité).
\item La résistance totale de la ligne est $R = \SI{10}{\Omega}$. Calculer la puissance perdue par effet Joule.
\item Que deviendrait cette perte si, à puissance égale, la tension de transport était divisée par 3 ?
\end{enumerate}
\end{exercice}

\begin{corrige}[Exercice : À toi de jouer]
\begin{enumerate}
\item $I = \dfrac{P}{U} = \dfrac{11\,000\,000}{90\,000} \approx \SI{122}{A}$.
\item $P_{\text{perdue}} = R \cdot I^{2} = 10 \times 122^{2} \approx \SI{148840}{W} \approx \SI{149}{kW}$.
\item Si $U$ est divisée par 3, alors $I$ est multipliée par 3 ($I' \approx \SI{367}{A}$), et la perte est multipliée par $3^{2} = 9$ : $P'_{\text{perdue}} \approx \SI{1,34}{MW}$. On retrouve que baisser la tension de transport augmente fortement les pertes.
\end{enumerate}
\end{corrige}

\begin{aretenir}[Bilan de la section]
\begin{itemize}
\item Le transport de l'électricité sur de longues distances provoque des pertes par \cle{effet Joule} : $P_{\text{perdue}} = R \cdot I^{2}$.
\item À puissance égale ($P = U \cdot I$), on \textbf{élève la tension} pour \textbf{réduire l'intensité} et donc les pertes : d'où les lignes haute tension (\SI{90}{kV}, \SI{225}{kV} au Cameroun).
\item Le \cle{transformateur} élève ou abaisse une tension alternative ; il fonctionne par \cle{induction} entre deux bobines sur un noyau de fer, et \textbf{uniquement en courant alternatif}.
\item La chaîne complète : centrale $\rightarrow$ transformateur élévateur $\rightarrow$ ligne haute tension $\rightarrow$ transformateurs abaisseurs $\rightarrow$ compteur $\rightarrow$ habitation (\SI{220}{V}).
\end{itemize}
\end{aretenir}

\coursec{De la centrale à l'habitation : la chaîne ENEO et le compteur}

Dans la section précédente, nous avons vu comment l'énergie électrique est transportée par les lignes à haute tension et comment les transformateurs modifient la tension. Mais comment cette énergie parvient-elle exactement jusqu'à la prise de courant de ta maison, à Yaoundé, à Douala ou dans ton village ? Et comment ENEO sait-elle combien tu dois payer à la fin du mois ? C'est tout l'objet de cette section : suivre le courant électrique, étape par étape, de la centrale jusqu'à ton ampoule, et comprendre le rôle du compteur.

\courssub{La chaîne complète de distribution}

\begin{experience}[Observer la chaîne autour de soi]
\textbf{Observation :} En te promenant dans ton quartier, tu peux remarquer sur les poteaux électriques des boîtiers cylindriques gris : ce sont des transformateurs de quartier. À l'entrée de ta maison ou de ta concession, tu trouves un compteur. Enfin, à l'intérieur, le courant alimente les lampes, le téléviseur, le réfrigérateur.

\textbf{Question :} Par quel chemin le courant produit à la centrale de Song Loulou ou d'Édéa arrive-t-il jusqu'à ta lampe ?

\textbf{Hypothèse :} Le courant passe par plusieurs étapes où sa tension est modifiée.

\textbf{Vérification :} En interrogeant un agent ENEO ou en lisant les documents du réseau, on reconstitue la chaîne complète.

\textbf{Conclusion :} Le courant traverse une succession d'étapes bien précises, avec des tensions différentes à chaque étape.
\end{experience}

\begin{definition}[La chaîne de distribution]
La \cle{chaîne de distribution} de l'énergie électrique est l'ensemble des étapes que parcourt le courant depuis la centrale de production jusqu'à l'installation domestique de l'abonné :
\begin{enumerate}
\item la \cle{centrale} produit le courant alternatif sous une tension moyenne (par exemple \SI{15}{kV}) ;
\item un \cle{transformateur élévateur}, placé à la sortie de la centrale, élève la tension (par exemple à \SI{225}{kV}) pour le transport ;
\item la \cle{ligne haute tension} transporte l'énergie sur de longues distances avec peu de pertes ;
\item un \cle{poste de transformation} abaisse la tension (retour à \SI{15}{kV} environ) à l'arrivée en ville ;
\item le \cle{transformateur de quartier}, fixé sur un poteau, abaisse encore la tension jusqu'à \SI{220}{V} ;
\item la \cle{ligne basse tension} dessert les maisons du quartier ;
\item le \cle{compteur} mesure l'énergie consommée par l'abonné ;
\item l'\cle{installation domestique} distribue le courant aux appareils de la maison.
\end{enumerate}
\end{definition}

\begin{popfigure}
\begin{center}
\begin{tikzpicture}[font=\small, >=Stealth,
  etape/.style={draw=popblue, fill=popblueL, rounded corners=2pt, align=center, minimum height=0.9cm, text width=2.6cm},
  tensionlbl/.style={font=\footnotesize\itshape, text=popdark}]
\node[etape, align=center] (c) at (0,0){Centrale\\ \SI{15}{kV}};
\node[etape, right=0.55cm of c] (te) {Transformateur élévateur};
\node[etape, right=0.55cm of te, align=center] (lht){Ligne haute tension\\ \SI{225}{kV}};
\node[etape, right=0.55cm of lht] (pt) {Poste de transformation};
\node[etape, below=0.8cm of pt] (tq) {Transformateur de quartier (poteau)};
\node[etape, left=0.55cm of tq, align=center] (lbt){Ligne basse tension\\ \SI{220}{V}};
\node[etape, left=0.55cm of lbt] (compt) {Compteur};
\node[etape, left=0.55cm of compt] (maison) {Installation domestique};
\draw[->, thick, poporange] (c) -- (te);
\draw[->, thick, poporange] (te) -- (lht);
\draw[->, thick, poporange] (lht) -- (pt);
\draw[->, thick, poporange] (pt) -- (tq);
\draw[->, thick, poporange] (tq) -- (lbt);
\draw[->, thick, poporange] (lbt) -- (compt);
\draw[->, thick, poporange] (compt) -- (maison);
\node[tensionlbl, below=0.05cm of te] {\SI{15}{kV} $\to$ \SI{225}{kV}};
\node[tensionlbl, above=0.05cm of pt] {\SI{225}{kV} $\to$ \SI{15}{kV}};
\node[tensionlbl, below=0.05cm of tq] {\SI{15}{kV} $\to$ \SI{220}{V}};
\end{tikzpicture}
\end{center}
\legende{La chaîne complète de distribution de l'énergie électrique, de la centrale à l'habitation, avec les valeurs de tension à chaque étape. Les transformateurs élèvent la tension pour le transport puis l'abaissent progressivement jusqu'à \SI{220}{V} pour l'usage domestique.}
\end{popfigure}

\begin{aretenir}
Le courant qui arrive dans ta maison a parcouru un long chemin : \textbf{centrale} $\to$ \textbf{transformateur élévateur} $\to$ \textbf{ligne haute tension} $\to$ \textbf{poste de transformation} $\to$ \textbf{transformateur de quartier} $\to$ \textbf{ligne basse tension} $\to$ \textbf{compteur} $\to$ \textbf{installation domestique}. La tension passe successivement de \SI{15}{kV} à \SI{225}{kV}, redescend à \SI{15}{kV}, puis à \SI{220}{V}.
\end{aretenir}

\courssub{Le compteur électrique : mesurer l'énergie consommée}

Lorsque tu paies ta facture ENEO, tu ne paies pas la puissance de tes appareils : tu paies l'\cle{énergie} qu'ils ont consommée. C'est le compteur qui mesure cette énergie.

\begin{definition}[Rôle du compteur]
Le \cle{compteur électrique} est un appareil placé à l'entrée de l'installation domestique. Il \textbf{mesure l'énergie électrique consommée} par l'abonné. Cette énergie est exprimée en \cle{kilowattheures} (symbole : kWh). Le cadran du compteur affiche un nombre appelé \cle{index}, qui augmente au fur et à mesure de la consommation.
\end{definition}

\begin{definition}[Le kilowattheure]
Le \cle{kilowattheure} (kWh) est l'unité d'énergie électrique utilisée pour la facturation. L'énergie consommée par un appareil se calcule par la relation :
\[ E = P \times t \]
où $E$ est l'énergie en kilowattheures (kWh), $P$ la puissance en kilowatts (kW) et $t$ la durée d'utilisation en heures (h).

Un kilowattheure est l'énergie consommée par un appareil de puissance \SI{1}{kW} fonctionnant pendant \SI{1}{h}.
\end{definition}

\begin{exemplebox}[Une ampoule de 100 W]
Une ampoule de puissance $P = \SI{100}{W} = \SI{0,1}{kW}$ reste allumée pendant $t = \SI{10}{h}$. L'énergie consommée est :
\[ E = P \times t = 0,1 \times 10 = \SI{1}{kWh}. \]
Cette ampoule a donc consommé exactement un kilowattheure.
\end{exemplebox}

\begin{methode}[Calculer une consommation d'énergie]
Pour calculer l'énergie consommée par un appareil :
\begin{enumerate}
\item Convertis la puissance en kilowatts : $P\ (\text{kW}) = \dfrac{P\ (\text{W})}{1000}$.
\item Exprime la durée d'utilisation en heures.
\item Applique la formule : $E\ (\text{kWh}) = P\ (\text{kW}) \times t\ (\text{h})$.
\end{enumerate}
Pour connaître la consommation d'un mois entier à partir du compteur :
\[ \text{Consommation du mois} = \text{nouvel index} - \text{ancien index}. \]
\end{methode}

\begin{attention}[Ne pas confondre kW et kWh]
C'est l'erreur la plus fréquente au BEPC ! Le \textbf{kilowatt (kW)} est une unité de \textbf{puissance} : il caractérise l'appareil (par exemple un fer à repasser de \SI{1}{kW}). Le \textbf{kilowattheure (kWh)} est une unité d'\textbf{énergie} : il dépend de la puissance \emph{et} de la durée d'utilisation. Dire « mon fer consomme 1 kW par heure » est incorrect ; on dit « mon fer consomme \SI{1}{kWh} en une heure de fonctionnement ».
\end{attention}

\begin{exoresolu}[Le fer à repasser de maman]
\textbf{Énoncé :} Un fer à repasser de puissance $P = \SI{1}{kW}$ est utilisé $2$ heures par jour pendant $30$ jours. Le kilowattheure est facturé $50$ FCFA. Calcule l'énergie consommée dans le mois, puis le coût correspondant.
\tcblower
\textbf{Solution détaillée :}

Étape 1 : durée totale d'utilisation dans le mois.
\[ t = 2 \times 30 = \SI{60}{h}. \]

Étape 2 : énergie consommée.
\[ E = P \times t = 1 \times 60 = \SI{60}{kWh}. \]

Étape 3 : coût de la consommation.
\[ \text{Coût} = 60 \times 50 = 3000\ \text{FCFA}. \]

\textbf{Conclusion :} Le fer à repasser consomme \SI{60}{kWh} dans le mois, ce qui coûte $3000$ FCFA sur la facture ENEO.
\end{exoresolu}

\begin{popfigure}
\begin{center}
\begin{tikzpicture}[font=\small]
\draw[thick, popdark, fill=popcream, rounded corners=4pt] (0,0) rectangle (7,2.6);
\draw[fill=popblueL, draw=popblue] (0.3,0.5) rectangle (6.7,2.2);
\node[font=\footnotesize\bfseries, text=popdark] at (3.5,2.42) {COMPTEUR ENEO -- INDEX (kWh)};
\foreach \x/\chiffre in {0.9/0, 1.9/0, 2.9/8, 3.9/4, 4.9/7, 5.9/3} {
  \draw[fill=white, draw=popdark] (\x-0.45,0.8) rectangle (\x+0.45,1.9);
  \node[font=\Large\bfseries] at (\x,1.35) {\chiffre};
}
\node[font=\footnotesize\itshape, text=popmuted] at (3.5,0.25) {Lecture : index = 008473 kWh};
\end{tikzpicture}
\end{center}
\legende{Reproduction simplifiée du cadran d'un compteur électrique. Les chiffres affichés forment l'index, exprimé en kilowattheures. La consommation d'une période s'obtient en soustrayant l'ancien index du nouvel index.}
\end{popfigure}

\courssub{La facture ENEO : lecture et calcul}

Chaque mois, un agent relève l'index de ton compteur, ou tu le communiques toi-même. La facture ENEO repose sur un principe simple : le prix payé est \textbf{proportionnel au nombre de kilowattheures consommés}.

\begin{propriete}[Calcul de la consommation facturée]
Sur une facture ENEO, on lit deux nombres :
\begin{itemize}
\item l'\cle{ancien index} : la valeur affichée par le compteur lors de la relève précédente ;
\item le \cle{nouvel index} : la valeur affichée lors de la relève actuelle.
\end{itemize}
La consommation de la période est la différence des deux index :
\[ \text{Consommation (kWh)} = \text{nouvel index} - \text{ancien index}. \]
Le montant à payer s'obtient en multipliant cette consommation par le prix du kilowattheure. En pratique, ENEO applique un \cle{tarif par tranche} : les premiers kilowattheures consommés sont facturés moins cher que les suivants, afin que les petits consommateurs paient moins.
\end{propriete}

\begin{exemplebox}[Lecture d'une facture simplifiée]
Voici une facture ENEO simplifiée de la famille Ngo Bell :

\begin{center}
\begin{tabularx}{\linewidth}{|l|X|}\hline
\rowcolor{popblueL} \textbf{Élément} & \textbf{Valeur} \\ \hline
Ancien index & \SI{8473}{kWh} \\ \hline
Nouvel index & \SI{8623}{kWh} \\ \hline
Consommation du mois & $8623 - 8473 = \SI{150}{kWh}$ \\ \hline
Prix moyen du kWh & $50$ FCFA \\ \hline
Montant de la consommation & $150 \times 50 = 7500$ FCFA \\ \hline
\end{tabularx}
\end{center}

La famille Ngo Bell paie donc $7500$ FCFA pour sa consommation du mois (hors taxes et redevances).
\end{exemplebox}

\begin{exoresolu}[Vérifier sa facture]
\textbf{Énoncé :} Le compteur de la famille Etame affichait \SI{12040}{kWh} fin mars et \SI{12160}{kWh} fin avril. Le kilowattheure coûte $50$ FCFA. Calcule la consommation du mois d'avril et le montant à payer.
\tcblower
\textbf{Solution détaillée :}

Étape 1 : consommation du mois.
\[ \text{Consommation} = \text{nouvel index} - \text{ancien index} = 12160 - 12040 = \SI{120}{kWh}. \]

Étape 2 : montant à payer.
\[ \text{Montant} = 120 \times 50 = 6000\ \text{FCFA}. \]

\textbf{Conclusion :} La famille Etame a consommé \SI{120}{kWh} en avril et doit payer $6000$ FCFA.
\end{exoresolu}

\begin{savaistu}[Réduire sa facture ENEO]
De nombreux appareils (téléviseur, décodeur, chargeur branché) continuent de consommer du courant lorsqu'ils sont en \textbf{veille} : c'est de l'énergie payée pour rien ! Éteindre complètement les appareils, débrancher les chargeurs et remplacer les ampoules à incandescence par des \textbf{ampoules LED} (qui consomment environ $8$ fois moins pour le même éclairage) permet de réduire sensiblement la facture ENEO. Par exemple, remplacer une ampoule de \SI{100}{W} par une LED de \SI{12}{W} allumée $5$ h par jour économise environ \SI{13}{kWh} par mois, soit plus de $600$ FCFA !
\end{savaistu}

\courssub{Complément : pourquoi y a-t-il des délestages ?}

\emph{Ce complément dépasse le strict programme, mais il t'aide à comprendre la vie quotidienne.}

Tu as certainement remarqué que le courant est parfois coupé dans ton quartier alors que le quartier voisin reste éclairé : c'est le \cle{délestage}. Le délestage est une coupure volontaire et tournante organisée par ENEO lorsque la \textbf{production} des centrales est \textbf{insuffisante par rapport à la demande} sur le réseau. Si tous les abonnés demandaient en même temps plus d'énergie que les centrales ne peuvent en fournir, tout le réseau risquerait une panne générale. Pour l'éviter, on coupe provisoirement certains quartiers, à tour de rôle. Le délestage n'est donc pas une panne de ton installation : c'est un déséquilibre entre la production et la consommation à l'échelle du pays.

\begin{exercice}[Pour t'entraîner]
Un téléviseur de puissance $P = \SI{150}{W}$ fonctionne en moyenne $4$ heures par jour pendant un mois de $30$ jours.
\begin{enumerate}
\item Convertis la puissance en kilowatts.
\item Calcule l'énergie consommée dans le mois, en kWh.
\item Calcule le coût de cette consommation si le kWh vaut $50$ FCFA.
\end{enumerate}
\end{exercice}

\begin{corrige}[Exercice : le téléviseur]
\begin{enumerate}
\item $P = \dfrac{150}{1000} = \SI{0,15}{kW}$.
\item Durée totale : $t = 4 \times 30 = \SI{120}{h}$. Énergie : $E = P \times t = 0,15 \times 120 = \SI{18}{kWh}$.
\item Coût : $18 \times 50 = 900$ FCFA.
\end{enumerate}
Le téléviseur consomme \SI{18}{kWh} dans le mois, pour un coût de $900$ FCFA.
\end{corrige}

\begin{aretenir}[L'essentiel de la section]
\begin{itemize}
\item Le courant parcourt une chaîne : centrale $\to$ transformateur élévateur $\to$ ligne haute tension $\to$ poste de transformation $\to$ transformateur de quartier $\to$ ligne basse tension $\to$ compteur $\to$ maison.
\item Le compteur mesure l'énergie électrique consommée, en kilowattheures (kWh).
\item L'énergie se calcule par $E = P \times t$, avec $P$ en kW et $t$ en heures.
\item La consommation du mois est la différence entre le nouvel index et l'ancien index ; la facture est proportionnelle au nombre de kWh consommés.
\item Le kW mesure une puissance, le kWh mesure une énergie : ne les confonds jamais !
\end{itemize}
\end{aretenir}

\coursec{Sécurité avec le courant électrique domestique}

Dans la section précédente, nous avons suivi le courant depuis la centrale jusqu'au compteur de ta maison : après le transformateur de quartier, la tension de \SI{220}{V} arrive chez toi, prête à alimenter la télévision, le fer à repasser ou le chargeur de téléphone. Cette énergie est formidablement utile... mais elle peut aussi tuer. Chaque année, des accidents domestiques liés à l'électricité font des victimes au Cameroun, souvent par ignorance de quelques règles simples. Cette dernière section est peut-être la plus importante de toute la leçon : elle peut sauver ta vie et celle de tes proches.

\courssub{Le courant domestique est dangereux pour l'être humain}

\textbf{Une observation qui doit faire réfléchir.}\par

Tu as peut-être déjà ressenti une « secousse » en touchant un appareil défectueux ou une prise mal isolée. Ce picotement désagréable est un avertissement : le courant a traversé ton corps. Pourquoi ? Parce que \cle{le corps humain est conducteur} : il est constitué à environ 60 \% d'eau contenant des sels dissous (des ions), et l'eau salée conduit le courant électrique.

\begin{attention}[L'eau conduit le courant]
Contrairement à une idée répandue, même l'eau du robinet conduit le courant électrique, car elle contient des sels minéraux dissous. Seule l'eau parfaitement pure (distillée) est très mauvaise conductrice. Conséquence pratique : \textbf{la peau mouillée laisse passer le courant beaucoup plus facilement que la peau sèche}. Sèche-toi toujours les mains avant de toucher une prise, un interrupteur ou un appareil électrique.
\end{attention}

\textbf{Les effets du courant sur le corps humain.}\par

Les médecins et les secouristes ont constaté, à partir de très nombreux accidents, que le passage du courant à travers le corps provoque des effets d'autant plus graves que l'intensité est grande et que la durée de contact est longue :

\begin{itemize}
\item les \cle{brûlures} : le courant échauffe les tissus qu'il traverse (effet Joule), surtout aux points d'entrée et de sortie du courant ;
\item la \cle{tétanisation} des muscles : le courant provoque des contractions musculaires violentes et incontrôlables ; la victime ne peut plus lâcher l'objet conducteur qu'elle tient en main ;
\item l'\cle{arrêt cardiaque} : si le courant traverse le cœur, il perturbe son rythme et peut l'arrêter. C'est l'\cle{électrocution}, qui peut être mortelle.
\end{itemize}

\begin{aretenir}[Ce qu'il faut retenir]
La tension domestique de \SI{220}{V} est \textbf{dangereuse} pour l'être humain. Le corps humain conduit le courant, surtout lorsqu'il est mouillé. Le passage du courant peut provoquer brûlures, tétanisation des muscles et arrêt cardiaque. Il n'y a pas de « petite » prise de courant sans danger : toute installation sous tension doit être respectée.
\end{aretenir}

\courssub{Les règles de sécurité à la maison}

La plupart des accidents peuvent être évités en appliquant des règles simples, que tu dois connaître et faire respecter autour de toi.

\begin{methode}[Les gestes qui protègent]
\begin{enumerate}
\item Ne jamais toucher une prise, un interrupteur ou un appareil électrique avec les \textbf{mains mouillées} ou les pieds nus sur un sol humide.
\item Ne jamais insérer d'objet (crayon, clou, fourchette, doigt) dans les trous d'une prise.
\item Ne jamais réparer ou ouvrir un appareil encore \textbf{branché} : débrancher d'abord la fiche de la prise.
\item Ne pas utiliser un appareil dont le fil est dénudé ou abîmé : le faire réparer par une personne compétente.
\item Ne jamais utiliser d'appareil électrique (sèche-cheveux, radio branchée) près d'une baignoire ou d'un lavabo rempli d'eau.
\end{enumerate}
\end{methode}

\textbf{Un danger extérieur méconnu : les câbles tombés au sol.}\par

Après une tempête ou un accident, il arrive qu'un câble électrique tombe dans la rue ou dans la cour. Même s'il semble inerte, sans étincelle ni bruit, il peut rester \textbf{sous tension}. Le sol humide autour du câble peut lui-même conduire le courant.

\begin{attention}[Câble tombé au sol : danger mortel]
Ne jamais approcher ni toucher un câble électrique tombé au sol, même s'il paraît inerte, même avec un bâton ou un objet sec. Il faut éloigner les personnes et les animaux, puis \textbf{prévenir immédiatement un adulte ou ENEO}. Plusieurs personnes sont mortes au Cameroun en voulant ramasser ou déplacer un câble « qui ne faisait rien ».
\end{attention}

\begin{popfigure}
\begin{center}
\begin{tikzpicture}[scale=1.0]
% Pictogramme 1 : mains mouillées interdites
\draw[line width=2pt,popink] (0,0) circle (1.1);
\draw[line width=2pt,popink] (-0.78,-0.78) -- (0.78,0.78);
% prise
\draw[fill=popcream,draw=popdark,rounded corners=2pt] (-0.35,-0.55) rectangle (0.35,0.15);
\fill[popdark] (-0.13,-0.25) circle (0.05);
\fill[popdark] (0.13,-0.25) circle (0.05);
% gouttes d'eau
\draw[fill=popblue,draw=popblue] (-0.55,0.55) .. controls (-0.65,0.35) and (-0.45,0.35) .. (-0.55,0.55);
\draw[fill=popblue,draw=popblue] (0.1,0.75) .. controls (0.0,0.55) and (0.2,0.55) .. (0.1,0.75);
\node[popdark,font=\small\bfseries] at (0,-1.6) {Mains mouillées : interdit};
% Pictogramme 2 : danger haute tension
\begin{scope}[xshift=5cm]
\draw[line width=2pt,popdark,fill=popgoldL] (0,1.1) -- (-1.05,-0.75) -- (1.05,-0.75) -- cycle;
\draw[fill=popdark,draw=popdark] (0.05,0.65) -- (-0.25,0.0) -- (0.05,0.0) -- (-0.05,-0.5) -- (0.3,0.2) -- (0.0,0.2) -- cycle;
\node[popdark,font=\small\bfseries] at (0,-1.6) {Danger : haute tension};
\end{scope}
\end{tikzpicture}
\end{center}
\legende{Pictogrammes de sécurité : à gauche, il est interdit de toucher une prise avec les mains mouillées ; à droite, le symbole international de danger électrique (éclair dans un triangle) que l'on voit sur les transformateurs et les postes ENEO.}
\end{popfigure}

\courssub{Les dispositifs de protection : fusible, disjoncteur, prise de terre}

\textbf{Le problème des surintensités.}\par

Observe la multiprise du salon : télévision, décodeur, chargeurs, ventilateur... Si l'on branche trop d'appareils puissants sur un même circuit, l'intensité totale devient trop grande pour les fils : ils s'échauffent (effet Joule), l'isolant fond, et un incendie peut se déclarer. C'est une \cle{surintensité}. Elle peut aussi provenir d'un \cle{court-circuit} (contact direct entre les deux fils d'alimentation dans un appareil défectueux : l'intensité devient alors très grande).

\begin{definition}[Fusible et disjoncteur]
Un \cle{fusible} ou un \cle{disjoncteur} est un \textbf{dispositif de protection qui coupe automatiquement le circuit lorsque l'intensité du courant dépasse une valeur limite} prévue à l'avance, appelée \cle{calibre}. Le fusible contient un fil fin qui fond (il est alors détruit et doit être remplacé) ; le disjoncteur bascule et peut être réarmé après élimination de la panne.
\end{definition}

\begin{exemplebox}[Un fusible de 10 A]
Un fusible de \SI{10}{A} protège le circuit des prises d'une chambre. Tant que l'intensité reste inférieure à \SI{10}{A}, tout fonctionne normalement. Si l'on branche trop d'appareils et que l'intensité dépasse \SI{10}{A}, le fil du fusible \textbf{fond} et coupe le circuit : les appareils s'éteignent, mais les fils de l'installation ne chauffent plus et l'incendie est évité. Il faut alors débrancher des appareils et remplacer le fusible par un fusible \textbf{de même calibre} — jamais par un fil de fer !
\end{exemplebox}

\begin{attention}[Les erreurs qui tuent]
\begin{itemize}
\item \textbf{Surcharger les multiprises} en branchant fer à repasser, bouilloire et chauffe-eau sur la même rallonge : l'échauffement peut provoquer un incendie.
\item \textbf{Remplacer un fusible grillé par un fil métallique quelconque} : le circuit n'est alors plus protégé du tout.
\item Installer un fusible de calibre trop grand (par exemple \SI{30}{A} sur un circuit prévu pour \SI{10}{A}) : il ne fondra pas à temps.
\end{itemize}
\end{attention}

\textbf{La prise de terre (complément utile).}\par

Certains appareils (réfrigérateur, machine à laver, fer à repasser) ont une carcasse métallique. Si un fil interne dénudé touche cette carcasse, celle-ci se retrouve sous tension : celui qui la touche est électrisé. La \cle{prise de terre} est un fil (souvent vert-jaune) relié à une tige métallique enfoncée dans le sol. Son rôle : \textbf{évacuer les courants de fuite vers la terre} au lieu de les laisser traverser le corps de l'utilisateur. Associée à un disjoncteur adapté, elle protège efficacement les personnes.

\begin{popfigure}
\begin{center}
\begin{tikzpicture}[scale=0.95]
% Arrivée ENEO
\draw[line width=1.5pt,popdark] (0,2.5) -- (1.2,2.5);
\node[popdark,font=\small,align=center] at (0.4,3.0) {Arrivée\\ENEO \SI{220}{V}};
% Compteur
\draw[fill=popblueL,draw=popdark,rounded corners=3pt] (1.2,2.0) rectangle (2.8,3.0);
\node[popdark,font=\small\bfseries] at (2.0,2.5) {Compteur};
% Disjoncteur général
\draw[fill=poporangeL,draw=popdark,rounded corners=3pt] (3.6,2.0) rectangle (5.4,3.0);
\node[popdark,font=\small\bfseries,align=center] at (4.5,2.5) {Disjoncteur\\général};
\draw[line width=1.5pt,popdark] (2.8,2.5) -- (3.6,2.5);
% Départ vers les pièces
\draw[line width=1.5pt,popdark] (5.4,2.5) -- (7.0,2.5);
% Circuit prises
\draw[line width=1.2pt,popdark] (7.0,2.5) -- (7.0,1.2) -- (8.6,1.2);
\draw[fill=popgreenL,draw=popdark,rounded corners=2pt] (8.6,0.8) rectangle (10.0,1.6);
\node[popdark,font=\small,align=center] at (9.3,1.2) {Prise\\avec terre};
% fusible sur le circuit
\draw[fill=popgoldL,draw=popdark] (6.3,2.3) rectangle (6.9,2.7);
\node[popdark,font=\scriptsize] at (6.6,3.05) {Fusible \SI{10}{A}};
% Circuit lampe
\draw[line width=1.2pt,popdark] (7.0,2.5) -- (7.0,3.8) -- (8.6,3.8);
\draw[fill=poppurpleL,draw=popdark] (8.6,3.8) circle (0.45);
\draw[popdark,line width=1pt] (8.3,3.5) -- (8.9,4.1);
\draw[popdark,line width=1pt] (8.3,4.1) -- (8.9,3.5);
\node[popdark,font=\small] at (9.3,4.5) {Lampe};
% Prise de terre
\draw[line width=1.2pt,popgreen!60!black] (9.3,0.8) -- (9.3,-0.3);
\draw[line width=1.5pt,popgreen!60!black] (8.8,-0.3) -- (9.8,-0.3);
\draw[line width=1.2pt,popgreen!60!black] (8.95,-0.5) -- (9.65,-0.5);
\draw[line width=1pt,popgreen!60!black] (9.1,-0.7) -- (9.5,-0.7);
\node[popgreen!60!black,font=\small] at (10.6,-0.45) {Terre (sol)};
\end{tikzpicture}
\end{center}
\legende{Schéma simplifié d'une installation domestique sécurisée : le courant ENEO traverse le compteur (qui mesure la consommation), puis le disjoncteur général (protection de toute la maison), avant d'alimenter les circuits (prises protégées par un fusible adapté, éclairage). La prise de terre évacue les courants de fuite vers le sol.}
\end{popfigure}

\begin{savaistu}[Les installations sauvages tuent]
Au Cameroun, une grande partie des accidents électriques domestiques est liée à des \textbf{installations vétustes} (fils dénudés, prises cassées) ou à des \textbf{branchements sauvages} réalisés sans normes, parfois directement sur les lignes ENEO. Ces branchements frauduleux provoquent électrocutions et incendies. Le respect des normes d'installation, l'emploi de matériel de qualité et l'intervention de techniciens qualifiés sauvent littéralement des vies. Signale à un adulte toute installation dangereuse que tu observes !
\end{savaistu}

\courssub{Conduite à tenir en cas d'accident électrique}

Une personne de ta famille touche un fil dénudé et reste « collée », secouée de tremblements : que faire ? Agir vite, mais \textbf{dans le bon ordre}, car un secouriste maladroit devient lui-même une victime.

\begin{methode}[Secourir une victime d'électrocution]
\begin{enumerate}
\item \textbf{Couper le courant d'abord} : actionner le disjoncteur général ou débrancher la fiche. C'est le premier geste, avant tout contact.
\item \textbf{Ne jamais toucher la victime directement} tant qu'elle est encore en contact avec le circuit : le courant traverserait aussi ton corps. Si l'on ne peut pas couper le courant, repousser la victime avec un objet \textbf{sec et isolant} (bâton en bois sec, chaise en plastique).
\item \textbf{Secourir ensuite} : une fois le courant coupé, allonger la victime, vérifier sa respiration, la placer sur le côté si elle respire.
\item \textbf{Alerter les secours} : appeler les pompiers (numéro d'urgence 118 au Cameroun) ou conduire rapidement la victime à l'hôpital, même si elle semble remise : un examen médical est toujours nécessaire.
\end{enumerate}
\end{methode}

\begin{exoresolu}[Que faire dans cette situation ?]
Énoncé : pendant la pluie, ton petit frère veut brancher le chargeur du téléphone avec les mains mouillées. Au même moment, tu remarques que la multiprise du salon alimente déjà le fer à repasser, la bouilloire et le ventilateur, et que le fusible de ce circuit a été remplacé par un fil de fer. Relever trois dangers et donner la conduite correcte pour chacun.
\tcblower
Solution détaillée :
\begin{enumerate}
\item \textbf{Mains mouillées sur la prise} : la peau humide conduit beaucoup mieux le courant ; risque d'électrisation, voire d'électrocution. Conduite : lui faire sécher les mains avant toute manipulation (ou brancher soi-même le chargeur).
\item \textbf{Multiprise surchargée} : fer à repasser, bouilloire et ventilateur sont des appareils puissants ; l'intensité totale peut dépasser la limite du circuit, d'où échauffement des fils et risque d'incendie. Conduite : débrancher une partie des appareils et répartir les branchements sur des prises différentes.
\item \textbf{Fusible remplacé par un fil de fer} : le circuit n'a plus aucune protection contre les surintensités. Conduite : faire remplacer le fil de fer par un fusible du calibre prévu (par exemple \SI{10}{A}) par une personne compétente, courant coupé.
\end{enumerate}
\end{exoresolu}

\begin{aretenir}[L'essentiel de la sécurité électrique]
\begin{itemize}
\item Le courant domestique de \SI{220}{V} est dangereux ; le corps humain, surtout mouillé, conduit le courant.
\item Effets du courant : brûlures, tétanisation, arrêt cardiaque.
\item Mains sèches, aucun objet dans les prises, aucun appareil réparé sous tension, aucun câble au sol approché.
\item Fusibles et disjoncteurs adaptés : ils coupent le circuit en cas de surintensité ; la prise de terre évacue les courants de fuite.
\item En cas d'accident : \textbf{couper le courant d'abord, secourir ensuite, alerter les secours}.
\end{itemize}
\end{aretenir}

\begin{exercice}[À toi de jouer]
\begin{enumerate}
\item Pourquoi est-il plus dangereux de toucher une prise avec les mains mouillées qu'avec les mains sèches ?
\item Citer les trois effets du courant électrique sur le corps humain.
\item Un fusible de \SI{16}{A} protège le circuit d'une cuisine. Expliquer ce qui se passe si un court-circuit y produit une intensité de \SI{50}{A}.
\item Ton voisin veut déplacer un câble ENEO tombé dans sa cour avec un bâton en bois humide. Quels conseils lui donnes-tu ?
\end{enumerate}
\end{exercice}

\begin{corrige}[Exercice : À toi de jouer]
\begin{enumerate}
\item L'eau (même du robinet) contient des sels dissous et conduit le courant : la peau mouillée laisse passer le courant beaucoup plus facilement que la peau sèche, donc l'intensité qui traverse le corps est plus grande et les effets plus graves.
\item Brûlures, tétanisation des muscles, arrêt cardiaque (électrocution).
\item L'intensité dépasse largement le calibre du fusible : le fil du fusible fond presque instantanément et coupe le circuit, ce qui protège l'installation contre l'incendie. Il faudra réparer la cause du court-circuit avant de remplacer le fusible par un fusible de \SI{16}{A}.
\item Le bois humide conduit le courant : il risque l'électrocution. Il faut lui dire de ne pas approcher le câble, d'éloigner les personnes et les animaux, et de prévenir immédiatement ENEO ou un adulte responsable.
\end{enumerate}
\end{corrige}

\newpage

\coursec{Activité d'intégration}

\coursec{Leçon N12 : Le courant alternatif}

\courssub{Objectifs de l'activité}

À la fin de cette activité, tu seras capable de :
\begin{itemize}
\item exploiter un oscillogramme pour déterminer la période $T$, la fréquence $f$ et la valeur maximale $U_{\max}$ d'une tension alternative ;
\item vérifier qu'une tension correspond au réseau domestique camerounais ($\SI{220}{V}$ ; $\SI{50}{Hz}$) ;
\item compléter le schéma de la chaîne de production, de transport et de distribution de l'énergie électrique en plaçant correctement les transformateurs ;
\item justifier le transport en haute tension par un calcul d'intensité à puissance égale ;
\item calculer une consommation d'énergie en kilowattheures et le montant d'une facture ENEO ;
\item formuler des consignes de sécurité et d'économie d'énergie adaptées à une habitation.
\end{itemize}

\courssub{Situation}

\textbf{Le voyage d'un kilowattheure : de Songloulou à ta maison.}

La famille de Marie habite un quartier de Yaoundé. En regardant le compteur ENEO fixé au mur de la maison, son père lui pose une question : « Sais-tu d'où vient l'électricité que tu utilises chaque soir pour faire tes devoirs ? » Marie sait que le courant du réseau est un courant \cle{alternatif}, mais elle ignore tout du chemin parcouru par l'énergie électrique. Pour répondre à son père, elle rassemble trois documents.

\textbf{Document 1 : oscillogramme de la tension du réseau.} Un technicien a branché un oscilloscope aux bornes d'une prise (avec un dispositif de sécurité adapté). Réglages : balayage horizontal $b = \SI{5}{ms/div}$, sensibilité verticale $s = \SI{100}{V/div}$. La courbe observée est une sinusoïde : une période complète occupe $4$ divisions horizontales et la crête atteint $\num{3,1}$ divisions verticales.

\textbf{Document 2 : extrait de carte simplifiée.} La centrale hydroélectrique de \cle{Songloulou} (sur le fleuve Sanaga) produit de l'énergie électrique. Une ligne \cle{haute tension} de $\SI{225}{kV}$ achemine cette énergie jusqu'à un poste de transformation situé à l'entrée de la ville. De là, des lignes moyenne tension alimentent le transformateur du quartier, qui dessert les maisons en $\SI{220}{V}$.

\textbf{Document 3 : facture ENEO simplifiée du mois.} Ancien index du compteur : $\SI{1240}{kWh}$ ; nouvel index : $\SI{1420}{kWh}$ ; prix du kilowattheure : $\SI{50}{FCFA/kWh}$.

\begin{popfigure}
\begin{center}
\begin{tikzpicture}[scale=0.9, >=Stealth]
% Centrale
\draw[fill=popblueL, draw=popblue] (0,0) rectangle (2.2,1.4);
\node[align=center] at (1.1,0.7) {\small Centrale de\\ \small Songloulou};
% Ligne HT
\draw[very thick, poporange, ->] (2.2,0.7) -- (5.2,0.7) node[midway, above, black]{\small Ligne $\SI{225}{kV}$};
% Poste
\draw[fill=popgoldL, draw=popgold] (5.2,0) rectangle (7.4,1.4);
\node[align=center] at (6.3,0.7) {\small Poste de\\ \small transformation};
% Ligne MT
\draw[very thick, popgreen, ->] (7.4,0.7) -- (9.6,0.7) node[midway, above, black]{\small Moyenne tension};
% Transfo quartier
\draw[fill=poppurpleL, draw=poppurple] (9.6,0) rectangle (11.6,1.4);
\node[align=center] at (10.6,0.7) {\small Transfo du\\ \small quartier};
% Maison
\draw[very thick, poppink, ->] (11.6,0.7) -- (13.4,0.7);
\draw[fill=popcream, draw=popdark] (13.4,0) rectangle (15.2,1.4);
\node[align=center] at (14.3,0.7) {\small Maison\\ \small $\SI{220}{V}$};
\end{tikzpicture}
\end{center}
\legende{Chaîne simplifiée de l'énergie électrique : centrale, transport, distribution, habitation (schéma de base à compléter).}
\end{popfigure}

\courssub{Consignes}

\begin{enumerate}
\item \textbf{Exploitation de l'oscillogramme.} À partir du document 1, détermine la période $T$ de la tension, puis calcule sa fréquence $f$. Détermine ensuite la valeur maximale $U_{\max}$ de la tension.
\item \textbf{Identification du réseau.} Calcule la valeur efficace $U$ de la tension à l'aide de la relation $U = \dfrac{U_{\max}}{\sqrt{2}}$ (avec $\sqrt{2} \approx \num{1,41}$). Conclus : s'agit-il bien du réseau domestique camerounais $\SI{220}{V}$ ; $\SI{50}{Hz}$ ?
\item \textbf{Schéma de la chaîne de distribution.} Recopie le schéma ci-dessus et place, aux bons endroits, un transformateur \cle{élévateur} (à la sortie de la centrale) et des transformateurs \cle{abaisseurs} (au poste de transformation et dans le quartier). Justifie chaque choix en une phrase.
\item \textbf{Pourquoi transporter en haute tension ?} La centrale fournit une puissance $P = \SI{3960}{kW}$ à la ligne. Calcule l'intensité $I$ du courant transporté dans deux cas : (a) si le transport se faisait sous $U_1 = \SI{220}{V}$ ; (b) sous $U_2 = \SI{225}{kV}$. Compare les deux résultats et explique pourquoi une intensité faible est préférable pour limiter les pertes par échauffement (effet Joule) dans les câbles.
\item \textbf{La facture de la famille.} Calcule la consommation d'énergie du mois en kilowattheures, puis le montant à payer.
\item \textbf{Conseils à afficher.} Rédige trois conseils clairs et courts : au moins un conseil de \cle{sécurité} et un conseil d'\cle{économie d'énergie}, à afficher dans la maison de Marie.
\item \textbf{Restitution.} Par groupe, présentez tous vos résultats sous forme d'une affiche : oscillogramme commenté, schéma de la chaîne complété, calculs de la facture et les trois conseils illustrés.
\end{enumerate}

\courssub{Ressources à mobiliser}

\begin{definition}[Courant continu et courant alternatif sinusoïdal]
Un \textbf{courant continu} (noté $=$) a une intensité constante dans le temps (pile, batterie). Un \textbf{courant alternatif} (noté $\sim$) a une intensité qui varie périodiquement (prise secteur). S'il varie suivant une sinusoïde, on parle de \textbf{courant alternatif sinusoïdal}. Il est caractérisé par :
\begin{itemize}
\item sa \textbf{période} $T$ (en $\SI{}{s}$) : durée d'un motif complet ;
\item sa \textbf{fréquence} $f$ (en $\SI{}{Hz}$) : nombre de périodes par seconde, $f = \dfrac{1}{T}$ ;
\item son \textbf{amplitude} (ou valeur maximale) $U_{\max}$ (en $\SI{}{V}$) ;
\item sa \textbf{valeur efficace} $U$ (en $\SI{}{V}$) : valeur du courant continu qui produirait le même effet thermique ; $U = \dfrac{U_{\max}}{\sqrt{2}}$.
\end{itemize}
Au Cameroun, le réseau domestique délivre une tension alternative sinusoïdale de valeur efficace $U = \SI{220}{V}$ et de fréquence $f = \SI{50}{Hz}$.
\end{definition}

\begin{aretenir}[Formules utiles]
\begin{itemize}
\item Période et fréquence : $T = b \times x$ (balayage $\times$ nombre de divisions) et $f = \dfrac{1}{T}$.
\item Tension maximale : $U_{\max} = s \times y$ (sensibilité $\times$ nombre de divisions) ; tension efficace : $U = \dfrac{U_{\max}}{\sqrt{2}}$.
\item Puissance électrique : $P = U \times I$, donc $I = \dfrac{P}{U}$.
\item Énergie consommée : $E = \text{nouvel index} - \text{ancien index}$ ; montant $= E \times \text{prix du kWh}$.
\item Rôle des transformateurs : \cle{élévateur} en sortie de centrale (augmente $U$, diminue $I$), \cle{abaisseur} près des utilisateurs (diminue $U$ pour la sécurité).
\end{itemize}
\end{aretenir}

\begin{savaistu}[Le saviez-vous ? La Sanaga, poumon énergétique du Cameroun]
Le fleuve Sanaga, avec ses chutes (Songloulou, Édéa, Nachtigal), fournit plus de $\SI{90}{\%}$ de l'électricité d'origine hydraulique du pays. La centrale de \textbf{Songloulou} (mise en service en 1980, $\SI{384}{MW}$) et celle de \textbf{Nachtigal} (en service depuis 2023-2024, $\SI{420}{MW}$) sont les piliers de la production d'ENEO. L'électricité y est produite en \textbf{courant alternatif triphasé} à $\SI{15}{kV}$ ou $\SI{20}{kV}$, puis élevée à $\SI{225}{kV}$ (voire $\SI{400}{kV}$ pour les futures lignes) pour le transport sur de longues distances vers Yaoundé, Douala et le Nord.
\end{savaistu}

\courssub{Démarche et corrigé commenté}

\begin{exoresolu}[Le voyage d'un kilowattheure : corrigé complet]
\textbf{Question 1.} Déterminer $T$, $f$ et $U_{\max}$ à partir de l'oscillogramme (balayage $\SI{5}{ms/div}$, sensibilité $\SI{100}{V/div}$, période sur $4$ divisions, crête à $\num{3,1}$ divisions).

\textbf{Question 2.} Calculer $U = U_{\max}/\sqrt{2}$ et identifier le réseau.

\textbf{Question 3.} Placer les transformateurs sur la chaîne et justifier.

\textbf{Question 4.} Calculer $I$ sous $\SI{220}{V}$ puis sous $\SI{225}{kV}$ pour $P = \SI{3960}{kW}$, et conclure.

\textbf{Question 5.} Calculer la consommation et le montant de la facture.

\textbf{Question 6.} Rédiger trois conseils de sécurité et d'économie d'énergie.
\tcblower
\textbf{Étape 1 : exploitation de l'oscillogramme.}

La période est la durée d'un motif complet. Une période occupe $x = 4$ divisions avec un balayage $b = \SI{5}{ms/div}$ :
\[
T = b \times x = 5 \times 4 = \SI{20}{ms} = \num{0,02}\ \SI{}{s}.
\]
La fréquence est l'inverse de la période :
\[
f = \frac{1}{T} = \frac{1}{\num{0,02}} = \SI{50}{Hz}.
\]
La crête atteint $y = \num{3,1}$ divisions avec une sensibilité $s = \SI{100}{V/div}$ :
\[
U_{\max} = s \times y = 100 \times \num{3,1} = \SI{310}{V}.
\]

\textbf{Étape 2 : identification du réseau.}
\[
U = \frac{U_{\max}}{\sqrt{2}} = \frac{310}{\num{1,41}} \approx \SI{220}{V}.
\]
La tension efficace est $\SI{220}{V}$ et la fréquence $\SI{50}{Hz}$ : c'est bien la tension alternative sinusoïdale du réseau domestique camerounais.

\textbf{Étape 3 : les transformateurs sur la chaîne.}
\begin{itemize}
\item À la sortie de la centrale : un transformateur \textbf{élévateur} porte la tension à $\SI{225}{kV}$ pour le transport sur longue distance.
\item Au poste de transformation : un transformateur \textbf{abaisseur} ramène la tension à la moyenne tension (ex. $\SI{30}{kV}$ ou $\SI{15}{kV}$).
\item Dans le quartier : un transformateur \textbf{abaisseur} fournit le $\SI{220}{V}$ utilisable sans danger relatif dans les maisons.
\end{itemize}
Justification : on élève la tension pour transporter avec une faible intensité (moins de pertes par effet Joule), puis on l'abaisse par paliers pour la rendre utilisable en sécurité.

\begin{popfigure}
\begin{center}
\begin{tikzpicture}[scale=0.9, >=Stealth, every node/.style={font=\small}]
% Centrale
\draw[fill=popblueL, draw=popblue] (0,0) rectangle (2.2,1.4);
\node[align=center] at (1.1,0.7) {Centrale de\\ Songloulou};
% Transfo élévateur
\draw[fill=poporangeL, draw=poporange, thick] (2.4,0.1) rectangle (3.6,1.3);
\node[align=center] at (3.0,0.7) {Transfo\\ \textbf{élévateur}};
\node[align=center, poporange] at (3.0, -0.5) {$\SI{15}{kV} \to \SI{225}{kV}$};
% Ligne HT
\draw[very thick, poporange, ->] (3.6,0.7) -- (5.8,0.7) node[midway, above, black]{Ligne $\SI{225}{kV}$};
% Poste
\draw[fill=popgoldL, draw=popgold] (5.8,0) rectangle (8.0,1.4);
\node[align=center] at (6.9,1.0) {Poste de};
\node[align=center] at (6.9,0.4) {transformation};
% Transfo abaisseur HT/MT
\draw[fill=popgreenL, draw=popgreen, thick] (8.2,0.1) rectangle (9.4,1.3);
\node[align=center] at (8.8,0.7) {Transfo\\ \textbf{abaisseur}};
\node[align=center, popgreen] at (8.8, -0.5) {$\SI{225}{kV} \to \SI{30}{kV}$};
% Ligne MT
\draw[very thick, popgreen, ->] (9.4,0.7) -- (11.6,0.7) node[midway, above, black]{Moyenne tension};
% Transfo quartier
\draw[fill=poppurpleL, draw=poppurple] (11.6,0) rectangle (13.8,1.4);
\node[align=center] at (12.7,1.0) {Transfo du};
\node[align=center] at (12.7,0.4) {quartier};
% Transfo abaisseur MT/BT
\draw[fill=poppinkL, draw=poppink, thick] (14.0,0.1) rectangle (15.2,1.3);
\node[align=center] at (14.6,0.7) {Transfo\\ \textbf{abaisseur}};
\node[align=center, poppink] at (14.6, -0.5) {$\SI{30}{kV} \to \SI{220}{V}$};
% Maison
\draw[very thick, poppink, ->] (15.2,0.7) -- (16.5,0.7);
\draw[fill=popcream, draw=popdark] (16.5,0) rectangle (18.3,1.4);
\node[align=center] at (17.4,0.7) {Maison\\ $\SI{220}{V}$};
\end{tikzpicture}
\end{center}
\legende{Chaîne complète : production ($\sim\SI{15}{kV}$), élévation ($\SI{225}{kV}$), transport HT, abaisseur HT/MT, distribution MT, abaisseur MT/BT ($\SI{220}{V}$), utilisation.}
\end{popfigure}

\textbf{Étape 4 : pourquoi la haute tension ?}

On convertit la puissance en watts : $P = \SI{3960}{kW} = \num{3960000}\ \SI{}{W}$.
\begin{align*}
\text{(a) Sous } U_1 = \SI{220}{V} : \quad I_1 &= \frac{P}{U_1} = \frac{\num{3960000}}{220} = \SI{18000}{A}.\\
\text{(b) Sous } U_2 = \SI{225}{kV} = \num{225000}\ \SI{}{V} : \quad I_2 &= \frac{P}{U_2} = \frac{\num{3960000}}{\num{225000}} = \SI{17,6}{A}.
\end{align*}
L'intensité est environ \textbf{mille fois plus faible} en haute tension ($\SI{17,6}{A}$ contre $\SI{18000}{A}$). Or les pertes par effet Joule (échauffement des câbles) sont proportionnelles au carré de l'intensité ($P_{\text{pertes}} = R \cdot I^2$). Transporter sous $\SI{225}{kV}$ permet donc de réduire considérablement les pertes d'énergie et l'échauffement des lignes, ce qui rend le transport rentable sur de grandes distances.

\textbf{Étape 5 : la facture.}
\[
E = \SI{1420}{kWh} - \SI{1240}{kWh} = \SI{180}{kWh}.
\]
\[
\text{Montant} = \SI{180}{kWh} \times \SI{50}{FCFA/kWh} = \SI{9000}{FCFA}.
\]
La famille a consommé $\SI{180}{kWh}$ et doit payer $\SI{9000}{FCFA}$ (hors taxes et redevances fixes).

\textbf{Étape 6 : trois conseils à afficher.}
\begin{enumerate}
\item \textbf{Sécurité :} Ne jamais toucher une prise ou un appareil électrique avec les mains mouillées, et ne pas brancher trop d'appareils sur une même rallonge (risque de surchauffe et d'incendie).
\item \textbf{Sécurité :} Couper le disjoncteur général (ou le différentiel) avant toute intervention sur l'installation, et ne jamais réparer soi-même un fil dénudé : appeler un électricien qualifié.
\item \textbf{Économie :} Éteindre systématiquement les lampes, le téléviseur et le ventilateur en quittant une pièce ; débrancher les chargeurs inutilisés : chaque kilowattheure économisé allège la facture et préserve les ressources.
\end{enumerate}
\end{exoresolu}

\begin{attention}[Pièges fréquents]
\begin{itemize}
\item Ne pas oublier de convertir les millisecondes en secondes avant de calculer $f = 1/T$ : $\SI{20}{ms} = \num{0,02}\ \SI{}{s}$.
\item Bien convertir les kilowatts en watts ($\times 1000$) et les kilovolts en volts ($\times 1000$) avant d'appliquer $I = P/U$ ; sinon l'intensité est fausse de facteur $1000$.
\item La valeur lue directement sur l'oscillogramme (crête) est $U_{\max}$, pas la tension efficace : il faut diviser par $\sqrt{2}$ pour retrouver les $\SI{220}{V}$ du réseau.
\item La consommation se calcule par différence des index (nouvel $-$ ancien), pas en lisant seulement le nouvel index.
\item Sur un schéma, le transformateur élévateur est \textbf{avant} la ligne haute tension ; les abaisseurs sont \textbf{après} (poste source, puis poste de quartier).
\end{itemize}
\end{attention}

\courssub{Bilan de l'activité}

Cette activité t'a permis de relier toutes les notions de la leçon dans une situation réelle. Tu as d'abord \textbf{mesuré} : l'exploitation de l'oscillogramme a montré que la tension du réseau est alternative sinusoïdale, de période $\SI{20}{ms}$, de fréquence $\SI{50}{Hz}$ et de valeur efficace $\SI{220}{V}$. Tu as ensuite \textbf{schématisé} la chaîne complète de l'énergie : la centrale de Songloulou produit l'électricité en $\SI{15}{kV}$, un transformateur élévateur permet son transport sous $\SI{225}{kV}$ (haute tension), puis des transformateurs abaisseurs (HT/MT au poste source, MT/BT au quartier) la distribuent jusqu'aux habitations en $\SI{220}{V}$. Le calcul des intensités à puissance égale a \textbf{justifié} scientifiquement le choix de la haute tension : plus la tension est élevée, plus l'intensité est faible, et moins les pertes par effet Joule ($R I^2$) sont importantes. Enfin, le calcul de la facture et la rédaction des conseils t'ont montré que comprendre l'électricité aide à la \textbf{consommer en sécurité} et à l'\textbf{économiser}. Le kilowattheure affiché sur ton compteur ENEO a donc fait un long voyage : du fleuve Sanaga jusqu'à ta lampe de bureau.

\newpage

\coursec{Méthodes et exercices résolus}

\coursec{Courant continu et courant alternatif}

\courssub{Définitions et symboles normalisés}

\begin{definition}[Courant continu]
Un \cle{courant continu} est un courant dont la direction (le sens de circulation) et l'intensité sont constantes dans le temps. Il est fourni par des générateurs électrochimiques : \cle{pile}, \cle{batterie}, \cle{accumulateur}.
\textbf{Symbole :} $\mathbf{=}$ (ou une ligne pleine au-dessus d'une ligne pointillée).
\end{definition}

\begin{definition}[Courant alternatif]
Un \cle{courant alternatif} est un courant dont la direction change de façon périodique : il circule alternativement dans un sens puis dans l'autre. Son intensité varie dans le temps selon une loi souvent \cle{sinusoïdale}. Il est fourni par le \cle{réseau de distribution publique} (ENEO au Cameroun) et par les \cle{alternateurs}.
\textbf{Symbole :} $\sim$
\end{definition}

\begin{aretenir}[Repérage rapide]
\begin{itemize}
\item Appareil à piles / batterie (lampe torche, téléphone, moto au démarreur) $\rightarrow$ \textbf{Courant continu} ($\mathbf{=}$).
\item Appareil branché sur prise murale ENEO (ventilateur, fer, TV, réfrigérateur) $\rightarrow$ \textbf{Courant alternatif} ($\sim$).
\end{itemize}
\end{aretenir}

\courssub{Observation à l'oscilloscope : expérience de cours}

\begin{experience}[Visualiser la nature d'une tension]
\textbf{Matériel :} oscilloscope, pile plate \SI{4,5}{V} (ou accumulateur), transformateur de sécurité basse tension (sortie \SI{6}{V} $\sim$), fils de connexion.
\textbf{Protocole :}
\begin{enumerate}
\item Régler l'oscilloscope : balayage \SI{5}{ms/div}, sensibilité \SI{2}{V/div}, couplage DC.
\item Brancher la pile aux bornes de l'entrée Y. Observer le trait. Noter sa position verticale. Inverser les bornes de la pile : le trait saute de l'autre côté du zéro.
\item Brancher le transformateur BT (secondaire) aux bornes de l'entrée Y. Observer la courbe. Inverser les bornes : la courbe ne change pas d'aspect.
\end{enumerate}
\textbf{Observations et interprétation :}
\begin{itemize}
\item \textbf{Pile :} un trait horizontal stable. La tension est constante : $u(t) = U_0$. C'est une \cle{tension continue}. L'inversion des bornes décale le trait (polarité).
\item \textbf{Transformateur :} une courbe en « vagues » régulières (sinusoïde) centrée sur le zéro. La tension change de signe périodiquement : $u(t) = U_{max} \sin(2\pi f t)$. C'est une \cle{tension alternative}.
\end{itemize}
\legende{À gauche : tension continue (pile) -- trait horizontal. À droite : tension alternative (secteur) -- sinusoïde centrée sur 0.}
\begin{popfigure}
\begin{tikzpicture}[scale=0.8, >=Stealth]
% Grille commune
\def\drawgrid{
  \draw[popmuted, step=1cm] (0,-2) grid (6,2);
  \draw[popline, thick] (0,0) -- (6,0) node[right] {$t$};
  \draw[popline, thick] (0,-2) -- (0,2) node[above] {$u$ (V)};
}
% Gauche : Continu
\begin{scope}[xshift=0cm]
  \drawgrid
  \draw[popblue, very thick] (0,1) -- (6,1);
  \node[popblue, right] at (6,1) {$U_0$};
  \node[popdark, align=center, font=\bfseries\small] at (3, -2.5) {Tension continue (pile)};
\end{scope}
% Droite : Alternatif
\begin{scope}[xshift=8cm]
  \drawgrid
  \draw[poporange, very thick, domain=0:6, samples=200] plot (\x, {1.5*sin(360*\x)});
  \draw[poporange, dashed] (0,1.5) -- (6,1.5) node[right] {$U_{max}$};
  \draw[poporange, dashed] (0,-1.5) -- (6,-1.5) node[right] {$-U_{max}$};
  \draw[<->, popdark] (0.5, -1.8) -- (2.5, -1.8) node[midway, below] {$T$};
  \node[popdark, align=center, font=\bfseries\small] at (3, -2.5) {Tension alternative (secteur)};
\end{scope}
\end{tikzpicture}
\legende{Visualisation à l'oscilloscope : tension continue (pile) et tension alternative sinusoïdale (secteur).}
\end{popfigure}
\end{experience}

\courssub{Méthode : Distinguer courant continu et courant alternatif}

\begin{methode}[Identifier la nature d'un courant ou d'une tension]
Pour savoir si une tension (ou un courant) est continue ou alternative :
\begin{enumerate}
\item \textbf{Observer l'appareil qui la produit :} une pile, une batterie ou un accumulateur fournit un courant \cle{continu} ; le secteur ENEO (prises de courant) fournit un courant \cle{alternatif}.
\item \textbf{Observer le symbole :} le symbole $\mathbf{=}$ (ou une ligne pleine et une ligne pointillée) indique le continu ; le symbole $\sim$ indique l'alternatif.
\item \textbf{Observer l'écran d'un oscilloscope :} une tension continue donne une droite horizontale ; une tension alternative donne une courbe qui monte et descend régulièrement (sinusoïde).
\item \textbf{Vérifier le sens du courant :} en continu, le courant circule toujours dans le même sens ; en alternatif, il change de sens périodiquement.
\end{enumerate}
\textbf{Réflexe :} à la maison, tout appareil branché sur une prise fonctionne en alternatif ; tout appareil à piles (lampe torche, téléphone) fonctionne en continu.
\end{methode}

\begin{exoresolu}[Pile ou prise de courant ?]
Un élève de 3ème observe deux appareils : une lampe torche fonctionnant avec deux piles et un ventilateur branché sur une prise ENEO.
\begin{enumerate}
\item Préciser la nature du courant utilisé par chaque appareil. Justifier.
\item Sur l'étiquette du ventilateur, on lit « 220 V $\sim$ 50 Hz ». Que signifie le symbole $\sim$ ?
\end{enumerate}
\tcblower
\textbf{1. Nature des courants.}

La lampe torche est alimentée par des \cle{piles}. Or une pile est un générateur de tension \cle{continue} : le courant qu'elle fournit circule toujours dans le même sens, du pôle positif vers le pôle négatif à l'extérieur de la pile. La lampe torche utilise donc un \textbf{courant continu}.

Le ventilateur est branché sur une prise du réseau ENEO. Le réseau de distribution publique fournit une tension \cle{alternative} : le courant change de sens de façon périodique. Le ventilateur utilise donc un \textbf{courant alternatif}.

\textbf{2. Signification du symbole $\sim$.}

Le symbole $\sim$ est le symbole normalisé du courant (ou de la tension) \textbf{alternatif}. L'inscription « 220 V $\sim$ 50 Hz » signifie donc que le ventilateur doit être alimenté par une tension alternative de valeur efficace \SI{220}{V} et de fréquence \SI{50}{Hz}, c'est-à-dire exactement les caractéristiques du réseau ENEO au Cameroun.

\textbf{Conclusion :} la lampe torche fonctionne en courant continu (piles), le ventilateur en courant alternatif (secteur ENEO, 220 V $\sim$ 50 Hz).
\end{exoresolu}

\coursec{Caractéristiques d'une tension alternative sinusoïdale}

\courssub{Grandeurs caractéristiques : définitions}

\begin{definition}[Période $T$]
Durée d'un motif complet de la courbe (le plus petit intervalle de temps après lequel la tension reprend la même valeur et la même variation). Unité : le \cle{seconde} (s) ou ses sous-multiples (ms, $\mu$s).
\end{definition}

\begin{definition}[Fréquence $f$]
Nombre de périodes effectuées en une seconde. $f = \dfrac{1}{T}$. Unité : le \cle{hertz} (Hz).
\end{definition}

\begin{definition}[Tension maximale (amplitude) $U_{max}$]
Valeur extrême atteinte par la tension (hauteur du pic par rapport à l'axe zéro). Unité : le \cle{volt} (V).
\end{definition}

\begin{definition}[Tension efficace $U_{eff}$]
Valeur de la tension continue qui produirait le même effet thermique (même puissance dissipée dans une résistance) que la tension alternative. C'est la valeur mesurée par un voltmètre et inscrite sur les appareils.
Relation : $U_{eff} = \dfrac{U_{max}}{\sqrt{2}} \approx \dfrac{U_{max}}{1,414}$.
\end{definition}

\begin{aretenir}[Valeurs nominales du réseau ENEO au Cameroun]
\[ U_{eff} = \SI{220}{V} \quad ; \quad f = \SI{50}{Hz} \quad ; \quad T = \dfrac{1}{50} = \SI{0,020}{s} = \SI{20}{ms} \quad ; \quad U_{max} = 220 \times \sqrt{2} \approx \SI{311}{V}. \]
\end{aretenir}

\courssub{Exploitation d'un oscillogramme : expérience guidée}

\begin{experience}[Mesurer $T$, $f$, $U_{max}$, $U_{eff}$ sur oscilloscope]
\textbf{Matériel :} oscilloscope, transformateur de sécurité BT (\SI{6}{V} $\sim$ ou \SI{12}{V} $\sim$).
\textbf{Protocole :}
\begin{enumerate}
\item Brancher le secondaire du transformateur sur l'entrée Y (couplage AC).
\item Régler le \textbf{balayage horizontal} (time/div) pour voir 2 à 3 périodes complètes.
\item Régler la \textbf{sensibilité verticale} (volt/div) pour que la sinusoïde occupe l'écran.
\item \textbf{Mesurer :}
  \begin{itemize}
  \item $n_T$ = nombre de divisions horizontales pour \textbf{une} période complète.
  \item $n_U$ = nombre de divisions verticales du zéro au sommet (crête).
  \end{itemize}
\item \textbf{Calculer :}
  \begin{align*}
  T &= n_T \times \text{balayage (s/div)} \\
  f &= \dfrac{1}{T} \\
  U_{max} &= n_U \times \text{sensibilité (V/div)} \\
  U_{eff} &= \dfrac{U_{max}}{\sqrt{2}}
  \end{align*}
\end{enumerate}
\textbf{Compétence :} savoir lire un oscillogramme et convertir les divisions en grandeurs physiques (attention aux unités : ms $\rightarrow$ s).
\end{experience}

\courssub{Méthode : Calculer période, fréquence et valeurs d'une tension alternative}

\begin{methode}[Exploiter les caractéristiques d'une tension sinusoïdale]
\begin{enumerate}
\item \textbf{Période $T$} : c'est la durée d'un motif complet de la courbe. Sur l'oscillogramme : $T = (\text{nombre de divisions d'un motif}) \times (\text{balayage en s/div})$.
\item \textbf{Fréquence $f$} : c'est le nombre de périodes par seconde : \[ f = \dfrac{1}{T} \qquad \text{avec } T \text{ en s et } f \text{ en Hz.} \]
\item \textbf{Tension maximale} (amplitude) $U_{max}$ : hauteur du pic : $U_{max} = (\text{nombre de divisions}) \times (\text{sensibilité en V/div})$.
\item \textbf{Tension efficace} : c'est la valeur indiquée par le voltmètre et inscrite sur les appareils : \[ U_{eff} = \dfrac{U_{max}}{\sqrt{2}} \qquad \text{avec } \sqrt{2} \approx 1{,}414. \]
\item \textbf{Toujours convertir les durées en secondes avant de calculer $f$.}
\end{enumerate}
\textbf{À retenir :} au Cameroun, le réseau ENEO a pour valeurs $U_{eff} = \SI{220}{V}$ et $f = \SI{50}{Hz}$, donc $T = \SI{0,02}{s}$.
\end{methode}

\begin{exoresolu}[Étude de la tension du secteur ENEO]
On branche un oscilloscope aux bornes d'une prise de courant. On obtient une sinusoïde dont un motif complet occupe 4 divisions horizontales, avec un balayage de \SI{5}{ms/div}. Le sommet de la courbe atteint 3,1 divisions verticales, avec une sensibilité de \SI{100}{V/div}.
\begin{enumerate}
\item Déterminer la période $T$ puis la fréquence $f$ de cette tension.
\item Calculer la tension maximale $U_{max}$ puis la tension efficace $U_{eff}$.
\item Conclure en comparant avec les valeurs du réseau ENEO.
\end{enumerate}
\tcblower
\textbf{1. Période et fréquence.}

La période est la durée d'un motif complet :
\[ T = 4 \text{ div} \times \SI{5}{ms/div} = \SI{20}{ms} = \SI{0,020}{s}. \]
La fréquence est l'inverse de la période :
\[ f = \dfrac{1}{T} = \dfrac{1}{0{,}020} = \SI{50}{Hz}. \]

\textbf{2. Tensions maximale et efficace.}

La tension maximale correspond à la hauteur du pic :
\[ U_{max} = 3{,}1 \text{ div} \times \SI{100}{V/div} = \SI{310}{V}. \]
La tension efficace se déduit de la tension maximale :
\[ U_{eff} = \dfrac{U_{max}}{\sqrt{2}} = \dfrac{310}{1{,}414} \approx \SI{220}{V}. \]

\textbf{3. Conclusion.}

On trouve $f = \SI{50}{Hz}$ et $U_{eff} = \SI{220}{V}$ : ce sont exactement les valeurs nominales du réseau de distribution ENEO au Cameroun. La tension mesurée à la prise est donc conforme.
\end{exoresolu}

\coursec{Production de l'énergie électrique : les centrales}

\courssub{Le principe de l'alternateur}

\begin{propriete}[Alternateur]
Un \cle{alternateur} est une machine tournante qui convertit de l'\cle{énergie mécanique} (rotation) en \cle{énergie électrique} sous forme de courant alternatif. Il est constitué d'un \cle{rotor} (aimant ou électroaimant tournant) et d'un \cle{stator} (bobines fixes). La variation du flux magnétique induit une tension alternative aux bornes du stator (loi de Faraday-Lenz).
\end{propriete}

\courssub{Les principales centrales au Cameroun}

\begin{savaistu}[Le mix électrique camerounais]
Le Cameroun produit son électricité principalement par :
\begin{itemize}
\item \textbf{Centrales hydrauliques} (barrages) : \textbf{Song Loulou} (\SI{384}{MW}), \textbf{Lagdo} (\SI{72}{MW}), \textbf{Mékin} (\SI{15}{MW}), \textbf{Nachtigal} (\SI{420}{MW} en mise en service). L'eau de retenue actionne des turbines Francis ou Kaplan. Énergie potentielle $\rightarrow$ cinétique $\rightarrow$ mécanique $\rightarrow$ électrique. Faible émission CO$_2$, mais dépendance à la pluviométrie.
\item \textbf{Centrales thermiques} (combustion) : \textbf{Dibamba} (fioul lourd, \SI{88}{MW}), \textbf{Kribi} (gaz naturel, \SI{216}{MW}), \textbf{Limbe} (fioul). Un moteur Diesel ou une turbine à gaz entraîne l'alternateur. Démarrage rapide, mais coût du carburant et pollution.
\item \textbf{Énergies renouvelables émergentes :} solaire (centrale de Maroua, \SI{20}{MW}), biomasse. Complémentaires pour l'électrification rurale.
\end{itemize}
\textbf{Point commun :} toutes ces centrales possèdent des \textbf{alternateurs} qui délivrent du courant alternatif triphasé (trois tensions décalées de 120°).
\end{savaistu}

\coursec{Transport et distribution : lignes haute tension et transformateurs}

\courssub{La chaîne de transport : pourquoi la haute tension ?}

\begin{methode}[Raconter la chaîne « centrale $\rightarrow$ habitation »]
Pour rédiger un parcours complet de l'énergie électrique, suivre les étapes dans l'ordre :
\begin{enumerate}
\item \textbf{Production} : une \cle{centrale} (hydraulique, thermique, solaire...) fait tourner un \cle{alternateur} qui transforme l'énergie mécanique en énergie électrique alternative (tension moyenne, ex. \SI{15}{kV}).
\item \textbf{Élévation de tension} : un transformateur \cle{élévateur} (poste de départ) porte la tension à la \cle{haute tension} (HT : \SI{90}{kV}, \SI{225}{kV} au Cameroun).
\item \textbf{Transport} : les \cle{lignes haute tension} (pylônes, câbles nus en aluminium-acier) acheminent l'électricité sur de longues distances.
\item \textbf{Abaissement de tension} : des transformateurs \cle{abaisseurs} (postes sources, postes de quartier) ramènent la tension successivement en moyenne tension (MT \SI{15}{kV} ou \SI{30}{kV}) puis en \cle{basse tension} (BT \SI{220}{V} monophasé / \SI{380}{V} triphasé) près des habitations.
\item \textbf{Distribution} : le réseau ENEO alimente les maisons ; le \cle{compteur} mesure l'énergie consommée en kWh.
\end{enumerate}
\textbf{Justification clé à citer (exigible au BEPC) :} on transporte l'électricité sous haute tension car, à puissance transportée égale ($P = U \times I$), la tension $U$ élevée impose une intensité $I$ faible ; or les pertes par \cle{effet Joule} dans les câbles ($P_{pertes} = R \cdot I^2$) sont proportionnelles au carré de l'intensité : elles sont donc fortement réduites.
\end{methode}

\begin{popfigure}
\begin{tikzpicture}[scale=0.9, >=Stealth, node distance=1.5cm and 2.5cm]
% Styles
\tikzset{
  block/.style={draw=popdark, thick, fill=popcream, rounded corners, minimum width=2.2cm, minimum height=1.2cm, align=center, font=\small},
  arrow/.style={->, thick, popblue},
  label/.style={font=\tiny, popdark, midway, above, sloped}
}
% Nœuds
\node[block, align=center] (centrale){Centrale\\(Song Loulou)\\Alternateur\\$U \approx \SI{15}{kV}$};
\node[block, right=of centrale, align=center] (transfo_up){Transformateur\\Élévateur\\$U \rightarrow \SI{225}{kV}$};
\node[block, right=of transfo_up, align=center] (ligne){Lignes Haute Tension\\Pylônes \SI{225}{kV}\\Faibles pertes ($I$ faible)};
\node[block, right=of ligne, align=center] (transfo_down){Transformateur\\Abaisseur (Poste)\\$U \rightarrow \SI{15}{kV} / \SI{220}{V}$};
\node[block, right=of transfo_down, align=center] (compteur){Compteur ENEO\\Mesure énergie (kWh)};
\node[block, below=of compteur, align=center] (maison){Installation\\Domestique\\Prises \SI{220}{V} $\sim$};

% Flèches
\draw[arrow] (centrale) -- node[label] {Courant alternatif} (transfo_up);
\draw[arrow] (transfo_up) -- node[label] {Transport HT} (ligne);
\draw[arrow] (ligne) -- node[label] {Arrivée poste} (transfo_down);
\draw[arrow] (transfo_down) -- node[label] {Distribution BT} (compteur);
\draw[arrow] (compteur) -- (maison);

% Annotations techniques
\node[poporange, font=\tiny, align=center] at ($(centrale.east)!0.5!(transfo_up.west)$) {Élévation $U$};
\node[poporange, font=\tiny, align=center] at ($(transfo_up.east)!0.5!(ligne.west)$) {Transport $P=UI$ cte};
\node[poporange, font=\tiny, align=center] at ($(ligne.east)!0.5!(transfo_down.west)$) {Abaissement $U$};
\end{tikzpicture}
\legende{Chaîne de production, transport et distribution de l'énergie électrique (exemple Song Loulou $\rightarrow$ Yaoundé).}
\end{popfigure}

\courssub{Le transformateur : principe et relation de transformation}

\begin{definition}[Transformateur]
Un \cle{transformateur} est un composant statique (pas de pièces mobiles) qui permet de modifier la valeur de la tension alternative. Il est constitué de deux enroulements de fil de cuivre (le \cle{primaire} et le \cle{secondaire}) enroulés sur un même \cle{noyau en fer laminé} (circuit magnétique fermé).
\end{definition}

\begin{propriete}[Relation de transformation (transformateur parfait)]
Pour un transformateur idéal (sans pertes) :
\[ \dfrac{U_2}{U_1} = \dfrac{N_2}{N_1} = \dfrac{I_1}{I_2} \]
où $U_1, N_1, I_1$ concernent le primaire (entrée) et $U_2, N_2, I_2$ le secondaire (sortie).
\begin{itemize}
\item Si $N_2 > N_1$ alors $U_2 > U_1$ : \cle{transformateur élévateur}.
\item Si $N_2 < N_1$ alors $U_2 < U_1$ : \cle{transformateur abaisseur}.
\end{itemize}
\end{propriete}

\begin{attention}[Condition d'application vitale]
Le transformateur ne fonctionne \textbf{qu'en courant alternatif}.
Le principe repose sur l'induction électromagnétique : un courant \textbf{variable} dans le primaire crée un champ magnétique \textbf{variable} dans le noyau, qui induit une tension au secondaire.
Avec un courant \textbf{continu} (pile, batterie) : le champ magnétique est constant $\rightarrow$ aucune tension induite au secondaire ($U_2 = 0$) $\rightarrow$ le primaire, de très faible résistance, est traversé par un courant intense $\rightarrow$ \textbf{échauffement dangereux et destruction}.
\end{attention}

\begin{popfigure}
\begin{tikzpicture}[scale=1.1, >=Stealth]
% Noyau
\draw[popdark, very thick, fill=popmuted!30] (0,0) rectangle (1,3);
\draw[popdark, very thick, fill=popmuted!30] (3,0) rectangle (4,3);
\draw[popdark, very thick, fill=popmuted!30] (0,3) -- (1,3) -- (1,3.5) -- (3,3.5) -- (3,3) -- (4,3) -- (4,0) -- (3,0) -- (3,-0.5) -- (1,-0.5) -- (1,0) -- cycle;
% Enroulements (représentation simplifiée)
\foreach \y in {0.4, 0.9, 1.4, 1.9, 2.4, 2.9} {
  \draw[popblue, thick] (0.1,\y) ellipse (0.35 and 0.15);
  \draw[poporange, thick] (3.1,\y) ellipse (0.35 and 0.15);
}
% Légendes
\node[popblue, font=\small\bfseries] at (0.5, -0.8) {Primaire ($N_1$ spires)};
\node[poporange, font=\small\bfseries] at (3.5, -0.8) {Secondaire ($N_2$ spires)};
\node[popdark, font=\small\bfseries] at (2, 3.8) {Noyau fer laminé};
% Flux magnétique
\draw[popgreen, ->, thick, decorate, decoration={snake, amplitude=1pt, segment length=4pt}] (1, 1.5) -- (3, 1.5) node[midway, above] {Flux $\Phi$ variable};
% Tensions
\draw[->, popblue] (-1, 1.5) -- (0, 1.5) node[midway, above] {$U_1 \sim$};
\draw[->, poporange] (4, 1.5) -- (5, 1.5) node[midway, above] {$U_2 \sim$};
\end{tikzpicture}
\legende{Schéma de principe d'un transformateur monophasé : noyau ferromagnétique, enroulements primaire et secondaire, flux magnétique variable.}
\end{popfigure}

\courssub{Méthode : Exploiter la relation du transformateur}

\begin{methode}[Appliquer la relation du transformateur]
Pour un transformateur parfait :
\begin{enumerate}
\item Repérer le primaire (enroulement d'entrée, $N_1$ spires, tension $U_1$) et le secondaire (sortie, $N_2$ spires, tension $U_2$).
\item Écrire la relation : \[ \dfrac{U_2}{U_1} = \dfrac{N_2}{N_1}. \]
\item Conclure sur le rôle : si $N_2 > N_1$, alors $U_2 > U_1$ : le transformateur est \cle{élévateur} ; si $N_2 < N_1$, il est \cle{abaisseur}.
\item Isoler l'inconnue avant de calculer : $U_2 = U_1 \times \dfrac{N_2}{N_1}$.
\end{enumerate}
\textbf{Condition d'application :} le transformateur ne fonctionne \textbf{qu'en courant alternatif}. Branché sur du continu, il ne transforme rien et peut chauffer dangereusement.
\end{methode}

\begin{exoresolu}[Le transformateur du quartier]
Un transformateur ENEO de quartier possède $N_1 = \num{15000}$ spires au primaire et $N_2 = 300$ spires au secondaire. Il est alimenté au primaire par une tension $U_1 = \SI{15}{kV}$.
\begin{enumerate}
\item Calculer la tension $U_2$ disponible au secondaire.
\item Ce transformateur est-il élévateur ou abaisseur ? Justifier.
\item Pourquoi ne peut-on pas alimenter ce transformateur avec une batterie ?
\end{enumerate}
\tcblower
\textbf{1. Calcul de $U_2$.}

D'après la relation du transformateur :
\[ U_2 = U_1 \times \dfrac{N_2}{N_1} = \SI{15000}{V} \times \dfrac{300}{\num{15000}} = \SI{300}{V}. \]
(On rappelle que $\SI{15}{kV} = \SI{15000}{V}$.) En tenant compte des pertes réelles et du réglage des prises du transformateur, la tension distribuée aux habitations est ramenée à la valeur normalisée de \SI{220}{V}.

\textbf{2. Rôle du transformateur.}

On a $N_2 = 300 < N_1 = \num{15000}$, donc $U_2 < U_1$ : la tension diminue entre le primaire et le secondaire. Ce transformateur est donc un \cle{abaisseur} de tension, ce qui correspond bien à son rôle : ramener la moyenne tension du réseau à la tension domestique.

\textbf{3. Impossibilité d'alimenter avec une batterie.}

Une batterie fournit un courant \cle{continu}. Or le fonctionnement du transformateur repose sur la variation du courant dans le primaire, qui crée un champ magnétique variable induisant une tension au secondaire. Avec un courant continu, le champ magnétique est constant : aucune tension n'apparaît au secondaire, et le primaire, parcouru par un courant continu intense (faible résistance ohmmique du cuivre), chauffe et peut être détruit.

\textbf{Conclusion :} $U_2 = \SI{300}{V}$ (ramenée à \SI{220}{V} pour les abonnés) ; c'est un abaisseur ; il ne fonctionne qu'en courant alternatif.
\end{exoresolu}

\coursec{De la centrale à l'habitation : la chaîne ENEO et le compteur}

\courssub{Étude de cas complète : Du barrage à l'ampoule}

\begin{exoresolu}[Du barrage de Song Loulou à la maison]
Décrire, en quatre étapes numérotées, le trajet de l'énergie électrique depuis la centrale hydroélectrique de Song Loulou (sur la Sanaga) jusqu'à une ampoule d'une maison de Yaoundé. On précisera le rôle du transformateur et on justifiera l'usage de la haute tension.
\tcblower
\textbf{Étape 1 — Production à la centrale.} L'eau du barrage de Song Loulou, retenue en hauteur, tombe sur des turbines. Celles-ci entraînent un \cle{alternateur} qui convertit l'énergie mécanique de rotation en énergie électrique sous forme de courant alternatif triphasé (tension $\approx \SI{15}{kV}$).

\textbf{Étape 2 — Élévation de la tension.} À la sortie de la centrale, un transformateur \cle{élévateur} augmente la tension jusqu'à la haute tension (\SI{225}{kV} sur l'axe Edéa-Yaoundé).

\textbf{Étape 3 — Transport par les lignes haute tension.} L'électricité parcourt les \SI{\approx 200}{km} entre Edéa et Yaoundé par les câbles des pylônes. On utilise la haute tension car, à puissance transportée égale ($P = U \times I$), une tension $U$ élevée impose une intensité $I$ faible ; or les pertes par effet Joule dans les câbles ($P_{pertes} = R I^2$) dépendent du carré de l'intensité : elles sont donc fortement réduites.

\textbf{Étape 4 — Abaissement et distribution.} En arrivant en ville, des transformateurs \cle{abaisseurs} (postes sources ENEO puis transformateurs de quartier sur poteaux ou en cabines) ramènent la tension successivement en moyenne tension puis en basse tension \SI{220}{V}. Le courant arrive alors au \cle{compteur} de la maison (mesure en kWh), puis au tableau de répartition (disjoncteur général, disjoncteurs divisionnaires) et enfin à l'ampoule.

\textbf{Conclusion :} l'énergie suit la chaîne : centrale $\rightarrow$ transformateur élévateur $\rightarrow$ ligne haute tension $\rightarrow$ transformateur abaisseur $\rightarrow$ compteur $\rightarrow$ installation domestique.
\end{exoresolu}

\courssub{Puissance, énergie et facture : le compteur ENEO}

\begin{methode}[Calculer une puissance et une énergie électrique]
\begin{enumerate}
\item Identifier les données : tension $U$ (en V), intensité $I$ (en A), durée $t$.
\item Calculer la \cle{puissance} absorbée : \[ P = U \times I \qquad (U \text{ en V},\ I \text{ en A},\ P \text{ en W}). \]
\item Calculer l'\cle{énergie} consommée : \[ E = P \times t. \]
\begin{itemize}
\item Avec $P$ en W et $t$ en s, $E$ est en \cle{joules} (J).
\item Avec $P$ en kW et $t$ en h, $E$ est en \cle{kilowattheures} (kWh), l'unité du compteur ENEO et de la facture.
\end{itemize}
\item \textbf{Conversions utiles :} $\SI{1}{kW} = \SI{1000}{W}$ ; $\SI{1}{h} = \SI{3600}{s}$ ; $\SI{1}{kWh} = \SI{3 600 000}{J}$.
\item Pour la facture : $\text{Coût (FCFA)} = E \text{ (kWh)} \times \text{prix du kWh (FCFA/kWh)}$.
\end{enumerate}
\textbf{Piège :} ne jamais mélanger les unités : convertir les minutes en heures ou en secondes, les watts en kilowatts, \textbf{avant} tout calcul.
\end{methode}

\begin{exoresolu}[Facture d'un fer à repasser]
Un fer à repasser porte les indications « 220 V ; 5 A ». Une ménagère de Douala l'utilise 2 heures par jour pendant 30 jours. Le kilowattheure est facturé \SI{99}{FCFA} par ENEO (tarif résidentiel tranche basse).
\begin{enumerate}
\item Calculer la puissance du fer.
\item Calculer, en kWh, l'énergie consommée en un mois.
\item En déduire le montant de la facture mensuelle due à ce fer.
\end{enumerate}
\tcblower
\textbf{1. Puissance du fer.}

\[ P = U \times I = 220 \times 5 = \SI{1100}{W} = \SI{1,1}{kW}. \]

\textbf{2. Énergie consommée en un mois.}

Durée totale d'utilisation :
\[ t = 2 \text{ h/jour} \times 30 \text{ jours} = \SI{60}{h}. \]
Énergie consommée (avec $P$ en kW et $t$ en h) :
\[ E = P \times t = 1{,}1 \times 60 = \SI{66}{kWh}. \]

\textbf{3. Montant de la facture.}

\[ \text{Coût} = E \times \text{prix du kWh} = 66 \times 99 = \SI{6534}{FCFA}. \]

\textbf{Conclusion :} le fer à repasser a une puissance de \SI{1100}{W} ; il consomme \SI{66}{kWh} par mois, ce qui coûte \SI{6534}{FCFA} à la ménagère (hors abonnement et taxes).
\end{exoresolu}

\coursec{Sécurité avec le courant électrique domestique}

\courssub{Dangers, risques et dispositifs de protection}

\begin{definition}[Électrisation et Électrocution]
L'\cle{électrisation} est le passage d'un courant électrique dans le corps humain. Elle peut provoquer des brûlures, des troubles du rythme cardiaque (fibrillation ventriculaire), des contractions musculaires tétanisantes. L'\cle{électrocution} est une électrisation mortelle. Le danger dépend de l'intensité, de la durée, du parcours (main-pied traverse le cœur) et de la fréquence (50 Hz très dangereux).
\end{definition}

\begin{definition}[Court-circuit]
Mise en contact direct des deux bornes d'un générateur (ou d'un circuit) par un conducteur de résistance quasi nulle. L'intensité devient énorme (limité seulement par la résistance interne) $\rightarrow$ échauffement violent, fusion, projection de métal fondu, incendie.
\end{definition}

\begin{propriete}[Dispositifs de protection obligatoires dans l'installation domestique (norme NF C 15-100 / norme camerounaise)]
\begin{itemize}
\item \textbf{Disjoncteur général (ou disjoncteur d'abonné) :} coupe l'alimentation de toute l'installation en cas de surcharge ou court-circuit sur le réseau amont. Calibre selon puissance souscrite (ex. 30 A, 45 A).
\item \textbf{Disjoncteurs divisionnaires (ou fusibles) :} un par circuit (éclairage, prises, four, machine à laver). Protègent les fils du circuit contre la surcharge et le court-circuit. Calibre adapté à la section du fil (ex. 10 A pour éclairage 1,5 mm², 16 A pour prises 2,5 mm², 20-32 A pour circuits spécialisés).
\item \textbf{Disjoncteur différentiel (interrupteur différentiel) :} \textbf{protection des personnes}. Détecte une différence de courant entre phase et neutre (courant de fuite vers la terre). Coupe instantanément (30 mA, 30 ms) en cas de défaut d'isolement (ex. machine à laver mouillée, fil dénudé touché). \textbf{Obligatoire} sur tous les circuits (sauf cas rares).
\item \textbf{Prise de terre (borne verte/jaune) :} relie les parties métalliques des appareils (classe I) à la terre. En cas de fuite, le courant part vers la terre $\rightarrow$ déclenche le différentiel. \textbf{Indispensable} pour la sécurité.
\item \textbf{Liaison équipotentielle :} relie entre elles toutes les arrivées métalliques (eau, gaz, chauffage) et la prise de terre pour éviter des différences de potentiel dangereuses.
\end{itemize}
\end{propriete}

\courssub{Méthode : Analyser une situation de danger domestique}

\begin{methode}[Analyser une situation de danger domestique]
Face à une situation concrète (exercice ou vie courante) :
\begin{enumerate}
\item \textbf{Repérer le danger} : fil dénudé, prise surchargée (rallonge multipliant les appareils), appareil électrique près de l'eau (salle de bain, évier), intervention sur une installation sous tension, absence de prise de terre sur appareil métallique.
\item \textbf{Nommer le risque} : \cle{électrisation} (passage du courant dans le corps), \cle{incendie} (échauffement par surintensité / effet Joule), \cle{court-circuit}.
\item \textbf{Citer la protection adaptée} : \cle{fusible} ou \cle{disjoncteur} (coupe le circuit en cas de surintensité), \cle{disjoncteur différentiel} (coupe en cas de fuite de courant vers la terre), \cle{prise de terre} (évacue les courants de fuite), gaines isolantes, obturateurs de prise.
\item \textbf{Citer le geste de sécurité} : couper le disjoncteur général avant toute intervention ; ne jamais toucher un appareil électrique avec les mains mouillées ; ne pas surcharger une multiprise ; en cas d'électrisation d'une personne, \textbf{couper d'abord le courant} (disjoncteur ou débrancher) avant de la secourir (ne pas toucher la victime tant qu'elle est sous tension).
\end{enumerate}
\textbf{Idée clé :} le fusible / disjoncteur se calibre à une intensité légèrement supérieure à l'intensité normale de fonctionnement du circuit, mais inférieure à l'intensité maximale admissible par les fils.
\end{methode}

\begin{exoresolu}[La rallonge surchargée]
Dans un salon à Bafoussam, une famille branche sur une même multiprise : un téléviseur ($I_1 = \SI{0,8}{A}$), un réfrigérateur ($I_2 = \SI{1,5}{A}$), un fer à repasser ($I_3 = \SI{5}{A}$) et une bouilloire ($I_4 = \SI{6}{A}$). La multiprise supporte au maximum \SI{10}{A}.
\begin{enumerate}
\item Calculer l'intensité totale traversant la multiprise.
\item Cette situation est-elle dangereuse ? Justifier et nommer le risque.
\item Citer deux mesures de sécurité à prendre.
\end{enumerate}
\tcblower
\textbf{1. Intensité totale.}

Les appareils branchés sur la même multiprise sont en dérivation : les intensités s'additionnent.
\[ I = I_1 + I_2 + I_3 + I_4 = 0{,}8 + 1{,}5 + 5 + 6 = \SI{13,3}{A}. \]

\textbf{2. Danger.}

L'intensité totale $I = \SI{13,3}{A}$ dépasse l'intensité maximale supportée par la multiprise (\SI{10}{A}). La situation est donc \textbf{dangereuse} : les fils de la rallonge, trop fins pour une telle intensité, s'échauffent par \cle{effet Joule} ($P = R I^2$). Cet échauffement peut fondre l'isolant, provoquer un court-circuit puis un \textbf{incendie}.

\textbf{3. Mesures de sécurité.}
\begin{itemize}
\item Ne pas brancher simultanément les appareils de forte puissance (fer, bouilloire) sur la même multiprise : utiliser des prises murales distinctes (circuits séparés protégés par disjoncteurs 16 A ou 20 A).
\item Installer (ou vérifier) un \cle{disjoncteur} ou un \cle{fusible} calibré à \SI{10}{A} (ou 16 A selon le câble) sur ce circuit : il coupera automatiquement le courant dès que l'intensité dépasse la valeur admissible.
\end{itemize}

\textbf{Conclusion :} avec $I = \SI{13,3}{A} > \SI{10}{A}$, la multiprise est surchargée : risque d'échauffement et d'incendie. Il faut répartir les appareils sur plusieurs prises et protéger le circuit par un disjoncteur adapté.
\end{exoresolu}

\courssub{Erreurs à éviter au BEPC}

\begin{attention}[Les 5 erreurs qui coûtent des points au BEPC]
\begin{enumerate}
\item \textbf{Confondre tension maximale et tension efficace.} La valeur \SI{220}{V} du réseau ENEO est une tension \emph{efficace}. Sur un oscillogramme, on mesure d'abord $U_{max}$, puis on calcule $U_{eff} = \dfrac{U_{max}}{\sqrt{2}}$. Oublier le $\sqrt{2}$ est l'erreur la plus fréquente.
\item \textbf{Oublier de convertir les unités de temps.} Avant de calculer $f = \dfrac{1}{T}$, convertir les millisecondes en secondes ($\SI{20}{ms} = \SI{0,020}{s}$). Pour $E = P \times t$ en kWh, convertir les minutes en heures et les watts en kilowatts.
\item \textbf{Additionner des puissances avec des intensités mal converties.} Dans $P = U \times I$, $U$ doit être en volts et $I$ en ampères : $\SI{15}{kV} = \SI{15000}{V}$. Une intensité en mA doit être convertie en A ($1 \text{ mA} = 10^{-3} \text{ A}$).
\item \textbf{Croire que le transformateur fonctionne en continu.} Le transformateur ne fonctionne \emph{qu'en courant alternatif}. Écrire qu'on peut le brancher sur une pile ou une batterie est une erreur grave de cours.
\item \textbf{Rédiger une réponse de sécurité sans nommer le dispositif ni le risque.} « C'est dangereux » ne suffit pas : il faut nommer le risque (électrisation, incendie par effet Joule, court-circuit) ET la protection (fusible, disjoncteur, disjoncteur différentiel, prise de terre) ET, si demandé, le geste (couper le courant avant d'intervenir ou de secourir une victime).
\end{enumerate}
\end{attention}

\newpage

\coursec{Exercices}

\courssub{Je vérifie mes connaissances}

\begin{exercice}[QCM : continu ou alternatif ?]
Pour chaque question, entoure la (ou les) bonne(s) réponse(s).
\begin{enumerate}
\item Un courant continu :
\begin{enumerate}
\item change de sens périodiquement ;
\item circule toujours dans le même sens ;
\item est fourni par une pile.
\end{enumerate}
\item Le courant distribué par ENEO dans les habitations est :
\begin{enumerate}
\item continu ;
\item alternatif sinusoïdal ;
\item sans danger.
\end{enumerate}
\item Quelle(s) source(s) fournit (fournissent) une tension alternative ?
\begin{enumerate}
\item la pile plate ;
\item l'alternateur d'une centrale ;
\item la dynamo d'une moto-taxi ;
\item le secteur ENEO.
\end{enumerate}
\item La fréquence du courant ENEO au Cameroun est :
\begin{enumerate}
\item \SI{50}{Hz} ;
\item \SI{220}{Hz} ;
\item \SI{225}{kHz}.
\end{enumerate}
\end{enumerate}
\end{exercice}

\begin{exercice}[Vrai ou faux ? Justifie]
Réponds par VRAI ou FAUX, puis justifie ta réponse en une ou deux phrases.
\begin{enumerate}
\item La tension efficace du secteur ENEO est égale à sa tension maximale.
\item On transporte l'énergie électrique sous haute tension pour réduire les pertes par effet Joule dans les lignes.
\item Un transformateur élévateur est placé à l'entrée de l'habitation, juste avant le compteur.
\item Le fusible protège les installations en coupant le circuit lorsque l'intensité devient trop grande.
\end{enumerate}
\end{exercice}

\begin{exercice}[Définitions]
Définis en une phrase chacune des notions suivantes, en utilisant le vocabulaire précis de la leçon :
\begin{enumerate}
\item courant alternatif sinusoïdal ;
\item période et fréquence d'une tension alternative ;
\item tension maximale et tension efficace ;
\item transformateur.
\end{enumerate}
\end{exercice}

\begin{exercice}[Calculs directs]
\begin{enumerate}
\item La tension du secteur ENEO a pour valeur efficace $U = \SI{220}{V}$. Calcule sa tension maximale $U_{\max}$ (on rappelle que $U_{\max} = U\sqrt{2}$ ; on prendra $\sqrt{2} \approx 1{,}41$).
\item La période du courant ENEO est $T = \SI{0,02}{s}$. Calcule sa fréquence $f$.
\item Une lampe porte les indications « \SI{220}{V} ; \SI{60}{W} ». Calcule l'intensité efficace du courant qui la traverse lorsqu'elle fonctionne normalement.
\end{enumerate}
\end{exercice}

\courssub{Je m'entraîne}

\begin{exercice}[Lecture d'un oscillogramme]
On visualise la tension du secteur sur l'écran d'un oscilloscope. Les réglages sont : balayage $b = \SI{5}{ms/div}$ ; sensibilité verticale $s = \SI{100}{V/div}$. On observe une sinusoïde dont une période complète occupe $4$ divisions horizontales et dont le sommet se situe à $3{,}1$ divisions au-dessus de l'axe.
\begin{enumerate}
\item Détermine la période $T$ de cette tension, puis sa fréquence $f$.
\item Détermine la tension maximale $U_{\max}$.
\item Calcule la tension efficace $U$. Le résultat est-il cohérent avec la tension du réseau ENEO ?
\end{enumerate}
\end{exercice}

\begin{exercice}[Pourquoi la haute tension ?]
Une centrale hydroélectrique fournit une puissance $P = \SI{110}{MW}$ à une ligne de transport.
\begin{enumerate}
\item Calcule l'intensité $I_1$ du courant dans la ligne si le transport s'effectuait sous $U_1 = \SI{220}{V}$.
\item Calcule l'intensité $I_2$ si le transport s'effectue sous haute tension $U_2 = \SI{225}{kV}$.
\item Les pertes par effet Joule dans la ligne sont proportionnelles à $I^2$. Calcule le rapport $\dfrac{I_1^2}{I_2^2}$ et conclus sur l'intérêt du transport en haute tension.
\end{enumerate}
\end{exercice}

\begin{exercice}[La chaîne de distribution à compléter]
Recopie et complète la chaîne suivante, qui décrit le parcours de l'énergie électrique, en indiquant le nom de chaque élément et la valeur de la tension à chaque étape :
\begin{center}
Centrale ($\ldots$ kV) $\longrightarrow$ $\ldots$ $\longrightarrow$ ligne haute tension ($\ldots$ kV) $\longrightarrow$ $\ldots$ $\longrightarrow$ quartier ($\ldots$ V) $\longrightarrow$ $\ldots$ $\longrightarrow$ prises de l'habitation ($\ldots$ V)
\end{center}
Mots et valeurs à placer : transformateur élévateur ; transformateur abaisseur ; compteur ; \SI{20}{kV} ; \SI{225}{kV} ; \SI{220}{V} (deux fois).
\end{exercice}

\begin{exercice}[Facture d'électricité]
Pendant le mois de juin (30 jours), une famille de Yaoundé utilise chaque jour :
\begin{itemize}
\item un téléviseur de puissance $P_1 = \SI{150}{W}$ pendant $5$ heures ;
\item cinq ampoules de puissance $P_2 = \SI{20}{W}$ chacune pendant $6$ heures.
\end{itemize}
\begin{enumerate}
\item Calcule, en kilowattheures, l'énergie électrique consommée par le téléviseur pendant le mois.
\item Calcule l'énergie consommée par l'ensemble des ampoules pendant le mois.
\item Calcule l'énergie totale consommée, puis le montant de la facture, le kilowattheure étant facturé \SI{50}{FCFA} par ENEO.
\end{enumerate}
\end{exercice}

\courssub{Je me prépare au Bac}

\begin{exercice}[Type BEPC : étude d'une tension alternative]
Lors d'une séance de travaux pratiques au collège, un groupe d'élèves de 3\up{e} branche un oscilloscope aux bornes d'un transformateur de laboratoire. Réglages de l'oscilloscope : balayage $b = \SI{2}{ms/div}$ ; sensibilité verticale $s = \SI{5}{V/div}$. Sur l'écran, la courbe observée est une sinusoïde : une période complète s'étend sur $10$ divisions horizontales, et la crête atteint $2{,}4$ divisions verticales.
\begin{enumerate}
\item Nomme l'appareil utilisé pour visualiser la tension et précise ce que représente chaque axe de l'écran.
\item Détermine la période $T$ de la tension observée, puis calcule sa fréquence $f$.
\item Détermine la tension maximale $U_{\max}$ du transformateur.
\item Calcule la tension efficace $U$ (on prendra $\sqrt{2} \approx 1{,}41$).
\item Cette tension est-elle continue ou alternative ? Justifie ta réponse à partir de l'allure de la courbe.
\item Le groupe veut alimenter une lampe « \SI{6}{V} ; \SI{3}{W} » avec ce transformateur. Est-ce possible sans danger pour la lampe ? Justifie par un calcul.
\end{enumerate}
\end{exercice}

\begin{exercice}[Type BEPC : de la centrale de Song-Loulou à Douala]
La centrale hydroélectrique de Song-Loulou produit de l'énergie électrique grâce à des alternateurs entraînés par des turbines. La puissance fournie à une ligne de transport est $P = \SI{90}{MW}$. L'énergie est transportée jusqu'à Douala sous une tension $U_2 = \SI{225}{kV}$, après élévation de la tension par un transformateur. À l'arrivée, un autre transformateur abaisse la tension à \SI{220}{V} pour les habitations.
\begin{enumerate}
\item Explique, à partir de l'expérience de l'aimant et de la bobine, pourquoi l'alternateur d'une centrale produit une tension alternative et non continue.
\item Quel est le rôle du transformateur placé à la sortie de la centrale ? Comment l'appelle-t-on ?
\item Calcule l'intensité $I$ du courant dans la ligne haute tension.
\item La résistance de la ligne est $R = \SI{10}{\Omega}$. Calcule la puissance $P_J$ perdue par effet Joule dans la ligne (on rappelle $P_J = R\,I^2$).
\item Calcule l'intensité $I'$ qu'il faudrait si l'on transportait la même puissance sous \SI{220}{V}, puis la puissance perdue $P'_J$ dans la même ligne. Conclus.
\item Cite deux autres types de centrales utilisées ou utilisables au Cameroun, en précisant la source d'énergie de chacune.
\end{enumerate}
\end{exercice}

\begin{exercice}[Type BEPC : installation domestique et sécurité]
Monsieur Mbarga fait installer l'électricité dans sa maison à Bafoussam. Le compteur ENEO est réglé pour une intensité maximale de \SI{15}{A} sous \SI{220}{V}. L'installation comprend : un réfrigérateur de \SI{200}{W}, un fer à repasser de \SI{1000}{W}, une bouilloire de \SI{1500}{W} et dix lampes de \SI{15}{W} chacune.
\begin{enumerate}
\item Calcule la puissance maximale $P_{\max}$ que le compteur peut débiter.
\item Tous les appareils cités fonctionnent en même temps. Calcule la puissance totale utilisée. Le compteur risque-t-il de disjoncter ? Justifie.
\item Calcule l'intensité totale du courant dans l'installation lorsque tous les appareils fonctionnent.
\item Monsieur Mbarga branche en plus un chauffe-eau de \SI{1200}{W}. Que se passe-t-il ? Explique le rôle protecteur du fusible (ou du disjoncteur) dans cette situation.
\item Cite trois consignes de sécurité à respecter avec le courant domestique, et explique le rôle de la prise de terre.
\item Le mois suivant, la famille consomme \SI{180}{kWh}. Calcule le montant de la facture à \SI{50}{FCFA} le kilowattheure.
\end{enumerate}
\end{exercice}

\newpage

\coursec{Corrigés des exercices}

\courssub{Corrigés des exercices d'application}

\begin{corrige}[Exercice 1 — QCM : continu ou alternatif ?]
\begin{enumerate}
\item Réponses \textbf{b} et \textbf{c}. Un courant continu circule toujours dans le même sens, du pôle $+$ vers le pôle $-$ à l'extérieur du générateur ; c'est le cas du courant fourni par une pile. La réponse a décrit un courant alternatif.
\item Réponse \textbf{b}. ENEO distribue un courant alternatif sinusoïdal de tension efficace \SI{220}{V} et de fréquence \SI{50}{Hz}. Il est dangereux : la réponse c est fausse.
\item Réponses \textbf{b}, \textbf{c} et \textbf{d}. L'alternateur d'une centrale, la dynamo d'une moto-taxi et le secteur ENEO fournissent des tensions alternatives (la dynamo d'une moto est en réalité un petit alternateur). La pile plate fournit une tension continue.
\item Réponse \textbf{a}. La fréquence du réseau ENEO est $f = \SI{50}{Hz}$, ce qui correspond à une période $T = \SI{0,02}{s}$.
\end{enumerate}
\end{corrige}

\begin{corrige}[Exercice 2 — Vrai ou faux ? Justifie]
\begin{enumerate}
\item \textbf{FAUX.} La tension efficace $U$ est inférieure à la tension maximale : $U = \dfrac{U_{\max}}{\sqrt{2}}$. Pour le secteur, $U = \SI{220}{V}$ alors que $U_{\max} \approx \SI{311}{V}$.
\item \textbf{VRAI.} Les pertes par effet Joule sont proportionnelles à $I^2$ ($P_J = R\,I^2$). À puissance transportée constante ($P = U\,I$), augmenter $U$ diminue $I$, donc diminue fortement les pertes.
\item \textbf{FAUX.} À l'entrée de l'habitation, c'est un transformateur \emph{abaisseur} (souvent sur un poteau du quartier) qui ramène la tension de \SI{20}{kV} à \SI{220}{V}. Le compteur ne fait que mesurer l'énergie consommée.
\item \textbf{VRAI.} Le fusible contient un fil qui fond lorsque l'intensité dépasse la valeur prévue : il coupe alors le circuit et protège l'installation contre les surintensités et les courts-circuits.
\end{enumerate}
\end{corrige}

\begin{corrige}[Exercice 3 — Définitions]
\begin{enumerate}
\item Un \cle{courant alternatif sinusoïdal} est un courant qui change de sens périodiquement, dont l'intensité varie au cours du temps selon une courbe en forme de sinusoïde.
\item La \cle{période} $T$ est la durée d'un motif élémentaire qui se répète identique à lui-même ; la \cle{fréquence} $f$ est le nombre de périodes par seconde : $f = \dfrac{1}{T}$, avec $T$ en secondes et $f$ en hertz (Hz).
\item La \cle{tension maximale} $U_{\max}$ est la plus grande valeur atteinte par la tension au cours d'une période ; la \cle{tension efficace} $U$ est la valeur qu'indique le voltmètre, liée à $U_{\max}$ par $U = \dfrac{U_{\max}}{\sqrt{2}}$.
\item Un \cle{transformateur} est un appareil qui modifie la valeur d'une tension alternative : il l'élève (transformateur élévateur) ou l'abaisse (transformateur abaisseur), sans en changer la fréquence.
\end{enumerate}
\end{corrige}

\begin{corrige}[Exercice 4 — Calculs directs]
\begin{enumerate}
\item $U_{\max} = U\sqrt{2} = 220 \times 1{,}41 = \SI{310,2}{V} \approx \SI{310}{V}$.
\item $f = \dfrac{1}{T} = \dfrac{1}{0{,}02} = \SI{50}{Hz}$. C'est bien la fréquence du réseau ENEO.
\item La lampe fonctionne normalement sous $U = \SI{220}{V}$ avec $P = \SI{60}{W}$. Or $P = U\,I$, donc :
\[ I = \frac{P}{U} = \frac{60}{220} \approx \SI{0,27}{A}. \]
\end{enumerate}
\end{corrige}

\begin{corrige}[Exercice 5 — Lecture d'un oscillogramme]
\begin{enumerate}
\item La période occupe 4 divisions horizontales avec un balayage de \SI{5}{ms/div} :
\[ T = 4 \times 5 = \SI{20}{ms} = \SI{0,02}{s}, \qquad f = \frac{1}{T} = \frac{1}{0{,}02} = \SI{50}{Hz}. \]
\item Le sommet est à $3{,}1$ divisions avec une sensibilité de \SI{100}{V/div} :
\[ U_{\max} = 3{,}1 \times 100 = \SI{310}{V}. \]
\item $U = \dfrac{U_{\max}}{\sqrt{2}} = \dfrac{310}{1{,}41} \approx \SI{220}{V}$. Le résultat est cohérent : c'est bien la tension efficace du réseau ENEO.
\end{enumerate}
\textbf{Point de vigilance :} ne pas oublier de convertir les millisecondes en secondes avant de calculer la fréquence.
\end{corrige}

\begin{corrige}[Exercice 6 — Pourquoi la haute tension ?]
\begin{enumerate}
\item $P = U_1\,I_1$ donne $I_1 = \dfrac{P}{U_1} = \dfrac{110\,000\,000}{220} = \SI{500000}{A}$. Une intensité énorme, irréaliste !
\item $I_2 = \dfrac{P}{U_2} = \dfrac{110\,000\,000}{225\,000} \approx \SI{489}{A}$.
\item Le rapport des pertes vaut :
\[ \frac{I_1^2}{I_2^2} = \left(\frac{I_1}{I_2}\right)^2 = \left(\frac{U_2}{U_1}\right)^2 = \left(\frac{225\,000}{220}\right)^2 \approx (1023)^2 \approx 1{,}05 \times 10^{6}. \]
Les pertes par effet Joule seraient environ \textbf{un million de fois plus grandes} sous \SI{220}{V} que sous \SI{225}{kV}. Conclusion : transporter l'énergie sous haute tension divise énormément l'intensité, donc les pertes en chaleur dans les lignes. C'est pourquoi on utilise des transformateurs élévateurs à la sortie des centrales.
\end{enumerate}
\end{corrige}

\begin{corrige}[Exercice 7 — La chaîne de distribution à compléter]
\begin{center}
Centrale (\SI{20}{kV}) $\longrightarrow$ \textbf{transformateur élévateur} $\longrightarrow$ ligne haute tension (\SI{225}{kV}) $\longrightarrow$ \textbf{transformateur abaisseur} $\longrightarrow$ quartier (\SI{220}{V}) $\longrightarrow$ \textbf{compteur} $\longrightarrow$ prises de l'habitation (\SI{220}{V})
\end{center}
\textbf{Justification :} à la sortie de la centrale, le transformateur élévateur porte la tension de \SI{20}{kV} à \SI{225}{kV} pour limiter les pertes pendant le transport. Près des habitations, le transformateur abaisseur ramène la tension à \SI{220}{V}, valeur utilisable et distribuée aux quartiers. Le compteur, placé à l'entrée de l'habitation, mesure l'énergie consommée ; la tension aux prises reste \SI{220}{V}.
\end{corrige}

\begin{corrige}[Exercice 8 — Facture d'électricité]
On utilise $E = P \times t$, avec $P$ en kilowatts et $t$ en heures pour obtenir $E$ en kilowattheures.
\begin{enumerate}
\item Téléviseur : $P_1 = \SI{150}{W} = \SI{0,15}{kW}$ ; durée mensuelle : $5 \times 30 = \SI{150}{h}$.
\[ E_1 = 0{,}15 \times 150 = \SI{22,5}{kWh}. \]
\item Ampoules : puissance totale $5 \times 20 = \SI{100}{W} = \SI{0,1}{kW}$ ; durée mensuelle : $6 \times 30 = \SI{180}{h}$.
\[ E_2 = 0{,}1 \times 180 = \SI{18}{kWh}. \]
\item Énergie totale : $E = E_1 + E_2 = 22{,}5 + 18 = \SI{40,5}{kWh}$.
Montant de la facture : $40{,}5 \times 50 = \SI{2025}{FCFA}$.
\end{enumerate}
\textbf{Point de vigilance :} convertir les watts en kilowatts avant d'appliquer $E = P\,t$, sinon le résultat est en wattheures et non en kilowattheures.
\end{corrige}

\begin{corrige}[Exercice 9 — Type BEPC : étude d'une tension alternative]
\begin{enumerate}
\item L'appareil est un \textbf{oscilloscope}. L'axe horizontal représente le \emph{temps} (gradué grâce au balayage) et l'axe vertical représente la \emph{tension} (graduée grâce à la sensibilité verticale).
\item Période : $T = 10 \times 2 = \SI{20}{ms} = \SI{0,02}{s}$. Fréquence :
\[ f = \frac{1}{T} = \frac{1}{0{,}02} = \SI{50}{Hz}. \]
\item Tension maximale : $U_{\max} = 2{,}4 \times 5 = \SI{12}{V}$.
\item Tension efficace : $U = \dfrac{U_{\max}}{\sqrt{2}} = \dfrac{12}{1{,}41} \approx \SI{8,5}{V}$.
\item Cette tension est \textbf{alternative} : la courbe est une sinusoïde qui passe alternativement au-dessus et en dessous de l'axe horizontal ; la tension change donc de signe (le courant change de sens) périodiquement. Une tension continue donnerait une droite horizontale.
\item La lampe est prévue pour \SI{6}{V}. Or la tension efficace du transformateur est $U \approx \SI{8,5}{V} > \SI{6}{V}$ : la lampe serait \textbf{survolée}, elle grillerait. Ce n'est donc pas possible sans danger. (La tension efficace d'alimentation doit être égale à la tension nominale de la lampe.)
\end{enumerate}
\textbf{Barème indicatif (sur 10) :} 1 pt pour la question 1 ; 2 pts pour la question 2 (1 pt pour $T$, 1 pt pour $f$ avec conversion) ; 1 pt pour la question 3 ; 1,5 pt pour la question 4 ; 1,5 pt pour la question 5 (allure de la courbe exigée) ; 2 pts pour la question 6 (comparaison de tensions + conclusion) ; 1 pt pour la présentation et les unités.

\textbf{Points de vigilance :} citer la conversion ms $\to$ s ; justifier « alternatif » par le \emph{changement de signe} de la tension, pas seulement par le mot « sinusoïde » ; comparer la tension \emph{efficace} à la tension nominale de la lampe.
\end{corrige}

\begin{corrige}[Exercice 10 — Type BEPC : de la centrale de Song-Loulou à Douala]
\begin{enumerate}
\item Dans l'expérience de l'aimant et de la bobine, lorsqu'on fait tourner l'aimant devant la bobine (ou la bobine devant l'aimant), le circuit est le siège d'une tension induite qui change de sens à chaque demi-tour : le courant oscille alternativement dans un sens puis dans l'autre. Dans l'alternateur de la centrale, la turbine entraîne un aimant (ou un électroaimant) en rotation devant des bobines fixes : la tension produite est donc \textbf{alternative}, de fréquence liée à la vitesse de rotation.
\item Le transformateur placé à la sortie de la centrale élève la tension (de \SI{20}{kV} à \SI{225}{kV}) afin de réduire l'intensité et donc les pertes par effet Joule pendant le transport. C'est un \textbf{transformateur élévateur}.
\item $I = \dfrac{P}{U_2} = \dfrac{90\,000\,000}{225\,000} = \SI{400}{A}$.
\item $P_J = R\,I^2 = 10 \times 400^2 = 10 \times 160\,000 = \SI{1600000}{W} = \SI{1,6}{MW}$, soit environ $1{,}8\,\%$ de la puissance transportée.
\item Sous \SI{220}{V} : $I' = \dfrac{90\,000\,000}{220} \approx \SI{409091}{A}$. Alors :
\[ P'_J = R\,I'^2 = 10 \times (409\,091)^2 \approx 1{,}67 \times 10^{12}\ \text{W}. \]
Cette puissance perdue serait des millions de fois supérieure à la puissance produite : le transport serait \textbf{impossible}. Conclusion : la haute tension est indispensable pour transporter l'énergie électrique sur de longues distances.
\item Exemples : une \textbf{centrale thermique} (à gasoil ou au gaz, comme les centrales de la SONARA ou de Kribi : énergie chimique des combustibles) ; une \textbf{centrale solaire} (énergie lumineuse du Soleil, convertie par des panneaux photovoltaïques). On peut aussi citer une autre centrale hydroélectrique (Édéa, Lagdo : énergie de l'eau en chute).
\end{enumerate}
\textbf{Barème indicatif (sur 12) :} 2 pts pour la question 1 (rotation + changement de sens exigés) ; 1,5 pt pour la question 2 ; 1,5 pt pour la question 3 ; 1,5 pt pour la question 4 ; 2,5 pts pour la question 5 (calculs + conclusion) ; 2 pts pour la question 6 ; 1 pt pour la présentation et les unités.

\textbf{Points de vigilance :} écrire les puissances en watts avant les calculs ; citer la formule $P = U\,I$ avant de l'appliquer ; la conclusion de la question 5 doit comparer $P'_J$ à la puissance produite ; pour la question 6, la source d'énergie de chaque centrale est exigée.
\end{corrige}

\begin{corrige}[Exercice 11 — Type BEPC : installation domestique et sécurité]
\begin{enumerate}
\item $P_{\max} = U \times I_{\max} = 220 \times 15 = \SI{3300}{W}$.
\item Puissance totale : $P = 200 + 1000 + 1500 + (10 \times 15) = 200 + 1000 + 1500 + 150 = \SI{2850}{W}$.
Comme $2850 < 3300$, le compteur \textbf{ne disjoncte pas}.
\item $I = \dfrac{P}{U} = \dfrac{2850}{220} \approx \SI{13}{A}$. On vérifie : $13 < 15$ A, cohérent avec la question 2.
\item Avec le chauffe-eau : $P' = 2850 + 1200 = \SI{4050}{W}$, soit $I' = \dfrac{4050}{220} \approx \SI{18,4}{A} > \SI{15}{A}$. L'intensité dépasse le calibre du compteur : le disjoncteur \textbf{coupe le circuit}. C'est son rôle protecteur : en cas de surintensité (trop d'appareils branchés ou court-circuit), le fusible fond ou le disjoncteur s'ouvre, ce qui interrompt le courant et évite l'échauffement des fils et les incendies.
\item Trois consignes de sécurité :
\begin{itemize}
\item ne jamais toucher un appareil électrique avec les mains mouillées ;
\item ne pas brancher trop d'appareils sur une même prise ou une même rallonge ;
\item ne jamais introduire d'objet métallique dans une prise, et faire réparer les fils dénudés.
\end{itemize}
La \textbf{prise de terre} relie les carcasses métalliques des appareils au sol : en cas de défaut d'isolement, le courant de fuite s'écoule vers la terre au lieu de traverser le corps de l'utilisateur, et le disjoncteur coupe le circuit. Elle protège les personnes contre l'électrisation.
\item Facture : $180 \times 50 = \SI{9000}{FCFA}$.
\end{enumerate}
\textbf{Barème indicatif (sur 12) :} 1,5 pt pour la question 1 ; 2 pts pour la question 2 (somme + comparaison) ; 1,5 pt pour la question 3 ; 2 pts pour la question 4 (calcul + rôle protecteur) ; 2 pts pour la question 5 (trois consignes + rôle de la terre) ; 1 pt pour la question 6 ; 2 pts pour la rigueur (formules citées, unités, présentation).

\textbf{Points de vigilance :} citer la formule $P = U\,I$ avant chaque application ; la conclusion de la question 2 exige la \emph{comparaison} avec \SI{3300}{W} ; pour la question 4, il faut montrer que $I' > \SI{15}{A}$ et expliquer la coupure ; le rôle de la prise de terre doit mentionner l'évacuation du courant de défaut vers le sol.
\end{corrige}

\newpage

\coursec{Fiche bilan}

\begin{popfigure}
\begin{tikzpicture}[mindmap, grow cyclic, every node/.style={concept, font=\small},
  root concept/.append style={concept color=popblue, font=\bfseries},
  level 1/.append style={level distance=3.4cm, sibling angle=60},
  level 1 concept/.append style={font=\footnotesize}]
\node[root concept]{Le courant alternatif}
  child[concept color=poporange]{ node {Deux types de courant} 
    child[concept color=poporangeL]{ node[concept, font=\scriptsize]{Continu : sens fixe (pile)} }
    child[concept color=poporangeL]{ node[concept, font=\scriptsize]{Alternatif : sens change (secteur)} }
  }
  child[concept color=popgreen]{ node {Tension sinuso\"idale}
    child[concept color=popgreenL]{ node[concept, font=\scriptsize]{$U_{\max}$, $U_{\mathrm{eff}} = U_{\max}/\sqrt{2}$} }
    child[concept color=popgreenL]{ node[concept, font=\scriptsize]{P\'eriode $T$, fr\'equence $f = 1/T$} }
  }
  child[concept color=popgold]{ node {Production : centrales}
    child[concept color=popgoldL]{ node[concept, font=\scriptsize]{Hydraulique, thermique, solaire} }
    child[concept color=popgoldL]{ node[concept, font=\scriptsize]{Alternateur = conversion d'\'energie} }
  }
  child[concept color=poppurple]{ node {Transport}
    child[concept color=poppurpleL]{ node[concept, font=\scriptsize]{Lignes haute tension : moins de pertes} }
    child[concept color=poppurpleL]{ node[concept, font=\scriptsize]{Transformateur : \'el\`eve ou abaisse $U$} }
  }
  child[concept color=poppink]{ node {Distribution ENEO}
    child[concept color=poppinkL]{ node[concept, font=\scriptsize]{\SI{220}{V}, \SI{50}{Hz} au compteur} }
  }
  child[concept color=popdark]{ node {S\'ecurit\'e}
    child[concept color=popdark!30]{ node[concept, font=\scriptsize]{Fusible, disjoncteur, terre} }
    child[concept color=popdark!30]{ node[concept, font=\scriptsize]{Danger : \'electrisation} }
  };
\end{tikzpicture}
\legende{Carte mentale de la le\c{c}on : les six id\'ees forces du courant alternatif.}
\end{popfigure}

\begin{bilan}
\textbf{D\'efinitions cl\'es}
\begin{itemize}
  \item \cle{Courant continu} : courant qui circule toujours dans le m\^eme sens (pile, batterie, accumulateur).
  \item \cle{Courant alternatif} : courant dont le sens change r\'eguli\`erement au cours du temps (secteur ENEO, groupe \'electrog\`ene).
  \item \cle{Tension alternative sinuso\"idale} : tension qui varie en suivant une sinuso\"ide ; elle est caract\'eris\'ee par sa \cle{tension maximale} $U_{\max}$, sa \cle{tension efficace} $U_{\mathrm{eff}}$, sa \cle{p\'eriode} $T$ et sa \cle{fr\'equence} $f$.
  \item \cle{Alternateur} : machine qui convertit l'\'energie m\'ecanique (rotation) en \'energie \'electrique alternative.
  \item \cle{Transformateur} : appareil qui modifie la valeur d'une tension alternative (il ne fonctionne \emph{pas} en courant continu).
\end{itemize}

\textbf{Formules et valeurs \`a conna\^itre par c\oe ur}
\begin{center}
\begin{tabularx}{\linewidth}{|l|X|l|}\hline
\rowcolor{popblueL} \textbf{Grandeur} & \textbf{Relation} & \textbf{Unit\'e} \\ \hline
Tension efficace & $U_{\mathrm{eff}} = \dfrac{U_{\max}}{\sqrt{2}}$ & volt (V) \\ \hline
Fr\'equence & $f = \dfrac{1}{T}$ & hertz (Hz) \\ \hline
P\'eriode & $T = \dfrac{1}{f}$ & seconde (s) \\ \hline
Secteur camerounais & $U_{\mathrm{eff}} = \SI{220}{V}$ ; $f = \SI{50}{Hz}$ & --- \\ \hline
\end{tabularx}
\end{center}

\textbf{La cha\^ine de l'\'energie \'electrique (m\'ethode express)}
\begin{enumerate}
  \item \textbf{Production} : centrale (hydraulique, thermique, solaire) $\rightarrow$ alternateur.
  \item \textbf{\'El\'evation} : transformateur \emph{\'el\'evateur} $\rightarrow$ lignes haute tension (pertes par effet Joule r\'eduites car $P_{\mathrm{pertes}} = R I^2$ diminue quand $I$ diminue).
  \item \textbf{Abaissement} : transformateurs \emph{abaisseurs} successifs jusqu'\`a \SI{220}{V}.
  \item \textbf{Distribution} : compteur ENEO $\rightarrow$ disjoncteur $\rightarrow$ installation de l'habitation.
\end{enumerate}

\textbf{S\'ecurit\'e : r\'eflexes indispensables}
\begin{itemize}
  \item Ne jamais toucher un appareil \'electrique avec les mains mouill\'ees ; ne jamais surcharger une rallonge.
  \item Le \cle{fusible} et le \cle{disjoncteur} coupent le circuit en cas de surintensit\'e ; la \cle{prise de terre} prot\`ege des d\'efauts d'isolement.
  \item Couper le courant au disjoncteur g\'en\'eral avant toute intervention sur l'installation.
\end{itemize}
\end{bilan}

\begin{aretenir}[Checklist avant le BEPC]
\begin{itemize}
  \item $\square$ Je sais distinguer courant continu et courant alternatif et citer un exemple de chaque.
  \item $\square$ Je connais l'allure d'une tension sinuso\"idale et je sais rep\'erer $U_{\max}$ et $T$ sur le graphique.
  \item $\square$ Je sais calculer $U_{\mathrm{eff}} = U_{\max}/\sqrt{2}$ et $f = 1/T$.
  \item $\square$ Je retiens les valeurs du secteur : \SI{220}{V} et \SI{50}{Hz}.
  \item $\square$ Je sais citer trois types de centrales et le r\^ole de l'alternateur.
  \item $\square$ Je sais expliquer pourquoi on transporte l'\'energie en haute tension (r\'eduction des pertes par effet Joule).
  \item $\square$ Je sais que le transformateur ne fonctionne qu'en courant alternatif.
  \item $\square$ Je sais ordonner la cha\^ine : centrale $\rightarrow$ haute tension $\rightarrow$ transformateurs $\rightarrow$ compteur $\rightarrow$ habitation.
  \item $\square$ Je connais le r\^ole du fusible, du disjoncteur et de la prise de terre.
  \item $\square$ Je sais citer trois r\`egles de s\'ecurit\'e avec l'\'electricit\'e domestique.
\end{itemize}
\end{aretenir}

\findecours

\end{document}
